\input{preamble}

% OK, start here.
%
\begin{document}

\title{Réduction semi-stable}


\maketitle

\phantomsection
\label{section-phantom}

\tableofcontents

\section{Introduction}
\label{section-introduction}

\noindent
Dans ce chapitre, nous démontrons le théorème de réduction semi-stable pour les courbes.
Nous utiliserons la méthode d'Artin et Winters exposée dans leur article
\cite{Artin-Winters}.

\medskip\noindent
On peut en fait démontrer le théorème de réduction semi-stable
pour les courbes sans aucun résultat de désingularisation. En effet, il
existe une méthode pour établir l'existence et la projectivité des espaces de modules
de courbes semi-stables au moyen de la théorie géométrique des invariants (GIT)
développée par Mumford, voir \cite{GIT}. Cette méthode a été
défendue par Gieseker, qui a démontré le résultat complet dans ses
notes de cours \cite{Gieseker}\footnote{Les notes de cours de Gieseker
sont rédigées sur un corps algébriquement clos, mais la même méthode
s'applique sur $\mathbf{Z}$.}. C'est là une prouesse
assez remarquable : il semble quelque peu contraire à l'intuition que l'on puisse
démontrer un tel résultat sans jamais étudier véritablement les familles de courbes sur
une base de dimension strictement positive.

\medskip\noindent
Historiquement, la première démonstration du théorème de réduction semi-stable
pour les courbes figure dans l'article \cite{DM} de Deligne et Mumford.
Elle démontre le théorème en ramenant le problème au cas des
variétés abéliennes, déjà connu à l'époque grâce
à Grothendieck et à d'autres, voir \cite{SGA7-I} et \cite{SGA7-II}).
Le théorème de réduction semi-stable pour les variétés abéliennes utilise la théorie
des modèles de Néron, qui repose à son tour sur une étude des lois de groupe
birationnelles sur une base.

\medskip\noindent
La méthode de l'article d'Artin et Winters repose sur la désingularisation
des surfaces pour obtenir des modèles réguliers. Une fois acquise l'existence
de modèles réguliers, la démonstration consiste à analyser les
possibilités pour la fibre spéciale, puis à conclure au moyen d'une inégalité
sur la torsion du groupe de Picard d'un schéma de dimension $1$ sur un corps.
Un argument analogue figure dans un article \cite{Saito} de Saito, qui utilise
directement la cohomologie étale et obtient un résultat plus fort, puisqu'il
peut caractériser la réduction semi-stable en termes de l'action de
l'inertie sur la cohomologie étale $\ell$-adique.

\medskip\noindent
Une autre approche pour démontrer le théorème consiste à utiliser
les techniques de géométrie analytique rigide. Nous renvoyons ici le lecteur à
\cite{vanderPut} et \cite{Arzdorf-Wewers}.

\medskip\noindent
L'article \cite{Temkin} de Temkin utilise des techniques de théorie des valuations
(et démontre bien d'autres résultats) ; son appendice A donne également
un bon aperçu des différentes démonstrations et de leur rapport avec
les désingularisations des schémas de dimension $2$.

\medskip\noindent
Le lecteur pourra également consulter l'article de synthèse
\cite{Abbes-ssr} d'Ahmed Abbes.






\section{Algèbre linéaire}
\label{section-linear-algebra}

\noindent
Voici quelques lemmes que nous utiliserons par la suite.

\begin{lemma}
\label{lemma-recurring}
\begin{reference}
\cite[Theorem I]{Taussky}
\end{reference}
Soit $A = (a_{ij})$ une matrice complexe $n \times n$.
\begin{enumerate}
\item Si $|a_{ii}| > \sum_{j \not = i} |a_{ij}|$ pour tout $i$, alors
$\det(A)$ est non nul.
\item S'il existe un vecteur réel $m = (m_1, \ldots, m_n)$
avec $m_i > 0$ tel que $|a_{ii} m_i| > \sum_{j \not = i} |a_{ij}m_j|$
pour tout $i$, alors $\det(A)$ est non nul.
\end{enumerate}
\end{lemma}

\begin{proof}
Si $A$ est comme en (1) et $\det(A) = 0$, il existe un vecteur non nul
$z$ tel que $Az = 0$. Choisissons $r$ tel que $|z_r|$ soit maximal. Alors
$$
|a_{rr} z_r| = |\sum\nolimits_{k \not = r} a_{rk}z_k| \leq
\sum\nolimits_{k \not = r} |a_{rk}||z_k| \leq
|z_r| \sum\nolimits_{k \not = r} |a_{rk}| < |a_{rr}||z_r|
$$
ce qui est contradictoire. Pour démontrer (2), appliquons (1) à la matrice
$(a_{ij}m_j)$, dont le déterminant est $m_1 \ldots m_n \det(A)$.
\end{proof}

\begin{lemma}
\label{lemma-recurring-real}
Soit $A = (a_{ij})$ une matrice réelle $n \times n$ telle que
$a_{ij} \geq 0$ pour $i \not = j$. Soit $m = (m_1, \ldots, m_n)$ un vecteur réel
avec $m_i > 0$. Pour $I \subset \{1, \ldots, n\}$, soit
$x_I \in \mathbf{R}^n$
le vecteur dont la coordonnée d'indice $i$ est $m_i$ si $i \in I$
et $0$ sinon. Si
\begin{equation}
\label{equation-ineq}
-a_{ii}m_i \geq \sum\nolimits_{j \not = i} a_{ij}m_j
\end{equation}
pour tout $i$, alors $\Ker(A)$ est l'espace vectoriel
engendré par les vecteurs $x_I$ tels que
\begin{enumerate}
\item $a_{ij} = 0$ pour $i \in I$, $j \not \in I$, et
\item on a l'égalité dans (\ref{equation-ineq}) pour $i \in I$.
\end{enumerate}
\end{lemma}

\begin{proof}
En remplaçant $a_{ij}$ par $a_{ij}m_j$, on peut supposer $m_i = 1$ pour tout $i$.
Si $I \subset \{1, \ldots, n\}$ est tel que (1) et (2) soient vérifiées,
un calcul immédiat montre que $x_I$ appartient au noyau de $A$.
Réciproquement, soit $x = (x_1, \ldots, x_n) \in \mathbf{R}^n$ un
vecteur non nul du noyau de $A$. Nous allons montrer, par récurrence
sur le nombre de coordonnées non nulles de $x$, que $x$ appartient au
sous-espace engendré par les vecteurs $x_I$ vérifiant (1) et (2). Soit
$I \subset \{1, \ldots, n\}$ l'ensemble des indices $r$ pour lesquels $|x_r|$ est maximal.
Pour $r \in I$, on a
$$
|a_{rr} x_r| = |\sum\nolimits_{k \not = r} a_{rk}x_k| \leq
\sum\nolimits_{k \not = r} a_{rk}|x_k| \leq
|x_r| \sum\nolimits_{k \not = r} a_{rk} \leq |a_{rr}||x_r|
$$
On a donc l'égalité partout. En particulier, on voit que
$a_{rk} = 0$ si $r \in I$, $k \not \in I$, et que l'égalité est vérifiée
dans (\ref{equation-ineq}) pour $r \in I$. On voit alors que l'on
peut retrancher à $x$ un multiple convenable de $x_I$ pour diminuer
le nombre de coordonnées non nulles.
\end{proof}

\begin{lemma}
\label{lemma-recurring-symmetric-real}
Soit $A = (a_{ij})$ une matrice réelle symétrique $n \times n$ telle que
$a_{ij} \geq 0$ pour $i \not = j$.
Soit $m = (m_1, \ldots, m_n)$ un vecteur réel avec $m_i > 0$.
Supposons que
\begin{enumerate}
\item $Am = 0$,
\item il n'existe aucune partie propre non vide $I \subset \{1, \ldots, n\}$
telle que $a_{ij} = 0$ pour $i \in I$ et $j \not \in I$.
\end{enumerate}
Alors $x^t A x \leq 0$, avec égalité si et seulement si $x = qm$
pour un certain $q \in \mathbf{R}$.
\end{lemma}

\begin{proof}[Première démonstration]
En remplaçant $a_{ij}$ par $a_{ij}m_im_j$, on peut supposer $m_i = 1$
pour tout $i$. La condition (1) signifie que $-a_{ii} = \sum_{j \not = i} a_{ij}$
pour tout $i$. Rappelons que $x^tAx = \sum_{i, j} x_ia_{ij}x_j$.
Alors
\begin{align*}
\sum\nolimits_{i \not = j} -a_{ij}(x_j - x_i)^2 & =
\sum\nolimits_{i \not = j} -a_{ij}x_j^2 + 2a_{ij}x_ix_i - a_{ij}x_i^2 \\
& =
\sum\nolimits_j a_{jj} x_j^2 +
\sum\nolimits_{i \not = j} 2a_{ij}x_ix_i +
\sum\nolimits_j a_{jj} x_i^2 \\
& = 2x^tAx
\end{align*}
Cette expression est manifestement $\leq 0$. S'il y a égalité, soit $I$ l'ensemble
des indices $i$ tels que $x_i \not = x_1$. Alors $a_{ij} = 0$ pour $i \in I$
et $j \not \in I$. Ainsi $I = \{1, \ldots, n\}$ par la condition (2), et
$x$ est un multiple de $m = (1, \ldots, 1)$.
\end{proof}

\begin{proof}[Seconde démonstration]
La matrice $A$ a des valeurs propres réelles, d'après le théorème spectral.
Nous affirmons que toutes ses valeurs propres sont $\leq 0$.
En effet, puisque la propriété (1) signifie que
$-a_{ii}m_i = \sum_{j \not = i} a_{ij}m_j$ pour tout $i$,
la matrice $A' = A - \lambda I$, pour $\lambda > 0$,
vérifie $|a'_{ii}m_i| > \sum a'_{ij}m_j = \sum |a'_{ij}m_j|$ pour tout $i$.
Ainsi $A'$ est inversible par le lemme \ref{lemma-recurring}.
Il en résulte que la forme bilinéaire symétrique $x^tAy$
est semi-définie négative, c'est-à-dire que $x^tAx \leq 0$ pour tout $x$.
Le noyau de $A$ est donc égal
à l'ensemble des vecteurs $x$ tels que $x^tAx = 0$.
La description du noyau dans le lemme \ref{lemma-recurring-real}
donne la dernière assertion du lemme.
\end{proof}

\begin{lemma}
\label{lemma-orthogonal-direct-sum}
Soit $L$ un $\mathbf{Z}$-module libre de type fini muni
d'une forme bilinéaire symétrique entière définie positive
$\langle\ ,\ \rangle : L \times L \to \mathbf{Z}$.
Soit $A \subset L$ un sous-module tel que $L/A$ soit sans torsion. Posons
$B = \{b \in L \mid \langle a, b\rangle = 0,\ \forall a \in A\}$.
On dispose alors d'applications injectives
$$
A^\#/A \leftarrow L/(A \oplus B) \rightarrow B^\#/B
$$
dont les conoyaux sont des quotients de $L^\#/L$. Ici
$A^\# = \{a' \in A \otimes \mathbf{Q} \mid
\langle a, a'\rangle \in \mathbf{Z},\ \forall a \in A\}$
et de même pour $B$ et $L$.
\end{lemma}

\begin{proof}
Observons que
$L \otimes \mathbf{Q} = A \otimes \mathbf{Q} \oplus B \otimes \mathbf{Q}$
car la forme est non dégénérée sur $A$ (par positivité).
On note $\pi_B : L \otimes \mathbf{Q} \to B \otimes \mathbf{Q}$
la projection. Observons que $\pi_B(x) \in B^\#$ pour $x \in L$,
car la forme est entière. On obtient ainsi une suite exacte
$$
0 \to A \to L \xrightarrow{\pi_B} B^\# \to Q \to 0
$$
où $Q$ est le conoyau de $L \to B^\#$. Observons que $Q$
est un quotient de $L^\#/L$, car l'application $L^\# \to B^\#$ est surjective,
étant la duale $\mathbf{Z}$-linéaire de $B \to L$, qui est scindée
comme application de $\mathbf{Z}$-modules.
En quotientant par $A \oplus B$, on obtient une suite exacte courte
$$
0 \to L/(A \oplus B) \to B^\#/B \to Q \to 0
$$
Ceci démontre le lemme.
\end{proof}

\begin{lemma}
\label{lemma-coker}
Soient $L_0$, $L_1$ des $\mathbf{Z}$-modules libres de type fini munis
de formes bilinéaires symétriques entières définies positives
$\langle\ ,\ \rangle : L_i \times L_i \to \mathbf{Z}$.
Soient $\text{d} : L_0 \to L_1$ et $\text{d}^* : L_1 \to L_0$
des applications adjointes. Si $\langle\ ,\ \rangle$ sur $L_0$ est unimodulaire,
il existe un isomorphisme
$$
\Phi :
\Coker(\text{d}^*\text{d})_{torsion}
\longrightarrow
\Im(\text{d})^\#/\Im(\text{d})
$$
avec les notations du lemme \ref{lemma-orthogonal-direct-sum}.
\end{lemma}

\begin{proof}
Soit $x \in L_0$ un élément représentant une classe de torsion
dans $\Coker(\text{d}^*\text{d})$.
Pour un certain $a > 0$, on peut donc écrire $ax = \text{d}^*\text{d}(y)$.
Pour tout $z \in \Im(\text{d})$, disons $z = \text{d}(y')$, on a
$$
\langle (1/a)\text{d}(y), z \rangle =
\langle (1/a)\text{d}(y), \text{d}(y') \rangle =
\langle x, y' \rangle \in \mathbf{Z}
$$
Ainsi $(1/a)\text{d}(y) \in \Im(\text{d})^\#$. On définit
$\Phi(x) = (1/a)\text{d}(y) \bmod \Im(\text{d})$.
Nous omettons la preuve que $\Phi$ est bien définie, additive et injective.

\medskip\noindent
Pour montrer que $\Phi$ est surjective, soit $z \in \Im(\text{d})^\#$.
Alors $z$ définit une application linéaire $L_0 \to \mathbf{Z}$
par la règle $x \mapsto \langle z, \text{d}(x)\rangle$.
L'accouplement sur $L_0$ étant unimodulaire par hypothèse,
on peut trouver un $x' \in L_0$ tel que
$\langle x', x \rangle = \langle z, \text{d}(x)\rangle$
pour tout $x \in L_0$. En particulier, on voit
que $x'$ est orthogonal à $\Ker(\text{d})$.
Puisque $\Im(\text{d}^*\text{d}) \otimes \mathbf{Q}$
est le supplémentaire orthogonal de $\Ker(\text{d}) \otimes \mathbf{Q}$,
cela signifie que $x'$ définit une classe de torsion dans
$\Coker(\text{d}^*\text{d})$. Nous affirmons que $\Phi(x') = z$.
En effet, écrivons $a x' = \text{d}^*\text{d}(y)$
pour certains $y \in L_0$ et $a > 0$.
Pour tout $x \in L_0$, on obtient
$$
\langle z, \text{d}(x)\rangle =
\langle x', x \rangle =
\langle (1/a)\text{d}^*\text{d}(y), x \rangle =
\langle (1/a)\text{d}(y),\text{d}(x) \rangle
$$
Ainsi $z = \Phi(x')$, ce qui achève la démonstration.
\end{proof}

\begin{lemma}
\label{lemma-recurring-symmetric-integer}
Soit $A = (a_{ij})$ une matrice entière symétrique $n \times n$ telle que
$a_{ij} \geq 0$ pour $i \not = j$. Soit $m = (m_1, \ldots, m_n)$ un
vecteur entier avec $m_i > 0$. Supposons que
\begin{enumerate}
\item $Am = 0$,
\item il n'existe aucune partie propre non vide $I \subset \{1, \ldots, n\}$
telle que $a_{ij} = 0$ pour $i \in I$ et $j \not \in I$.
\end{enumerate}
Soit $e$ le nombre de couples $(i, j)$ tels que $i < j$ et $a_{ij} > 0$.
Alors, pour un nombre premier $\ell$ premier à tous les $a_{ij}$ et $m_i$,
on a
$$
\dim_{\mathbf{F}_\ell}(\Coker(A)[\ell]) \leq 1 - n + e
$$
\end{lemma}

\begin{proof}
D'après le lemme \ref{lemma-recurring-symmetric-real}, le rang de $A$ est $n - 1$.
La composée
$$
\mathbf{Z}^{\oplus n} \xrightarrow{\text{diag}(m_1, \ldots, m_n)}
\mathbf{Z}^{\oplus n} \xrightarrow{(a_{ij})}
\mathbf{Z}^{\oplus n} \xrightarrow{\text{diag}(m_1, \ldots, m_n)}
\mathbf{Z}^{\oplus n}
$$
a pour matrice $a_{ij}m_im_j$. Puisque les conoyaux de la première et de la dernière
application sont de torsion d'ordre premier à $\ell$, par notre condition sur $\ell$,
il suffit de démontrer le lemme pour la matrice
de coefficients $a_{ij}m_im_j$. On peut donc supposer $m = (1, \ldots, 1)$.

\medskip\noindent
Supposons $m = (1, \ldots, 1)$. Posons $V = \{1, \ldots, n\}$ et
$E = \{(i, j) \mid i < j\text{ et }a_{ij} > 0\}$. Pour
$e = (i, j) \in E$, posons $a_e = a_{ij}$. Définissons des applications
$s, t : E \to V$ en posant $s(i, j) = i$ et $t(i, j) = j$.
Posons
$\mathbf{Z}(V) = \bigoplus_{i \in V} \mathbf{Z}i$ et
$\mathbf{Z}(E) = \bigoplus_{e \in E} \mathbf{Z}e$.
On définit des accouplements symétriques définis positifs à valeurs entières
sur $\mathbf{Z}(V)$ et $\mathbf{Z}(E)$ en posant
$$
\langle i, i \rangle = 1\text{ pour }i \in V, \quad
\langle e, e \rangle = a_e\text{ pour }e \in E
$$
et en prenant tous les autres accouplements nuls. Considérons les applications
$$
\text{d} : \mathbf{Z}(V) \to \mathbf{Z}(E), \quad
i \longmapsto
\sum\nolimits_{e \in E,\ s(e) = i} e - \sum\nolimits_{e \in E,\ t(e) = i} e
$$
et
$$
\text{d}^*(e) = a_e(s(e) - t(e))
$$
Un calcul montre que
$$
\langle d(x), y\rangle = \langle x, \text{d}^*(y) \rangle
$$
autrement dit, $\text{d}$ et $\text{d}^*$ sont adjointes. On calcule ensuite
\begin{align*}
\text{d}^*\text{d}(i)
& = 
\text{d}^*(
\sum\nolimits_{e \in E,\ s(e) = i} e - \sum\nolimits_{e \in E,\ t(e) = i} e) \\
& =
\sum\nolimits_{e \in E,\ s(e) = i} a_e(s(e) - t(e)) -
\sum\nolimits_{e \in E,\ t(e) = i} a_e(s(e) - t(e))
\end{align*}
Le coefficient de $i$ dans $\text{d}^*\text{d}(i)$ est
$$
\sum\nolimits_{e \in E,\ s(e) = i} a_e +
\sum\nolimits_{e \in E,\ t(e) = i} a_e = - a_{ii}
$$
car $\sum_j a_{ij} = 0$, et le coefficient de
$j \not = i$ dans $\text{d}^*\text{d}(i)$ est $-a_{ij}$.
Ainsi $\Coker(A) = \Coker(\text{d}^*\text{d})$.

\medskip\noindent
Considérons l'inclusion
$$
\Im(\text{d}) \oplus \Ker(\text{d}^*) \subset \mathbf{Z}(E)
$$
Le membre de gauche est une somme directe orthogonale. Il est clair que
$\mathbf{Z}(E)/\Ker(\text{d}^*)$ est sans torsion.
Nous affirmons que $\mathbf{Z}(E)/\Im(\text{d})$ est également sans torsion.
En effet, soient $x = \sum x_e e \in \mathbf{Z}(E)$ et $a > 1$ tels
que $ax = \text{d}y$ pour un certain $y = \sum y_i i \in \mathbf{Z}(V)$.
Alors $a x_e = y_{s(e)} - y_{t(e)}$. La propriété (2) entraîne
que tous les $y_i$ ont la même classe de congruence modulo $a$.
On peut donc écrire $y = a y' + (y_1, y_1, \ldots, y_1)$.
Puisque $\text{d}(y_1, y_1, \ldots, y_1) = 0$, on en déduit
que $x = \text{d}(y')$, ce qu'il fallait démontrer.

\medskip\noindent
On peut donc appliquer le lemme \ref{lemma-orthogonal-direct-sum}
pour obtenir des applications injectives
$$
\Im(\text{d})^\#/\Im(\text{d}) \leftarrow
\mathbf{Z}(E)/(\Im(\text{d}) \oplus \Ker(\text{d}^*)) \rightarrow
\Ker(\text{d}^*)^\#/\Ker(\text{d}^*)
$$
dont les conoyaux sont annulés par le produit des $a_e$
(qui est premier à $\ell$). Puisque $\Ker(\text{d}^*)$ est
un réseau de rang $1 - n + e$, il suffit, pour achever la démonstration,
de montrer qu'il existe un isomorphisme
$$
\Phi : M_{torsion} \longrightarrow \Im(\text{d})^\#/\Im(\text{d})
$$
Ceci est démontré dans le lemme \ref{lemma-coker}.
\end{proof}









\section{Types numériques}
\label{section-numerical-types}

\noindent
Une partie des arguments fera intervenir la combinatoire des données suivantes.

\begin{definition}
\label{definition-type}
Un {\it type numérique} $T$ est donné par
$$
n, m_i, a_{ij}, w_i, g_i
$$
où $n \geq 1$ est un entier et où $m_i$, $a_{ij}$, $w_i$, $g_i$
sont des entiers, pour $1 \leq i, j \leq n$, soumis aux conditions suivantes :
\begin{enumerate}
\item $m_i > 0$, $w_i > 0$, $g_i \geq 0$,
\item la matrice $A = (a_{ij})$ est symétrique et $a_{ij} \geq 0$
pour $i \not = j$,
\item il n'existe aucune partie propre non vide $I \subset \{1, \ldots, n\}$
telle que $a_{ij} = 0$ pour $i \in I$, $j \not \in I$,
\item pour tout $i$, on a $\sum_j a_{ij}m_j = 0$, et
\item $w_i | a_{ij}$.
\end{enumerate}
\end{definition}

\noindent
Cette sorte de structure est certes quelque peu malcommode à manipuler,
mais c'est exactement ce qui apparaît dans les fibres spéciales des modèles
propres réguliers des courbes lisses géométriquement connexes.
Naturellement, ces types ne nous intéressent qu'à permutation des indices près.

\begin{definition}
\label{definition-type-equivalent}
On dit que deux types numériques $n, m_i, a_{ij}, w_i, g_i$ et
$n', m'_i, a'_{ij}, w'_i, g'_i$ sont des {\it types équivalents} s'il
existe une permutation $\sigma$ de $\{1, \ldots, n\}$
telle que $m_i = m'_{\sigma(i)}$, $a_{ij} = a'_{\sigma(i)\sigma(j)}$,
$w_i = w'_{\sigma(i)}$ et $g_i = g'_{\sigma(i)}$.
\end{definition}

\noindent
Un type numérique possède un genre.

\begin{lemma}
\label{lemma-genus}
Soit $n, m_i, a_{ij}, w_i, g_i$ un type numérique. Alors l'expression
$$
g = 1 + \sum m_i(w_i(g_i - 1) - \frac{1}{2} a_{ii})
$$
est un entier.
\end{lemma}

\begin{proof}
Pour montrer que $g$ est un entier, il faut montrer que $\sum a_{ii}m_i$ est pair.
On le voit en calculant modulo $2$ comme suit :
\begin{align*}
\sum\nolimits_i a_{ii} m_i 
& \equiv
\sum\nolimits_{i,\ m_i\text{ impair}} a_{ii}m_i \\
& \equiv
\sum\nolimits_{i,\ m_i\text{ impair}} \sum\nolimits_{j \not = i} a_{ij}m_j \\
& \equiv
\sum\nolimits_{i,\ m_i\text{ impair}}
\sum\nolimits_{j \not = i,\ m_j\text{ impair}} a_{ij}m_j \\
& \equiv
\sum\nolimits_{i < j,\ m_i\text{ et }m_j\text{ impairs}} a_{ij}(m_i + m_j) \\
& \equiv
0
\end{align*}
où l'on a utilisé que $a_{ij} = a_{ji}$ et que $\sum_j a_{ij}m_j = 0$
pour tout $i$.
\end{proof}

\begin{definition}
\label{definition-genus}
On dit que $n, m_i, a_{ij}, w_i, g_i$ est un {\it type numérique de genre $g$}
si $g = 1 + \sum m_i(w_i(g_i - 1) - \frac{1}{2} a_{ii})$ est l'entier
du lemme \ref{lemma-genus}.
\end{definition}

\noindent
Nous démontrerons plus bas (lemme \ref{lemma-genus-nonnegative}) que le genre
est presque toujours $\geq 0$. Il peut toutefois exister des types numériques de
genre négatif.

\begin{lemma}
\label{lemma-irreducible}
Soit $n, m_i, a_{ij}, w_i, g_i$ un type numérique de genre $g$.
Si $n = 1$, alors $a_{11} = 0$ et $g = 1 + m_1w_1(g_1 - 1)$.
De plus, on peut classifier tous ces types numériques comme suit :
\begin{enumerate}
\item Si $g < 0$, alors $g_1 = 0$ et il n'y a qu'un nombre fini de
types numériques de genre $g$ avec $n = 1$, correspondant aux factorisations
$m_1w_1 = 1 - g$.
\item Si $g = 0$, alors $m_1 = 1$, $w_1 = 1$, $g_1 = 0$,
comme dans le lemme \ref{lemma-genus-zero}.
\item Si $g = 1$, on en déduit $g_1 = 1$, mais $m_1, w_1$ peuvent être des entiers
strictement positifs arbitraires ; c'est le cas
(\ref{item-one}) du lemme \ref{lemma-genus-one}.
\item Si $g > 1$, alors $g_1 > 1$ et il n'y a qu'un nombre fini de
types numériques de genre $g$ avec $n = 1$, correspondant aux
factorisations $m_1w_1(g_1 - 1) = g - 1$.
\end{enumerate}
\end{lemma}

\begin{proof}
Le lemme est immédiat.
\end{proof}

\begin{lemma}
\label{lemma-diagonal-negative}
Soit $n, m_i, a_{ij}, w_i, g_i$ un type numérique de genre $g$.
Si $n > 1$, alors $a_{ii} < 0$ pour tout $i$.
\end{lemma}

\begin{proof}
Le lemme \ref{lemma-recurring-symmetric-real} s'applique à la matrice $A$.
\end{proof}

\begin{lemma}
\label{lemma-minus-one}
Soit $n, m_i, a_{ij}, w_i, g_i$ un type numérique de genre $g$.
Supposons $n > 1$. Si $i$ est tel que la contribution
$m_i(w_i(g_i - 1) - \frac{1}{2} a_{ii})$
au genre $g$ soit $< 0$, alors $g_i = 0$ et $a_{ii} = -w_i$.
\end{lemma}

\begin{proof}
Cela résulte immédiatement du lemme \ref{lemma-diagonal-negative} et des relations
$w_i > 0$, $g_i \geq 0$ et $w_i | a_{ii}$.
\end{proof}

\begin{definition}
\label{definition-type-minus-one}
Soit $n, m_i, a_{ij}, w_i, g_i$ un type numérique.
On dit que $i$ est un {\it indice $(-1)$} si $g_i = 0$ et $a_{ii} = -w_i$.
\end{definition}

\noindent
On peut « contracter » les indices $(-1)$.

\begin{lemma}
\label{lemma-contract}
Soit $n, m_i, a_{ij}, w_i, g_i$ un type numérique $T$.
Supposons que $n$ soit un indice $(-1)$. Il existe alors un type
numérique $T'$ donné par $n', m'_i, a'_{ij}, w'_i, g'_i$, avec
\begin{enumerate}
\item $n' = n - 1$,
\item $m'_i = m_i$,
\item $a'_{ij} = a_{ij} - a_{in}a_{jn}/a_{nn}$,
\item $w'_i = w_i/2$ si $a_{in}/w_n$ est pair et $a_{in}/w_i$ impair,
et $w'_i = w_i$ sinon,
\item $g'_i =
\frac{w_i}{w'_i}(g_i - 1) + 1 + \frac{a_{in}^2 - w_na_{in}}{2w'_iw_n}$.
\end{enumerate}
De plus, on a $g = g'$.
\end{lemma}

\begin{proof}
Observons que $n > 1$, par exemple d'après le lemme \ref{lemma-irreducible},
et donc que $n' \geq 1$. Vérifions les conditions (1) -- (5) de la
définition \ref{definition-type} pour $n', m'_i, a'_{ij}, w'_i, g'_i$.

\medskip\noindent
La condition (1) est immédiate.

\medskip\noindent
Condition (2). La symétrie de $A' = (a'_{ij})$ est immédiate
et, puisque $a_{nn} < 0$ d'après le lemme \ref{lemma-diagonal-negative},
on voit que $a'_{ij} \geq a_{ij} \geq 0$ si $i \not = j$.

\medskip\noindent
Condition (3). Supposons que $I \subset \{1, \ldots, n - 1\}$ soit tel que
$a'_{ii'} = 0$ pour $i \in I$ et $i' \in \{1, \ldots, n - 1\} \setminus I$.
On voit alors que, pour tout $i \in I$ et tout $i' \in I'$, on a
$a_{in}a_{i'n} = 0$. Ainsi, ou bien $a_{in} = 0$ pour tout $i \in I$, et
$I \subset \{1, \ldots, n\}$ contredit la propriété (3) pour $T$,
ou bien $a_{i'n} = 0$ pour tout $i' \in \{1, \ldots, n - 1\} \setminus I$,
et $I \cup \{n\} \subset \{1, \ldots, n\}$ contredit
la propriété (3) de $T$. La condition (3) est donc vérifiée pour $T'$.

\medskip\noindent
Condition (4). On calcule
$$
\sum\nolimits_{j = 1}^{n - 1} a'_{ij}m_j =
\sum\nolimits_{j = 1}^{n - 1}
(a_{ij}m_j - \frac{a_{in}a_{jn}m_j}{a_{nn}}) =
- a_{in}m_n - \frac{a_{in}}{a_{nn}}(-a_{nn}m_n) = 0
$$
comme voulu.

\medskip\noindent
Condition (5). Il faut montrer que $w'_i$ divise $a_{in}a_{jn}/a_{nn}$.
C'est clair puisque $a_{nn} = -w_n$, $w_n | a_{jn}$ et $w_i | a_{in}$.

\medskip\noindent
Pour montrer que $g = g'$, écrivons d'abord
\begin{align*}
g
& =
1 + \sum\nolimits_{i = 1}^n m_i(w_i(g_i - 1) - \frac{1}{2}a_{ii}) \\
& =
1 + \sum\nolimits_{i = 1}^{n - 1} m_i(w_i(g_i - 1) - \frac{1}{2}a_{ii})
-\frac{1}{2}m_nw_n \\
& =
1 +  \sum\nolimits_{i = 1}^{n - 1}
m_i(w_i(g_i - 1) - \frac{1}{2}a_{ii} - \frac{1}{2}a_{in})
\end{align*}
En comparant avec l'expression de $g'$, on voit qu'il suffit que
$$
w'_i(g'_i - 1) - \frac{1}{2}a'_{ii} =
w_i(g_i - 1) - \frac{1}{2}a_{in} - \frac{1}{2}a_{ii}
$$
pour $i \leq n - 1$. Autrement dit, on a
$$
g'_i = \frac{2w_i(g_i - 1) - a_{in} - a_{ii} + a'_{ii} + 2w'_i}{2w'_i} =
\frac{w_i}{w'_i}(g_i - 1) + 1 + \frac{a_{in}^2 - w_na_{in}}{2w'_iw_n}
$$
On vérifie élémentairement qu'il s'agit d'un entier $\geq 0$
si l'on choisit $w'_i$ comme en (4).
\end{proof}

\begin{lemma}
\label{lemma-top-genus}
Soit $n, m_i, a_{ij}, w_i, g_i$ un type numérique.
Soit $e$ le nombre de couples $(i, j)$ tels que $i < j$ et $a_{ij} > 0$.
Alors l'expression $g_{top} = 1 - n + e$ est $\geq 0$.
\end{lemma}

\begin{proof}
Sinon, $e < n - 1$, ce qui signifie qu'il existe un $i$ tel que
$a_{ij} = 0$ pour tout $j \not = i$. Cela contredit l'hypothèse
(3) de la définition \ref{definition-type}.
\end{proof}

\begin{definition}
\label{definition-top-genus}
Soit $n, m_i, a_{ij}, w_i, g_i$ un type numérique $T$. Le
{\it genre topologique de $T$} est l'entier positif ou nul
$g_{top} = 1 - n + e$ du lemme \ref{lemma-top-genus}.
\end{definition}

\noindent
Nous voulons minorer le genre par le genre topologique. Toutefois, cela
ne sera pas toujours possible, par exemple pour les types numériques
avec $n = 1$ du lemme \ref{lemma-irreducible}. Ce sera cependant
le cas pour les types numériques minimaux, définis comme suit.

\begin{definition}
\label{definition-type-minimal}
On dit que le type numérique $n, m_i, a_{ij}, w_i, g_i$ de genre $g$
est {\it minimal} s'il n'existe aucun $i$
tel que $g_i = 0$ et $a_{ii} = -w_i$, autrement dit, s'il
n'existe aucun indice $(-1)$.
\end{definition}

\noindent
Nous démontrerons que le genre $g$ d'un type minimal avec $n > 1$
est supérieur ou égal à $\max(1, g_{top})$.

\begin{lemma}
\label{lemma-non-irreducible-minimal-type-genus-at-least-one}
Si $n, m_i, a_{ij}, w_i, g_i$ est un type numérique minimal
avec $n > 1$, alors $g \geq 1$.
\end{lemma}

\begin{proof}
Cela résulte de ce que $g = 1 + \sum \Phi_i$, avec
$\Phi_i = m_i(w_i(g_i - 1) - \frac{1}{2} a_{ii})$ positif ou nul,
d'après le lemme \ref{lemma-minus-one} et la définition des types minimaux.
\end{proof}

\begin{lemma}
\label{lemma-genus-nonnegative}
Si $n, m_i, a_{ij}, w_i, g_i$ est un type numérique minimal
avec $n > 1$, alors $g \geq g_{top}$.
\end{lemma}

\begin{proof}
Le lecteur qui ne s'intéresse qu'aux types numériques
associés aux modèles propres réguliers peut omettre cette démonstration, car nous
redémontrerons le résultat plus loin dans la situation géométrique.
On peut écrire
$$
g_{top} = 1 - n + \frac{1}{2}\sum\nolimits_{a_{ij} > 0} 1 =
1 + \sum\nolimits_i (-1 +
\frac{1}{2}\sum\nolimits_{j \not = i,\ a_{ij} > 0} 1) 
$$
D'autre part, on a
\begin{align*}
g & =
1 + \sum m_i(w_i(g_i - 1) - \frac{1}{2} a_{ii}) \\
& =
1 + \sum m_iw_ig_i - \sum m_iw_i +
\frac{1}{2} \sum\nolimits_{i \not = j} a_{ij}m_j \\
& =
1 + \sum\nolimits_i
m_iw_i(-1 + g_i + \frac{1}{2} \sum\nolimits_{j \not = i} \frac{a_{ij}}{w_i})
\end{align*}
La première égalité est la définition, la deuxième utilise que
$\sum a_{ij}m_j = 0$, et la dernière utilise
la relation $a_{ij} = a_{ji}$ ainsi que l'interversion de l'ordre
des sommations. En comparant avec la formule pour $g_{top}$, on conclut
que le lemme est vrai si
$$
\Psi_i =
m_iw_i(-1 + g_i + \frac{1}{2} \sum\nolimits_{j \not = i} \frac{a_{ij}}{w_i})
- (-1 + \frac{1}{2}\sum\nolimits_{j \not = i,\ a_{ij} > 0} 1)
$$
est $\geq 0$ pour tout $i$. Ce n'est cependant pas toujours le cas.
Examinons pour quels indices on peut avoir $\Psi_i < 0$.
Observons d'abord que
$$
(-1 + g_i + \frac{1}{2}\sum\nolimits_{j \not = i} \frac{a_{ij}}{w_i}) \geq
(-1 + \frac{1}{2}\sum\nolimits_{j \not = i,\ a_{ij} > 0} 1)
$$
car $a_{ij}/w_i$ est un entier positif ou nul. Puisque $m_iw_i$ est
un entier strictement positif, on en déduit que $\Psi_i \geq 0$ dès que
$m_iw_i = 1$ ou que le membre de gauche de l'inégalité est $\geq 0$,
ce qui se produit si $g_i > 0$, ou si $a_{ij} > 0$ pour au moins deux indices $j$, ou
s'il existe un $j$ tel que $a_{ij} > w_i$. Ainsi
$$
P = \{i : \Psi_i < 0\}
$$
est l'ensemble des indices $i$ tels que $m_iw_i > 1$, $g_i = 0$,
$a_{ij} > 0$ pour un unique $j$ et $a_{ij} = w_i$ pour ce $j$.
De plus,
$$
i \in P \Rightarrow \Psi_i = \frac{1}{2}(-m_iw_i + 1)
$$
La stratégie consiste à montrer que, pour $i \in P$, on peut prélever une part
de $\Psi_j$, où $j$ est le voisin de $i$, c'est-à-dire où $a_{ij} > 0$.
Toutefois, cela ne suffit pas tout à fait, car $j$ peut être un indice tel que $\Psi_j = 0$.

\medskip\noindent
Considérons l'ensemble
$$
Z = \{j : g_j = 0\text{ et }
j\text{ a exactement deux voisins }i, k\text{ avec }
a_{ij} = w_j = a_{jk}\}
$$
Pour $j \in Z$, on a $\Psi_j = 0$. Nous considérerons des suites
$M = (i, j_1, \ldots, j_s)$ où $s \geq 0$,
$i \in P$, $j_1, \ldots, j_s \in Z$, et
$a_{ij_1} > 0, a_{j_1j_2} > 0, \ldots, a_{j_{s - 1}j_s} > 0$.
Si notre type numérique est constitué de deux indices appartenant à $P$,
ou, plus généralement, si notre type numérique est constitué de
deux indices appartenant à $P$ et de tous les autres indices dans $Z$, alors
$g_{top} = 0$ et on conclut par le
lemme \ref{lemma-non-irreducible-minimal-type-genus-at-least-one}.
On peut donc écarter ces cas, ce que nous faisons.

\medskip\noindent
Soit $M = (i, j_1, \ldots, j_s)$ une suite maximale, et soit
$k$ le second voisin de $j_s$. (Si $s = 0$, alors $k$
est l'unique voisin de $i$.) Par maximalité, $k \not \in Z$
et, d'après ce qui précède, $k \not \in P$. Observons que
$w_i = a_{ij_1} = w_{j_1} = a_{j_1j_2} = \ldots = w_{j_s} =
a_{j_sk}$. En examinant la définition
d'un type numérique, on voit que
\begin{align*}
m_ia_{ii} + m_{j_1}w_i & = 0,\\
m_iw_i + m_{j_1}a_{j_1j_1} + m_{j_2}w_i & = 0,\\
\ldots & \ldots \\
m_{j_{s - 1}}w_i + m_{j_s}a_{j_sj_s} + m_kw_i & = 0
\end{align*}
La première égalité entraîne $m_{j_1} \geq 2m_i$, car le
type numérique est minimal. La deuxième égalité entraîne alors
$m_{j_2} \geq 3m_i$, et ainsi de suite. Dans tous les cas, on en déduit que
$m_k \geq 2m_i$ (y compris lorsque $s = 0$).

\medskip\noindent
Soit $k$ un indice pour lequel on dispose d'un $t > 0$ et de suites
maximales deux à deux distinctes $M_1, \ldots, M_t$ comme ci-dessus, avec
$M_b = (i_b, j_{b, 1}, \ldots, j_{b, s_b})$
telles que $k$ soit un voisin de $j_{b, s_b}$ pour $b = 1, \ldots, t$.
Nous montrerons que $\Phi_j + \sum_{b = 1, \ldots, t} \Phi_{i_b} \geq 0$.
Cela achèvera la démonstration du lemme, d'après ce qui précède.
Soit $M$ la réunion des indices figurant dans $M_b$, $b = 1, \ldots, t$.
Écrivons
$$
\Psi_k =
-\sum\nolimits_{b = 1, \ldots, t} \Psi_{i_b} + \Psi_k'
$$
où
\begin{align*}
\Psi_k' & =
m_kw_k\left(-1 + g_k +
\frac{1}{2} \sum\nolimits_{b = 1, \ldots t}
(\frac{a_{kj_{b, s_b}}}{w_k} - \frac{m_{i_b}w_{i_b}}{m_kw_k}) +
\frac{1}{2} \sum\nolimits_{l \not = k,\ l \not \in M}
\frac{a_{kl}}{w_k}
\right) \\
&
-\left(
-1 + \frac{1}{2}\sum\nolimits_{l \not = k,\ l \not \in M,\ a_{kl} > 0} 1
\right)
\end{align*}
Supposons $\Psi_k' < 0$ pour obtenir une contradiction.
Si l'ensemble $\{l : l \not = k,\ l \not \in M,\ a_{kl} > 0\}$ est vide,
alors $\{1, \ldots, n\} = M \cup \{k\}$ et $g_{top} = 0$,
car $e = n - 1$ dans ce cas, et le résultat découle du
lemme \ref{lemma-non-irreducible-minimal-type-genus-at-least-one}.
On peut donc supposer qu'il existe au moins un tel $l$, qui contribue
pour $(1/2)a_{kl}/w_k \geq 1/2$ à la somme entre les premières parenthèses.
Pour tout $b = 1, \ldots, t$, on a
$$
\frac{a_{kj_{b, s_b}}}{w_k} - \frac{m_{i_b}w_{i_b}}{m_kw_k} =
\frac{w_{i_b}}{w_k}(1 - \frac{m_{i_b}}{m_k})
$$
Cette expression est $\geq \frac{1}{2}$, car $m_k \geq 2m_{i_b}$
d'après le paragraphe précédent, et elle est $\geq 1$ si $w_k < w_{i_b}$.
Il s'ensuit que $\Psi_k' < 0$ entraîne $g_k = 0$.
Si $t \geq 2$, ou si $t = 1$ et $w_k < w_{i_1}$, alors $\Psi_k' \geq 0$
(on utilise ici l'existence d'un $l$ établie ci-dessus), ce qui
est encore contradictoire.
Ainsi $t = 1$ et $w_k = w_{i_1}$. S'il y a au moins deux termes non nuls
dans la somme sur $l$, ou s'il existe un tel $k$ et que $a_{kl} > w_k$, alors
on a également $\Psi_k' \geq 0$. La dernière possibilité est que $t = 1$ et
qu'il existe un $l$ tel que $a_{kl} = w_k$. Ceci est exclu, car cela
signifierait que $k \in Z$, contrairement à la maximalité de $M_1$.
\end{proof}

\begin{lemma}
\label{lemma-minus-two}
Soit $n, m_i, a_{ij}, w_i, g_i$ un type numérique de genre $g$.
Supposons $n > 1$. Si $i$ est tel que la contribution
$m_i(w_i(g_i - 1) - \frac{1}{2} a_{ii})$
au genre $g$ soit $0$, alors $g_i = 0$ et $a_{ii} = -2w_i$.
\end{lemma}

\begin{proof}
Cela résulte immédiatement du lemme \ref{lemma-diagonal-negative} et des relations
$w_i > 0$, $g_i \geq 0$ et $w_i | a_{ii}$.
\end{proof}

\noindent
Les indices vérifiant cette relation jouent en fait un
rôle important dans la structure des types numériques minimaux.
Nous leur donnons donc un nom.

\begin{definition}
\label{definition-type-minus-two}
Soit $n, m_i, a_{ij}, w_i, g_i$ un type numérique de genre $g$.
On dit que $i$ est un {\it indice $(-2)$} si $g_i = 0$ et $a_{ii} = -2w_i$.
\end{definition}

\noindent
Étant donné un type numérique minimal de genre $g$, les indices $(-2)$
sont exactement ceux dont la contribution au genre n'est pas strictement positive
dans la formule
$$
g = 1 + \sum m_i(w_i(g_i - 1) - \frac{1}{2} a_{ii})
$$
Il sera donc quelque peu délicat de borner les quantités associées
aux indices $(-2)$, comme nous le verrons plus loin.

\begin{remark}
\label{remark-genus-equality}
Soit $n, m_i, a_{ij}, w_i, g_i$ un type numérique minimal avec $n > 1$.
L'égalité $g = g_{top}$ peut avoir lieu dans le lemme \ref{lemma-genus-nonnegative}.
Par exemple, si $m_i = w_i = 1$ et $g_i = 0$ pour tout $i$, et si
$a_{ij} \in \{0, 1\}$ pour $i < j$.
\end{remark}


\section{Le groupe de Picard d'un type numérique}
\label{section-picard-group}

\noindent
Voici la définition.

\begin{definition}
\label{definition-picard-group}
Soit $n, m_i, a_{ij}, w_i, g_i$ un type numérique $T$. Le
{\it groupe de Picard de $T$} est le conoyau de la matrice
$(a_{ij}/w_i)$, plus précisément
$$
\Pic(T) =
\Coker\left(
\mathbf{Z}^{\oplus n} \to \mathbf{Z}^{\oplus n},\quad
e_i
\mapsto
\sum \frac{a_{ij}}{w_j}e_j
\right)
$$
où $e_i$ désigne le vecteur d'indice $i$ de la base canonique de $\mathbf{Z}^{\oplus n}$.
\end{definition}

\begin{lemma}
\label{lemma-picard-rank-1}
Soit $n, m_i, a_{ij}, w_i, g_i$ un type numérique $T$.
Le groupe de Picard de $T$ est un groupe abélien de type fini de rang $1$.
\end{lemma}

\begin{proof}
Si $n = 1$, alors $A = (a_{ij})$ est la matrice nulle et
le résultat est clair. Pour $n > 1$, la matrice $A$ est de rang
$n - 1$ d'après le lemme \ref{lemma-recurring-real} ou le
lemme \ref{lemma-recurring-symmetric-real}.
Bien entendu, le rang ne change pas lorsqu'on multiplie les lignes
par $1/w_i$. Ceci démontre le lemme.
\end{proof}

\begin{lemma}
\label{lemma-picard-T-and-A}
Soit $n, m_i, a_{ij}, w_i, g_i$ un type numérique $T$.
Alors $\Pic(T) \subset \Coker(A)$, où $A = (a_{ij})$.
\end{lemma}

\begin{proof}
Puisque $\Pic(T)$ est le conoyau de $(a_{ij}/w_i)$,
on dispose d'un diagramme commutatif
$$
\xymatrix{
0 \ar[r] &
\mathbf{Z}^{\oplus n} \ar[rr]_A & &
\mathbf{Z}^{\oplus n} \ar[rr] & &
\Coker(A) \ar[r] & 0 \\
0 \ar[r] &
\mathbf{Z}^{\oplus n} \ar[rr]^{(a_{ij}/w_i)} \ar[u]_{\text{id}} & &
\mathbf{Z}^{\oplus n} \ar[rr] \ar[u]_{\text{diag}(w_1, \ldots, w_n)} & &
\Pic(T) \ar[r] \ar[u] & 0
}
$$
dont les lignes sont exactes. Le lemme du serpent entraîne que
$\Pic(T) \subset \Coker(A)$.
\end{proof}

\begin{lemma}
\label{lemma-contract-picard-group}
Soit $n, m_i, a_{ij}, w_i, g_i$ un type numérique $T$.
Supposons que $n$ soit un indice $(-1)$. Soit $T'$ le type numérique
construit dans le lemme \ref{lemma-contract}. Il existe une
application injective
$$
\Pic(T) \to \Pic(T')
$$
dont le conoyau est un $2$-groupe abélien élémentaire.
\end{lemma}

\begin{proof}
Rappelons que $n' = n - 1$. Soit $e_i$, resp. $e'_i$, le vecteur d'indice $i$ de la
base de $\mathbf{Z}^{\oplus n}$, resp.\ $\mathbf{Z}^{\oplus n - 1}$.
Notons d'abord
$$
q : \mathbf{Z}^{\oplus n} \to \mathbf{Z}^{\oplus n - 1},
\quad e_n \mapsto 0\text{ et }e_i \mapsto e'_i\text{ pour }i \leq n - 1
$$
et posons
$$
p : \mathbf{Z}^{\oplus n} \to \mathbf{Z}^{\oplus n - 1},\quad
e_n \mapsto \sum\nolimits_{j = 1}^{n - 1} \frac{a_{nj}}{w'_j} e'_j
\text{ et }
e_i \mapsto \frac{w_i}{w'_i} e'_i\text{ pour }i \leq n - 1
$$
Un calcul (que nous omettons) montre que l'on a un diagramme commutatif
$$
\xymatrix{
\mathbf{Z}^{\oplus n} \ar[rr]_{(a_{ij}/w_i)} \ar[d]_q & &
\mathbf{Z}^{\oplus n} \ar[d]^p \\
\mathbf{Z}^{\oplus n'} \ar[rr]^{(a'_{ij}/w'_i)} & &
\mathbf{Z}^{\oplus n'}
}
$$
Le conoyau de la flèche supérieure étant
$\Pic(T)$ et celui de la flèche inférieure
étant $\Pic(T')$, on obtient l'homomorphisme voulu
entre les groupes de Picard. Puisque $\frac{w_i}{w'_i} \in \{1, 2\}$,
le conoyau de $\Pic(T) \to \Pic(T')$
est annulé par $2$ (car $2e'_i$ appartient à l'image de $p$
pour tout $i \leq n - 1$).
Montrons enfin que $\Pic(T) \to \Pic(T')$ est injective.
Soit $L = (l_1, \ldots, l_n)$ un représentant
d'un élément de $\Pic(T)$ d'image nulle dans $\Pic(T')$.
Puisque $q$ est surjective, une chasse au diagramme montre que l'on peut supposer
$L$ dans le noyau de $p$. Cela signifie que
$l_na_{ni}/w'_i + l_iw_i/w'_i = 0$, c'est-à-dire $l_i = - a_{ni}/w_i l_n$.
Ainsi $L$ est l'image de $-l_ne_n$ par l'application $(a_{ij}/w_j)$,
ce qui démontre le lemme.
\end{proof}

\begin{lemma}
\label{lemma-picard-group-genus-nonpositive}
Soit $n, m_i, a_{ij}, w_i, g_i$ un type numérique $T$.
Si le genre $g$ de $T$ est $\leq 0$, alors $\Pic(T) = \mathbf{Z}$.
\end{lemma}

\begin{proof}
On raisonne par récurrence sur $n$. Si $n = 1$, l'assertion est claire.
Si $n > 1$, alors $T$ n'est pas minimal d'après le
lemme \ref{lemma-non-irreducible-minimal-type-genus-at-least-one}.
En remplaçant $T$ par un type équivalent,
on peut supposer que $n$ est un indice $(-1)$.
D'après le lemme \ref{lemma-contract-picard-group},
on a $\Pic(T) \subset \Pic(T')$.
D'après le lemme \ref{lemma-contract}, le genre
de $T'$ est égal à celui de $T$, et l'on conclut par
récurrence.
\end{proof}








\section{Classification des sous-graphes propres}
\label{section-classify-proper-subgraphs}

\noindent
Dans cette section, on suppose donné un type numérique
$n, m_i, a_{ij}, w_i, g_i$ de genre $g$. Nous dresserons
une liste complète des « sous-graphes » possibles constitués entièrement
d'indices $(-2)$ (définition \ref{definition-type-minus-two}) et,
en même temps, nous classifierons tous les types numériques
minimaux possibles de genre $1$. Autrement dit, nous démontrons ici la
proposition \ref{proposition-classify-subgraphs} et le
lemme \ref{lemma-genus-one}

\medskip\noindent
Notre stratégie sera la suivante. Soit $n, m_i, a_{ij}, w_i, g_i$
un type numérique de genre $g$. Soit $I \subset \{1, \ldots, n\}$ une partie
constituée d'indices $(-2)$ telle qu'il n'existe aucune partie propre non vide
$J \subset I$ avec $a_{jj'} = 0$ pour $j \in J$, $j' \in I \setminus J$.
On raisonne par récurrence sur le cardinal $|I|$ de $I$. Si $I = \{i\}$
est constitué de $1$ indice, les seules contraintes sur $m_i$, $a_{ii}$
et $w_i$ sont $w_i | a_{ii}$, d'après la définition \ref{definition-type},
et $a_{ii} < 0$, d'après le lemme \ref{lemma-diagonal-negative},
ce qui fournira le cas initial. Dans l'étape de récurrence,
on applique d'abord l'hypothèse de récurrence aux parties $I' \subset I$
de cardinal $|I'| < |I|$. On obtient ainsi des contraintes sur les valeurs possibles
de $m_i, a_{ij}, w_i$, $i, j \in I$. En particulier, puisque $|I'| < |I| \leq n$,
il résultera de $\sum a_{ij}m_j = 0$ et du
lemme \ref{lemma-recurring-symmetric-real} que les
sous-matrices $(a_{ij})_{i, j \in I'}$ sont définies négatives
et que leur déterminant est de signe $(-1)^m$.
Pour chaque possibilité restante, on calcule le déterminant de
$(a_{ij})_{i, j \in I}$. Si ce déterminant est de signe $-(-1)^{|I|}$,
on peut écarter ce cas, car le théorème de Sylvester montre
que la matrice $(a_{ij})_{i, j \in I}$ n'est pas semi-définie négative.
Si le déterminant est de signe $(-1)^{|I|}$, alors $|I| < n$, et l'on
conclut (provisoirement) que ce cas peut se présenter comme sous-graphe propre
et on l'inscrit dans l'un des lemmes de cette section.
Si le déterminant est $0$, on doit avoir $|I| = n$
(toujours d'après le lemme \ref{lemma-recurring-symmetric-real}) et $g = 0$.
Dans ces cas, on détermine toutes les valeurs possibles de $m_i, a_{ij}, w_i$, $i, j \in I$,
et on les énumère dans le lemme \ref{lemma-genus-one}. Une fois l'argument
achevé, on obtient tous les types numériques minimaux possibles de genre
$1$ avec $n > 1$, car chacun d'eux est nécessairement constitué entièrement
d'indices $(-2)$ (et apparaîtra donc au cours de la récurrence),
d'après la formule du genre et les remarques de la section précédente.
En définitive, le lecteur pourra parcourir la liste des
possibilités du lemme \ref{lemma-genus-one} pour vérifier que toutes les
configurations de sous-graphes propres présentées ici comme possibles
se réalisent effectivement déjà pour des types numériques de genre $1$.

\medskip\noindent
Supposons que $i$ et $j$ soient des indices $(-2)$ tels que $a_{ij} > 0$.
La matrice $A = (a_{ij})$ étant semi-définie négative d'après le
lemme \ref{lemma-recurring-symmetric-real}, on voit que la matrice
$$
\left(
\begin{matrix}
-2w_i & a_{ij} \\
a_{ij} & -2w_j
\end{matrix}
\right)
$$
est définie négative, sauf si $n = 2$. Le cas $n = 2$ peut se produire :
le déterminant $4w_1w_2 - a_{12}^2$ est alors nul. En utilisant que
$\text{lcm}(w_1, w_2)$ divise $a_{12}$, le lecteur trouve aisément
que les seules possibilités sont
$$
(w_1, w_2, a_{12}) = (w, w, 2w), (w, 4w, 4w), \text{ ou }(4w, w, 4w)
$$
Observons que le cas $(4w, w, 4w)$ s'obtient à partir du cas
$(w, 4w, 4w)$ en échangeant les indices $i, j$.
Dans ces cas, $g = 1$. Cela conduit aux cas
(\ref{item-two-cycle}) et (\ref{item-up4}) du lemme \ref{lemma-genus-one}.
En supposant $n > 2$, on voit
que le déterminant $4w_iw_j - a_{ij}^2$ de la matrice ci-dessus
est $> 0$, et l'on en déduit que $a_{ij}^2/w_iw_j < 4$.
D'autre part, on sait que $\text{lcm}(w_i, w_j) | a_{ij}$,
donc $a_{ij}^2/w_iw_j$ est un entier.
Ainsi $a_{ij}^2/w_iw_j \in \{1, 2, 3\}$ et $w_i | w_j$ ou
inversement. Cela conduit aux possibilités suivantes :
$$
(w_1, w_2, a_{12}) = (w, w, w), (w, 2w, 2w), (w, 3w, 3w),
(2w, w, 2w), \text{ ou }(3w, w, 3w)
$$
Observons que le cas $(2w, w, 2w)$ s'obtient à partir du cas
$(w, 2w, 2w)$ en échangeant les indices $i, j$, et de même
pour les cas $(3w, w, 3w)$ et $(w, 3w, 3w)$. Les trois premières
solutions conduisent aux cas (\ref{item-A2}), (\ref{item-B2}) et
(\ref{item-G2}) du lemme \ref{lemma-two-by-two}. Dans ce lemme,
on a explicité les conséquences pour les entiers $m_i$ et $m_j$,
en utilisant que $\sum_l a_{kl}m_l = 0$ pour tout $k$ entraîne en particulier
$a_{ii}m_i + a_{ij}m_j \leq 0$ pour $k = i$ et
$a_{ij}m_i + a_{jj}m_j \leq 0$ pour $k = j$.

\begin{lemma}
\label{lemma-two-by-two}
Classification des sous-graphes propres de la forme
$$
\xymatrix{
\bullet \ar@{-}[r] & \bullet
}
$$
Si $n > 2$, pour une paire $i, j$ d'indices $(-2)$ telle que $a_{ij} > 0$,
on a, à permutation près, les possibilités suivantes pour les $m$, les $a$ et les $w$ :
\begin{enumerate}
\item
\label{item-A2}
ils sont donnés par
$$
\left(
\begin{matrix}
m_1 \\
m_2
\end{matrix}
\right),
\quad
\left(
\begin{matrix}
-2w & w \\
w & -2w
\end{matrix}
\right),
\quad
\left(
\begin{matrix}
w \\
w
\end{matrix}
\right)
$$
avec $w$ arbitraire, $2m_1 \geq m_2$ et $2m_2 \geq m_1$, ou
\item
\label{item-B2}
ils sont donnés par
$$
\left(
\begin{matrix}
m_1 \\
m_2
\end{matrix}
\right),
\quad
\left(
\begin{matrix}
-2w & 2w \\
2w & -4w
\end{matrix}
\right),
\quad
\left(
\begin{matrix}
w \\
2w
\end{matrix}
\right)
$$
avec $w$ arbitraire, $m_1 \geq m_2$ et $2m_2 \geq m_1$, ou
\item
\label{item-G2}
ils sont donnés par
$$
\left(
\begin{matrix}
m_1 \\
m_2
\end{matrix}
\right),
\quad
\left(
\begin{matrix}
-2w & 3w \\
3w & -6w
\end{matrix}
\right),
\quad
\left(
\begin{matrix}
w \\
3w
\end{matrix}
\right)
$$
avec $w$ arbitraire, $2m_1 \geq 3m_2$ et $2m_2 \geq m_1$.
\end{enumerate}
\end{lemma}

\begin{proof}
Voir la discussion ci-dessus.
\end{proof}

\noindent
Supposons que $i$, $j$ et $k$ soient trois indices $(-2)$ tels que $a_{ij} > 0$
et $a_{jk} > 0$. Autrement dit, l'indice $i$ « rencontre » $j$, et
$j$ « rencontre » $k$. Nous utiliserons sans le rappeler que chaque paire
$(i, j)$, $(i, k)$ et $(j, k)$ figure dans la liste du lemme \ref{lemma-two-by-two}.
Puisque la matrice $A = (a_{ij})$ est semi-définie
négative d'après le lemme \ref{lemma-recurring-symmetric-real}, la matrice
$$
\left(
\begin{matrix}
-2w_i & a_{ij} & a_{ik} \\
a_{ij} & -2w_j & a_{jk} \\
a_{ik} & a_{jk} & -2w_k
\end{matrix}
\right)
$$
est définie négative, sauf si $n = 3$. Le cas $n = 3$ peut se produire :
le déterminant\footnote{Il vaut
$-8w_iw_jw_k + 2a_{ij}^2w_k + 2a_{jk}^2w_i + 2a_{ik}^2w_j +
2a_{ij}a_{jk}a_{ik}$.} de la matrice est alors nul,
et l'on obtient l'égalité
$$
4 = \frac{a_{ij}^2}{w_iw_j} +
\frac{a_{jk}^2}{w_jw_k} +
\frac{a_{ik}^2}{w_iw_k} +
\frac{a_{ij}a_{ik}a_{jk}}{w_iw_jw_k}
$$
entre entiers. Le dernier terme du membre de droite de cette égalité est déterminé
par les autres, car
$$
\left(\frac{a_{ij}a_{ik}a_{jk}}{w_iw_jw_k}\right)^2 =
\frac{a_{ij}^2}{w_iw_j} \frac{a_{jk}^2}{w_jw_k} \frac{a_{ik}^2}{w_iw_k}
$$
Puisque nous avons vu ci-dessus que
$\frac{a_{ij}^2}{w_iw_j}, \frac{a_{jk}^2}{w_jw_k}$ appartiennent à
$\{1, 2, 3\}$ et que $\frac{a_{ik}^2}{w_iw_k}$ appartient à $\{0, 1, 2, 3\}$,
on en déduit que les seules possibilités sont
$$
(\frac{a_{ij}^2}{w_iw_j}, \frac{a_{jk}^2}{w_jw_k}, \frac{a_{ik}^2}{w_iw_k}) =
(1, 1, 1), (1, 3, 0), (2, 2, 0),\text{ ou } (3, 1, 0)
$$
Observons que le cas $(3, 1, 0)$ s'obtient à partir du cas $(1, 3, 0)$
en renversant l'ordre des indices $i, j, k$.
Dans chacun de ces cas, $g = 1$ ; le lecteur les retrouvera aux cas
(\ref{item-three-cycle}), (\ref{item-equal-up3}), (\ref{item-equal-down3}),
(\ref{item-up-up}), (\ref{item-up-down}), (\ref{item-down-up}) du
lemme \ref{lemma-genus-one},
avec un cas correspondant à $(1, 1, 1)$,
deux cas correspondant à $(1, 3, 0)$ et
trois cas correspondant à $(2, 2, 0)$.
En supposant $n > 3$,
on obtient l'inégalité
$$
4 > \frac{a_{ij}^2}{w_iw_j} + \frac{a_{ik}^2}{w_iw_k} +
\frac{a_{jk}^2}{w_jw_k} + \frac{a_{ij}a_{ik}a_{jk}}{w_iw_jw_k}
$$
entre entiers. En utilisant les restrictions sur les nombres données ci-dessus,
on voit que les seules possibilités sont
$$
(\frac{a_{ij}^2}{w_iw_j}, \frac{a_{jk}^2}{w_jw_k}, \frac{a_{ik}^2}{w_iw_k}) =
(1, 1, 0), (1, 2, 0),\text{ ou }(2, 1, 0)
$$
en particulier $a_{ik} = 0$ (rappelons que l'on suppose
$a_{ij} > 0$ et $a_{jk} > 0$). Observons que le cas
$(2, 1, 0)$ s'obtient à partir du cas $(1, 2, 0)$ en renversant
l'ordre des indices $i, j, k$. Les deux premières solutions conduisent
aux cas (\ref{item-A3}), (\ref{item-C3}) et (\ref{item-B3})
du lemme \ref{lemma-three-by-three}, où l'on a aussi explicité les conséquences
pour les entiers $m_i$, $m_j$ et $m_k$.

\begin{lemma}
\label{lemma-three-by-three}
Classification des sous-graphes propres de la forme
$$
\xymatrix{
\bullet \ar@{-}[r] &
\bullet \ar@{-}[r] &
\bullet
}
$$
Si $n > 3$, pour un triplet $i, j, k$ d'indices $(-2)$
dont au moins deux des $a_{ij}, a_{ik}, a_{jk}$ sont non nuls, on a,
à permutation près, les possibilités suivantes pour les $m$, les $a$ et les $w$ :
\begin{enumerate}
\item
\label{item-A3}
ils sont donnés par
$$
\left(
\begin{matrix}
m_1 \\
m_2 \\
m_3
\end{matrix}
\right),
\quad
\left(
\begin{matrix}
-2w & w & 0 \\
w & -2w & w \\
0 & w & -2w
\end{matrix}
\right),
\quad
\left(
\begin{matrix}
w \\
w \\
w
\end{matrix}
\right)
$$
avec $2m_1 \geq m_2$, $2m_2 \geq m_1 + m_3$, $2m_3 \geq m_2$, ou
\item
\label{item-C3}
ils sont donnés par
$$
\left(
\begin{matrix}
m_1 \\
m_2 \\
m_3
\end{matrix}
\right),
\quad
\left(
\begin{matrix}
-2w & w & 0 \\
w & -2w & 2w \\
0 & 2w & -4w
\end{matrix}
\right),
\quad
\left(
\begin{matrix}
w \\
w \\
2w
\end{matrix}
\right)
$$
avec $2m_1 \geq m_2$, $2m_2 \geq m_1 + 2m_3$, $2m_3 \geq m_2$, ou
\item
\label{item-B3}
ils sont donnés par
$$
\left(
\begin{matrix}
m_1 \\
m_2 \\
m_3
\end{matrix}
\right),
\quad
\left(
\begin{matrix}
-4w & 2w & 0 \\
2w & -4w & 2w \\
0 & 2w & -2w
\end{matrix}
\right),
\quad
\left(
\begin{matrix}
2w \\
2w \\
w
\end{matrix}
\right)
$$
avec $2m_1 \geq m_2$, $2m_2 \geq m_1 + m_3$, $m_3 \geq m_2$.
\end{enumerate}
\end{lemma}

\begin{proof}
Voir la discussion ci-dessus.
\end{proof}

\noindent
Supposons que $i$, $j$, $k$ et $l$ soient quatre indices $(-2)$ tels que
$a_{ij} > 0$, $a_{jk} > 0$ et $a_{kl} > 0$. Autrement dit,
l'indice $i$ « rencontre » $j$, $j$ « rencontre » $k$, et $k$ « rencontre » $l$.
Le lemme \ref{lemma-three-by-three} donne alors $a_{ik} = a_{jl} = 0$.
La matrice $A = (a_{ij})$ étant semi-définie négative, on voit que la
matrice
$$
\left(
\begin{matrix}
-2w_i & a_{ij} & 0 & a_{il} \\
a_{ij} & -2w_j & a_{jk} & 0 \\
0 & a_{jk} & -2w_k & a_{kl} \\
a_{il} & 0 & a_{kl} & -2w_l
\end{matrix}
\right)
$$
est définie négative, sauf si $n = 4$. Le cas $n = 4$ peut se produire :
le déterminant\footnote{Il vaut
$16w_iw_jw_kw_l - 4a_{ij}^2w_kw_l - 4a_{jk}^2w_iw_l - 4a_{kl}^2w_iw_j -
4a_{il}^2w_jw_k + a_{ij}^2a_{kl}^2 + a_{jk}^2a_{il}^2 -
2a_{ij}a_{il}a_{jk}a_{kl}$.} de la matrice est alors nul, et l'on obtient l'égalité
$$
16 +
\frac{a_{ij}^2}{w_iw_j}\frac{a_{kl}^2}{w_kw_l} +
\frac{a_{jk}^2}{w_jw_k}\frac{a_{il}^2}{w_iw_l} =
4\frac{a_{ij}^2}{w_iw_j} + 4\frac{a_{jk}^2}{w_jw_k} + 4\frac{a_{kl}^2}{w_kw_l}
+ 4\frac{a_{il}^2}{w_iw_l} + 2\frac{a_{ij}a_{il}a_{jk}a_{kl}}{w_iw_jw_kw_l}
$$
entre entiers positifs ou nuls. Le dernier terme du membre de droite de cette égalité est
déterminé par les autres, car
$$
\left(\frac{a_{ij}a_{il}a_{jk}a_{kl}}{w_iw_jw_kw_l}\right)^2 =
\frac{a_{ij}^2}{w_iw_j} \frac{a_{jk}^2}{w_jw_k}
\frac{a_{kl}^2}{w_kw_l} \frac{a_{il}^2}{w_iw_l}
$$
Puisque nous avons vu ci-dessus que
$\frac{a_{ij}^2}{w_iw_j}, \frac{a_{jk}^2}{w_jw_k}, \frac{a_{kl}^2}{w_kw_l}$
appartiennent à $\{1, 2\}$ et que $\frac{a_{il}^2}{w_iw_l}$ appartient à $\{0, 1, 2\}$,
on en déduit que les seules solutions possibles sont
$$
(\frac{a_{ij}^2}{w_iw_j}, \frac{a_{jk}^2}{w_jw_k}, \frac{a_{kl}^2}{w_kw_l},
\frac{a_{il}^2}{w_iw_l}) =
(1, 1, 1, 1) \text{ ou } (2, 1, 2, 0)
$$
et que dans ce cas $g = 1$ ; le lecteur les retrouvera aux cas
(\ref{item-four-cycle}), (\ref{item-up-equal-up}),
(\ref{item-up-equal-down}) et (\ref{item-down-equal-up})
du lemme \ref{lemma-genus-one}. En supposant $n > 4$,
on obtient l'inégalité
$$
16 +
\frac{a_{ij}^2}{w_iw_j}\frac{a_{kl}^2}{w_kw_l} +
\frac{a_{jk}^2}{w_jw_k}\frac{a_{il}^2}{w_iw_l} >
4\frac{a_{ij}^2}{w_iw_j} + 4\frac{a_{jk}^2}{w_jw_k} + 4\frac{a_{kl}^2}{w_kw_l}
+ 4\frac{a_{il}^2}{w_iw_l} + 2\frac{a_{ij}a_{il}a_{jk}a_{kl}}{w_iw_jw_kw_l}
$$
entre entiers positifs ou nuls. En utilisant les restrictions sur les nombres données ci-dessus,
on voit que les seules possibilités sont
$$
(\frac{a_{ij}^2}{w_iw_j}, \frac{a_{jk}^2}{w_jw_k}, \frac{a_{kl}^2}{w_kw_l},
\frac{a_{il}^2}{w_iw_l}) =
(1, 1, 1, 0), (1, 1, 2, 0), (1, 2, 1, 0), \text{ ou }(2, 1, 1, 0)
$$
en particulier $a_{il} = 0$ (rappelons que les trois autres étaient supposés
non nuls). Observons que le cas $(2, 1, 1, 0)$ s'obtient
à partir du cas $(1, 1, 2, 0)$ en renversant l'ordre
des indices $i, j, k, l$. Les trois premières solutions conduisent
aux cas (\ref{item-A4}), (\ref{item-C4}), (\ref{item-B4}) et
(\ref{item-F4}) du lemme \ref{lemma-four-by-four},
où l'on a aussi explicité les conséquences pour les entiers $m_i$, $m_j$, $m_k$
et $m_l$.

\begin{lemma}
\label{lemma-four-by-four}
Classification des sous-graphes propres de la forme
$$
\xymatrix{
\bullet \ar@{-}[r] &
\bullet \ar@{-}[r] &
\bullet \ar@{-}[r] &
\bullet
}
$$
Si $n > 4$, pour quatre indices $(-2)$ $i, j, k, l$
tels que $a_{ij}, a_{jk}, a_{kl}$ soient non nuls, on a,
à permutation près, les possibilités suivantes pour les $m$, les $a$ et les $w$ :
\begin{enumerate}
\item
\label{item-A4}
ils sont donnés par
$$
\left(
\begin{matrix}
m_1 \\
m_2 \\
m_3 \\
m_4
\end{matrix}
\right),
\quad
\left(
\begin{matrix}
-2w & w & 0 & 0 \\
w & -2w & w & 0 \\
0 & w & -2w & w \\
0 & 0 & w & -2w 
\end{matrix}
\right),
\quad
\left(
\begin{matrix}
w \\
w \\
w \\
w
\end{matrix}
\right)
$$
avec $2m_1 \geq m_2$, $2m_2 \geq m_1 + m_3$, $2m_3 \geq m_2 + m_4$,
et $2m_4 \geq m_3$, ou
\item
\label{item-C4}
ils sont donnés par
$$
\left(
\begin{matrix}
m_1 \\
m_2 \\
m_3 \\
m_4
\end{matrix}
\right),
\quad
\left(
\begin{matrix}
-2w & w & 0 & 0 \\
w & -2w & w & 0 \\
0 & w & -2w & 2w \\
0 & 0 & 2w & -4w 
\end{matrix}
\right),
\quad
\left(
\begin{matrix}
w \\
w \\
w \\
2w
\end{matrix}
\right)
$$
avec $2m_1 \geq m_2$, $2m_2 \geq m_1 + m_3$, $2m_3 \geq m_2 + 2m_4$,
et $2m_4 \geq m_3$, ou
\item
\label{item-B4}
ils sont donnés par
$$
\left(
\begin{matrix}
m_1 \\
m_2 \\
m_3 \\
m_4
\end{matrix}
\right),
\quad
\left(
\begin{matrix}
-4w & 2w & 0 & 0 \\
2w & -4w & 2w & 0 \\
0 & 2w & -4w & 2w \\
0 & 0 & 2w & -2w 
\end{matrix}
\right),
\quad
\left(
\begin{matrix}
2w \\
2w \\
2w \\
w
\end{matrix}
\right)
$$
avec $2m_1 \geq m_2$, $2m_2 \geq m_1 + m_3$, $2m_3 \geq m_2 + m_4$,
et $m_4 \geq m_3$, ou
\item
\label{item-F4}
ils sont donnés par
$$
\left(
\begin{matrix}
m_1 \\
m_2 \\
m_3 \\
m_4
\end{matrix}
\right),
\quad
\left(
\begin{matrix}
-2w & w & 0 & 0 \\
w & -2w & 2w & 0 \\
0 & 2w & -4w & 2w \\
0 & 0 & 2w & -4w 
\end{matrix}
\right),
\quad
\left(
\begin{matrix}
w \\
w \\
2w \\
2w
\end{matrix}
\right)
$$
avec $2m_1 \geq m_2$, $2m_2 \geq m_1 + 2m_3$, $2m_3 \geq m_2 + m_4$,
et $2m_4 \geq m_3$.
\end{enumerate}
\end{lemma}

\begin{proof}
Voir la discussion ci-dessus.
\end{proof}

\noindent
Supposons que $i$, $j$, $k$
et $l$ soient quatre indices $(-2)$ tels que $a_{ij} > 0$, $a_{ij} > 0$ et
$a_{il} > 0$. Autrement dit, l'indice $i$ « rencontre » les indices
$j$, $k$, $l$. Le lemme \ref{lemma-three-by-three} donne alors
$a_{jk} = a_{jl} = a_{kl} = 0$. Puisque la matrice $A = (a_{ij})$ est
semi-définie négative, on voit que la matrice
$$
\left(
\begin{matrix}
-2w_i & a_{ij} & a_{ik} & a_{il} \\
a_{ij} & -2w_j & 0 & 0 \\
a_{ik} & 0 & -2w_k & 0 \\
a_{il} & 0 & 0 & -2w_l
\end{matrix}
\right)
$$
est définie négative, sauf si $n = 4$. Le cas $n = 4$ peut se produire :
le déterminant\footnote{Il vaut
$16w_iw_jw_kw_l - 4a_{ij}^2w_kw_l - 4a_{ik}^2w_jw_l - 4a_{il}^2w_jw_k$.}
de la matrice est alors nul, et l'on obtient l'égalité
$$
4 = \frac{a_{ij}^2}{w_iw_j} + \frac{a_{ik}^2}{w_iw_k} + \frac{a_{il}^2}{w_jw_l}
$$
entre entiers positifs ou nuls. Puisque nous avons vu ci-dessus que
$\frac{a_{ij}^2}{w_iw_j}, \frac{a_{ik}^2}{w_iw_k}, \frac{a_{il}^2}{w_iw_l}$
appartiennent à $\{1, 2\}$, on en déduit que la seule possibilité
est, à permutation près : $4 = 1 + 1 + 2$. Dans chacun de ces
cas, $g = 1$ ; le lecteur les retrouvera aux cas
(\ref{item-triple-with-up}) et (\ref{item-triple-with-down}) du
lemme \ref{lemma-genus-one}. En supposant $n > 4$,
on obtient l'inégalité
$$
4 > \frac{a_{ij}^2}{w_iw_j} + \frac{a_{ik}^2}{w_iw_k} + \frac{a_{il}^2}{w_jw_l}
$$
entre entiers positifs ou nuls. Il en résulte que $\frac{a_{ij}^2}{w_iw_j} =
\frac{a_{ik}^2}{w_iw_k} = \frac{a_{il}^2}{w_jw_l} = 1$
et que $w_i = w_j = w_k = w_l$. Cela conduit au cas
(\ref{item-D4}) du lemme \ref{lemma-D4},
où l'on a aussi explicité les conséquences pour les entiers $m_i$, $m_j$, $m_k$
et $m_l$.

\begin{lemma}
\label{lemma-D4}
Classification des sous-graphes propres de la forme
$$
\xymatrix{
\bullet \ar@{-}[r] & \bullet \ar@{-}[r] \ar@{-}[d] & \bullet \\
& \bullet
}
$$
Si $n > 4$, pour quatre indices $(-2)$ $i, j, k, l$
tels que $a_{ij}, a_{ik}, a_{il}$ soient non nuls, on a,
à permutation près, les possibilités suivantes pour les $m$, les $a$ et les $w$ :
\begin{enumerate}
\item
\label{item-D4}
ils sont donnés par
$$
\left(
\begin{matrix}
m_1 \\
m_2 \\
m_3 \\
m_4
\end{matrix}
\right),
\quad
\left(
\begin{matrix}
-2w & w & w & w \\
w & -2w & 0 & 0 \\
w & 0 & -2w & 0 \\
w & 0 & 0 & -2w 
\end{matrix}
\right),
\quad
\left(
\begin{matrix}
w \\
w \\
w \\
w
\end{matrix}
\right)
$$
avec $2m_1 \geq m_2 + m_3 + m_4$, $2m_2 \geq m_1$, $2m_3 \geq m_1$,
$2m_4 \geq m_1$. Observons que cela entraîne $m_1 \geq \max(m_2, m_3, m_4)$.
\end{enumerate}
\end{lemma}

\begin{proof}
Voir la discussion ci-dessus.
\end{proof}

\noindent
Supposons que $h$, $i$, $j$, $k$ et $l$ soient cinq indices $(-2)$
tels que $a_{hi} > 0$, $a_{ij} > 0$, $a_{jk} > 0$ et $a_{kl} > 0$.
Autrement dit, l'indice $h$ « rencontre » $i$, $i$ « rencontre » $j$,
$j$ « rencontre » $k$ et $k$ « rencontre » $l$.
On peut alors appliquer les lemmes \ref{lemma-three-by-three} et
\ref{lemma-four-by-four} pour voir que
$a_{hj} = a_{hk} = a_{ik} = a_{il} = a_{jl} = 0$
et que les fractions
$\frac{a_{hi}^2}{w_hw_i}, \frac{a_{ij}^2}{w_iw_j}, \frac{a_{jk}^2}{w_jw_k},
\frac{a_{kl}^2}{w_kw_l}$ appartiennent à $\{1, 2\}$, tandis que la fraction
$\frac{a_{hl}^2}{w_hw_l} \in \{0, 1, 2\}$.
Puisque la matrice $A = (a_{ij})$ est semi-définie négative, on voit que la
matrice
$$
\left(
\begin{matrix}
-2w_h & a_{hi} & 0 & 0 & a_{hl} \\
a_{hi} & -2w_i & a_{ij} & 0 & 0 \\
0 & a_{ij} & -2w_j & a_{jk} & 0 \\
0 & 0 & a_{jk} & -2w_k & a_{kl} \\
a_{hl} & 0 & 0 & a_{kl} & -2w_l
\end{matrix}
\right)
$$
est définie négative, sauf si $n = 5$. Le cas $n = 5$ peut se produire :
le déterminant\footnote{Il vaut
$-32w_hw_iw_jw_kw_l +
8a_{hi}^2w_jw_kw_l +
8a_{ij}^2w_hw_kw_l +
8a_{jk}^2w_hw_iw_l +
8a_{kl}^2w_hw_iw_j +
8a_{hl}^2w_iw_jw_k -
2a_{hi}^2a_{jk}^2w_l -
2a_{hi}^2a_{kl}^2w_j -
2a_{ij}^2a_{kl}^2w_h -
2a_{hl}^2a_{ij}^2w_k -
2a_{hl}^2a_{jk}^2w_i +
2a_{hi}a_{ij}a_{jk}a_{kl}a_{hl}
$
.}
de la matrice est nul et l'on obtient l'égalité
\begin{align*}
16 + 
\frac{a_{hi}^2}{w_hw_i}\frac{a_{jk}^2}{w_jw_k} +
\frac{a_{hi}^2}{w_hw_i}\frac{a_{kl}^2}{w_kw_l} +
\frac{a_{ij}^2}{w_iw_j}\frac{a_{kl}^2}{w_kw_l} +
\frac{a_{hl}^2}{w_hw_l}\frac{a_{ij}^2}{w_iw_j} +
\frac{a_{hl}^2}{w_hw_l}\frac{a_{jk}^2}{w_jw_k} \\
= 4\frac{a_{hi}^2}{w_hw_i} +
4\frac{a_{ij}^2}{w_iw_j} +
4\frac{a_{jk}^2}{w_jw_k} +
4\frac{a_{kl}^2}{w_kw_l} +
4\frac{a_{hl}^2}{w_hw_l} +
\frac{a_{hi}a_{ij}a_{jk}a_{kl}a_{hl}}{w_hw_iw_jw_kw_l}
\end{align*}
entre entiers positifs ou nuls. Le dernier terme du membre de droite de cette égalité
est déterminé par les autres, car
$$
\left(\frac{a_{hi}a_{ij}a_{jk}a_{kl}a_{hl}}{w_hw_iw_jw_kw_l} \right)^2 =
\frac{a_{hi}^2}{w_hw_i} \frac{a_{ij}^2}{w_iw_j}
\frac{a_{jk}^2}{w_jw_k} \frac{a_{kl}^2}{w_kw_l} \frac{a_{hl}^2}{w_hw_l}
$$
On conclut que les seules solutions possibles sont
$$
(\frac{a_{hi}^2}{w_hw_i}, \frac{a_{ij}^2}{w_iw_j},
\frac{a_{jk}^2}{w_jw_k}, \frac{a_{kl}^2}{w_kw_l}, \frac{a_{hl}^2}{w_hw_l}) =
(1, 1, 1, 1, 1), (1, 1, 2, 1, 0), (1, 2, 1, 1, 0), \text{ ou }(2, 1, 1, 2, 0)
$$
Observons que le cas $(1, 2, 1, 1, 0)$ s'obtient à partir du cas
$(1, 1, 2, 1, 0)$ en renversant l'ordre des indices $h, i, j, k, l$.
Dans ces cas, $g = 1$ ; le lecteur les retrouvera aux cas
(\ref{item-five-cycle}),
(\ref{item-equal-equal-up-equal}), (\ref{item-equal-equal-down-equal}),
(\ref{item-up-equal-equal-up}), (\ref{item-up-equal-equal-down}) et
(\ref{item-down-equal-equal-up}) du lemme \ref{lemma-genus-one},
avec un cas correspondant à $(1, 1, 1, 1, 1)$,
deux cas correspondant à $(1, 1, 2, 1, 0)$ et
trois cas correspondant à $(2, 1, 1, 2, 0)$.
En supposant $n > 5$, on obtient l'inégalité
\begin{align*}
16 + 
\frac{a_{hi}^2}{w_hw_i}\frac{a_{jk}^2}{w_jw_k} +
\frac{a_{hi}^2}{w_hw_i}\frac{a_{kl}^2}{w_kw_l} +
\frac{a_{ij}^2}{w_iw_j}\frac{a_{kl}^2}{w_kw_l} +
\frac{a_{hl}^2}{w_hw_l}\frac{a_{ij}^2}{w_iw_j} +
\frac{a_{hl}^2}{w_hw_l}\frac{a_{jk}^2}{w_jw_k} \\
> 4\frac{a_{hi}^2}{w_hw_i} +
4\frac{a_{ij}^2}{w_iw_j} +
4\frac{a_{jk}^2}{w_jw_k} +
4\frac{a_{kl}^2}{w_kw_l} +
4\frac{a_{hl}^2}{w_hw_l} +
\frac{a_{hi}a_{ij}a_{jk}a_{kl}a_{hl}}{w_hw_iw_jw_kw_l}
\end{align*}
entre entiers positifs ou nuls. En utilisant les restrictions données ci-dessus
sur ces nombres, on voit que les seules possibilités sont
$$
(\frac{a_{hi}^2}{w_hw_i}, \frac{a_{ij}^2}{w_iw_j},
\frac{a_{jk}^2}{w_jw_k}, \frac{a_{kl}^2}{w_kw_l}, \frac{a_{hl}^2}{w_hw_l}) =
(1, 1, 1, 1, 0), (1, 1, 1, 2, 0), \text{ ou } (2, 1, 1, 1, 0)
$$
en particulier, $a_{hl} = 0$ (rappelons que nous avons supposé les quatre autres
non nuls). Observons que le cas $(1, 1, 1, 2, 0)$
s'obtient à partir du cas $(2, 1, 1, 1, 0)$ en renversant l'ordre
des indices $h, i, j, k, l$. Les deux premières solutions conduisent aux
cas (\ref{item-A5}), (\ref{item-C5}) et (\ref{item-B5}) du
lemme \ref{lemma-five-by-five}, où l'on a aussi explicité les
conséquences pour les entiers $m_h$, $m_i$, $m_j$, $m_k$ et $m_l$.

\begin{lemma}
\label{lemma-five-by-five}
Classification des sous-graphes propres de la forme
$$
\xymatrix{
\bullet \ar@{-}[r] &
\bullet \ar@{-}[r] &
\bullet \ar@{-}[r] &
\bullet \ar@{-}[r] &
\bullet
}
$$
Si $n > 5$, pour cinq indices $(-2)$ $h, i, j, k, l$
tels que $a_{hi}, a_{ij}, a_{jk}, a_{kl}$ soient non nuls, on a,
à permutation près, les possibilités suivantes pour les $m$, les $a$ et les $w$ :
\begin{enumerate}
\item
\label{item-A5}
ils sont donnés par
$$
\left(
\begin{matrix}
m_1 \\
m_2 \\
m_3 \\
m_4 \\
m_5
\end{matrix}
\right),
\quad
\left(
\begin{matrix}
-2w & w & 0 & 0 & 0 \\
w & -2w & w & 0 & 0 \\
0 & w & -2w & w & 0 \\
0 & 0 & w & -2w & w \\
0 & 0 & 0 & w & -2w
\end{matrix}
\right),
\quad
\left(
\begin{matrix}
w \\
w \\
w \\
w \\
w
\end{matrix}
\right)
$$
avec $2m_1 \geq m_2$, $2m_2 \geq m_1 + m_3$, $2m_3 \geq m_2 + m_4$,
$2m_4 \geq m_3 + m_5$ et $2m_5 \geq m_4$, ou
\item
\label{item-C5}
ils sont donnés par
$$
\left(
\begin{matrix}
m_1 \\
m_2 \\
m_3 \\
m_4 \\
m_5
\end{matrix}
\right),
\quad
\left(
\begin{matrix}
-2w & w & 0 & 0 & 0 \\
w & -2w & w & 0 & 0 \\
0 & w & -2w & w & 0 \\
0 & 0 & w & -2w & 2w \\
0 & 0 & 0 & 2w & -4w
\end{matrix}
\right),
\quad
\left(
\begin{matrix}
w \\
w \\
w \\
w \\
2w
\end{matrix}
\right)
$$
avec $2m_1 \geq m_2$, $2m_2 \geq m_1 + m_3$, $2m_3 \geq m_2 + 2m_4$,
$2m_4 \geq m_3 + m_5$ et $2m_5 \geq m_4$, ou
\item
\label{item-B5}
ils sont donnés par
$$
\left(
\begin{matrix}
m_1 \\
m_2 \\
m_3 \\
m_4 \\
m_5
\end{matrix}
\right),
\quad
\left(
\begin{matrix}
-4w & 2w & 0 & 0 & 0 \\
2w & -4w & 2w & 0 & 0 \\
0 & 2w & -4w & 2w & 0 \\
0 & 0 & 2w & -4w & 2w \\
0 & 0 & 0 & 2w & -2w
\end{matrix}
\right),
\quad
\left(
\begin{matrix}
2w \\
2w \\
2w \\
2w \\
w
\end{matrix}
\right)
$$
avec $2m_1 \geq m_2$, $2m_2 \geq m_1 + m_3$, $2m_3 \geq m_2 + m_4$,
$2m_4 \geq m_3 + m_5$ et $m_4 \geq m_3$.
\end{enumerate}
\end{lemma}

\begin{proof}
Voir la discussion ci-dessus.
\end{proof}

\noindent
Supposons que $h$, $i$, $j$, $k$ et $l$ soient cinq indices $(-2)$
tels que $a_{hi} > 0$, $a_{hj} > 0$, $a_{hk} > 0$ et $a_{hl} > 0$.
Autrement dit, l'indice $h$ « rencontre » les indices
$i$, $j$, $k$, $l$. Le lemme \ref{lemma-three-by-three} montre alors
que $a_{ij} = a_{ik} = a_{il} = a_{jk} = a_{jl} = a_{kl} = 0$ et
le lemme \ref{lemma-D4} que
$w_h = w_i = w_j = w_k = w_l = w$ pour un entier $w > 0$ et
$a_{hi} = a_{hj} = a_{hk} = a_{hl} = -2w$.
La matrice correspondante
$$
\left(
\begin{matrix}
-2w & w & w & w & w \\
w & -2w & 0 & 0 & 0 \\
w & 0 & -2w & 0 & 0 \\
w & 0 & 0 & -2w & 0 \\
w & 0 & 0 & 0 & -2w
\end{matrix}
\right)
$$
est singulière. Cela ne peut donc se produire que si $n = 5$ et $g = 1$.
Le lecteur retrouvera cette possibilité au cas
(\ref{item-quadruple}) du lemme \ref{lemma-genus-one}.

\begin{lemma}
\label{lemma-fourfold}
Non-existence de sous-graphes propres de la forme
$$
\xymatrix{
\bullet \ar@{-}[r] & \bullet \ar@{-}[ld] \ar@{-}[r] \ar@{-}[d] & \bullet \\
\bullet & \bullet
}
$$
Si $n > 5$, il n'existe {\bf pas} cinq indices $(-2)$
$h$, $i$, $j$, $k$ tels que $a_{hi} > 0$, $a_{hj} > 0$, $a_{hk} > 0$ et
$a_{hl} > 0$.
\end{lemma}

\begin{proof}
Voir la discussion ci-dessus.
\end{proof}

\noindent
Supposons que $h$, $i$, $j$, $k$ et $l$ soient cinq indices $(-2)$
tels que $a_{hi} > 0$, $a_{ij} > 0$, $a_{jk} > 0$ et $a_{jl} > 0$.
Autrement dit, l'indice $h$ « rencontre » $i$, et l'indice $j$ « rencontre »
les indices $i$, $k$, $l$. Le
lemme \ref{lemma-D4} montre alors que $a_{ik} = a_{il} = a_{kl} = 0$,
$w_i = w_j = w_k = w_l = w$ et $a_{ij} = a_{jk} = a_{jl} = w$
pour un entier $w$. En appliquant le lemme \ref{lemma-four-by-four} aux
deux quadruplets
$h, i, j, k$ et $h, i, j, l$, on voit que $a_{hj} = a_{hk} = a_{hl} = 0$,
que $w_h = \frac{1}{2}w$, $w$ ou $2w$, et que
respectivement $a_{hi} = w$, $w$ ou $2w$.
Puisque $A$ est semi-définie négative, on voit que la matrice
$$
\left(
\begin{matrix}
-2w_h & a_{hi} & 0 & 0 & 0 \\
a_{hi} & -2w & w & 0 & 0 \\
0 & w & -2w & w & w \\
0 & 0 & w & -2w & 0 \\
0 & 0 & w & 0 & -2w
\end{matrix}
\right)
$$
est définie négative, sauf si $n = 5$. Le lecteur vérifiera que le
déterminant de la matrice vaut $0$ lorsque $w_h = \frac{1}{2}w$ ou $2w$.
Cela conduit aux cas (\ref{item-triple-extended-up}) et
(\ref{item-triple-extended-down}) du lemme \ref{lemma-genus-one}.
Pour $w_h = w$, on obtient le cas (\ref{item-D5}) du
lemme \ref{lemma-D5}.

\begin{lemma}
\label{lemma-D5}
Classification des sous-graphes propres de la forme
$$
\xymatrix{
\bullet \ar@{-}[r] & \bullet \ar@{-}[r] &
\bullet \ar@{-}[r] \ar@{-}[d] & \bullet \\
& & \bullet
}
$$
Si $n > 5$, pour cinq indices $(-2)$ $h, i, j, k, l$
tels que $a_{hi}, a_{ij}, a_{jk}, a_{jl}$ soient non nuls, on a,
à permutation près, les possibilités suivantes pour les $m$, les $a$ et les $w$ :
\begin{enumerate}
\item
\label{item-D5}
ils sont donnés par
$$
\left(
\begin{matrix}
m_1 \\
m_2 \\
m_3 \\
m_4 \\
m_5
\end{matrix}
\right),
\quad
\left(
\begin{matrix}
-2w & w & 0 & 0 & 0 \\
w & -2w & w & 0 & 0 \\
0 & w & -2w & w & w \\
0 & 0 & w & -2w & 0 \\
0 & 0 & w & 0 & -2w
\end{matrix}
\right),
\quad
\left(
\begin{matrix}
w \\
w \\
w \\
w \\
w
\end{matrix}
\right)
$$
avec $2m_1 \geq m_2$, $2m_2 \geq m_1 + m_3$, $2m_3 \geq m_2 + m_4 + m_5$,
$2m_4 \geq m_3$ et $2m_5 \geq m_3$.
\end{enumerate}
\end{lemma}

\begin{proof}
Voir la discussion ci-dessus.
\end{proof}

\noindent
Supposons que $t > 5$ et que $i_1, \ldots, i_t$ soient $t$ indices $(-2)$ distincts
tels que $a_{i_ji_{j + 1}}$ soit non nul pour $j = 1, \ldots, t - 1$. Nous allons
montrer par récurrence sur $t$ que, si $n = t$, cela conduit aux possibilités
(\ref{item-n-cycle}), (\ref{item-up-chain-equal-up}),
(\ref{item-up-chain-equal-down}), (\ref{item-down-chain-equal-up})
du lemme \ref{lemma-genus-one}, et que, si $n > t$, cela conduit aux cas
(\ref{item-An}), (\ref{item-Cn}) et (\ref{item-Bn}) du
lemme \ref{lemma-long}.
Premièrement, si $a_{i_1i_t}$ est non nul, il résulte clairement
du lemme \ref{lemma-five-by-five} que
$w_{i_1} = \ldots = w_{i_t} = w$, que
$a_{i_ji_{j + 1}} = w$ pour $j = 1, \ldots, t - 1$ et que
$a_{i_1i_t} = w$. Le vecteur $(1, \ldots, 1)$ appartient alors au noyau
de la matrice $t \times t$ correspondante. On doit donc avoir $n = t$ ;
on voit que le genre vaut $1$ et que l'on se trouve
dans le cas (\ref{item-n-cycle}) du lemme \ref{lemma-genus-one}.
On peut donc supposer $a_{i_1i_t} = 0$.
L'hypothèse de récurrence (ou le lemme \ref{lemma-five-by-five} si $t = 6$)
montre que $a_{i_ji_k} = 0$ si $k > j + 1$.
De plus, on a $w_{i_1} = \ldots = w_{i_{t - 1}} = w$
pour un entier $w$, et $w_{i_1}, w_{i_t} \in \{\frac{1}{2}w, w, 2w\}$.
En outre, si $w_{i_1}$, respectivement $w_{i_t}$, vaut
$\frac{1}{2}w$, $w$ ou $2w$, alors
$a_{i_1i_2}$, respectivement $a_{i_{t - 1}i_t}$,
vaut $w$, $w$ ou $2w$. Cela donne $9$ possibilités.
Dans chaque cas, il est facile de déterminer ce qui se passe :
\begin{enumerate}
\item si $(w_{i_1}, w_{i_t}) = (\frac{1}{2}w, \frac{1}{2}w)$, alors
on se trouve dans le cas (\ref{item-up-chain-equal-down}) du
lemme \ref{lemma-genus-one} ;
\item si $(w_{i_1}, w_{i_t}) = (\frac{1}{2}w, w)$ ou $(w, \frac{1}{2}w)$,
alors on se trouve dans le cas (\ref{item-Bn}) du lemme \ref{lemma-long} ;
\item si $(w_{i_1}, w_{i_t}) = (\frac{1}{2}w, 2w)$ ou $(2w, \frac{1}{2}w)$,
alors on se trouve dans le cas (\ref{item-up-chain-equal-up}) du
lemme \ref{lemma-genus-one} ;
\item si $(w_{i_1}, w_{i_t}) = (w, w)$, alors on se trouve dans le cas
(\ref{item-An}) du lemme \ref{lemma-long} ;
\item si $(w_{i_1}, w_{i_t}) = (w, 2w)$ ou $(2w, w)$, alors on se trouve
dans le cas (\ref{item-Cn}) du lemme \ref{lemma-long} ;
\item enfin, si $(w_{i_1}, w_{i_t}) = (2w, 2w)$, alors on se trouve dans le cas
(\ref{item-down-chain-equal-up}) du lemme \ref{lemma-genus-one}.
\end{enumerate}

\begin{lemma}
\label{lemma-long}
Classification des sous-graphes propres de la forme
$$
\xymatrix{
\bullet \ar@{-}[r] &
\bullet \ar@{-}[r] &
\bullet \ar@{..}[r] &
\bullet \ar@{-}[r] &
\bullet \ar@{-}[r] &
\bullet
}
$$
Soient $t > 5$ et $n > t$. Alors, pour $t$ indices $(-2)$ distincts
$i_1, \ldots, i_t$ tels que $a_{i_ji_{j + 1}}$ soit non nul pour
$j = 1, \ldots, t - 1$, on a, à renversement près de l'ordre de ces indices,
les possibilités suivantes pour les $a$ et les $w$ :
\begin{enumerate}
\item
\label{item-An}
ils sont donnés par $w_{i_1} = w_{i_2} = \ldots = w_{i_t} = w$,
$a_{i_ji_{j + 1}} = w$ et $a_{i_ji_k} = 0$ si $k > j + 1$, ou
\item
\label{item-Cn}
ils sont donnés par $w_{i_1} = w_{i_2} = \ldots = w_{i_{t - 1}} = w$,
$w_{j_t} = 2w$, $a_{i_ji_{j + 1}} = w$ pour $j < t - 1$,
$a_{i_{t - 1}i_t} = 2w$ et $a_{i_ji_k} = 0$ si $k > j + 1$, ou
\item
\label{item-Bn}
ils sont donnés par $w_{i_1} = w_{i_2} = \ldots = w_{i_{t - 1}} = 2w$,
$w_{j_t} = w$, $a_{i_ji_{j + 1}} = 2w$ et
$a_{i_{t - 1}i_t} = 2w$, et $a_{i_ji_k} = 0$ si $k > j + 1$.
\end{enumerate}
\end{lemma}

\begin{proof}
Voir la discussion ci-dessus.
\end{proof}

\noindent
Supposons que $t > 4$ et que $i_1, \ldots, i_{t + 1}$ soient $t + 1$
indices $(-2)$ distincts tels que $a_{i_ji_{j + 1}} > 0$
pour $j = 1, \ldots, t - 1$ et tels que $a_{j_{t - 1}j_{t + 1}} > 0$.
Voir la figure du lemme \ref{lemma-Dn}. Nous allons montrer par récurrence sur $t$
que, si $n = t + 1$, cela conduit aux possibilités
(\ref{item-Dn-extended-up}) et (\ref{item-Dn-extended-down})
du lemme \ref{lemma-genus-one}, et que, si $n > t + 1$, cela conduit au
cas (\ref{item-Dn}) du lemme \ref{lemma-Dn}.
L'hypothèse de récurrence (ou le lemme \ref{lemma-D5} lorsque $t = 5$)
montre que $a_{i_ji_k}$ est nul en dehors des termes requis
non nuls pour $j, k \geq 2$.
De plus, on voit que $w_2 = \ldots = w_{t + 1} = w$ pour un entier
$w$ et que les $a_{i_ji_k}$ non nuls pour $j, k \geq 2$ sont
égaux à $w$. En appliquant le lemme \ref{lemma-long}
(ou le lemme \ref{lemma-five-by-five} si $t = 5$) à la suite
$i_1, \ldots, i_t$ et à la suite
$i_1, \ldots, i_{t - 1}, i_{t + 1}$, on conclut que
$a_{i_1 i_j} = 0$ pour $j \geq 3$, que
$w_1$ vaut $\frac{1}{2}w$, $w$ ou $2w$
et que respectivement $a_{i_1i_2}$ vaut $w, w, 2w$.
Cela donne $3$ possibilités. Dans chaque cas, il est facile de
déterminer ce qui se passe :
\begin{enumerate}
\item Si $w_1 = \frac{1}{2}w$, alors on se trouve dans le cas
(\ref{item-Dn-extended-down}) du lemme \ref{lemma-genus-one}.
\item Si $w_1 = w$, alors on se trouve dans le cas
(\ref{item-Dn}) du lemme \ref{lemma-Dn}.
\item Si $w_1 = 2w$, alors on se trouve dans le cas
(\ref{item-Dn-extended-up}) du lemme \ref{lemma-genus-one}.
\end{enumerate}

\begin{lemma}
\label{lemma-Dn}
Classification des sous-graphes propres de la forme
$$
\xymatrix{
\bullet \ar@{-}[r] & \bullet \ar@{..}[r] & \bullet \ar@{-}[r] &
\bullet \ar@{-}[r] \ar@{-}[d] & \bullet \\
& & & \bullet
}
$$
Soient $t > 4$ et $n > t + 1$. Alors, pour $t + 1$ indices
$(-2)$ distincts $i_1, \ldots, i_{t + 1}$ tels que $a_{i_ji_{j + 1}}$
soit non nul pour $j = 1, \ldots, t - 1$ et que $a_{i_{t - 1}i_{t + 1}}$
soit non nul, on a les possibilités suivantes pour les $a$ et les $w$ :
\begin{enumerate}
\item
\label{item-Dn}
ils sont donnés par $w_{i_1} = w_{i_2} = \ldots = w_{i_{t + 1}} = w$,
$a_{i_ji_{j + 1}} = w$ pour $j = 1, \ldots, t - 1$,
$a_{i_{t - 1}i_{t + 1}} = w$ et $a_{i_ji_k} = 0$ pour les autres
paires $(j, k)$ telles que $j > k$.
\end{enumerate}
\end{lemma}

\begin{proof}
Voir la discussion ci-dessus.
\end{proof}

\noindent
Supposons donnés $6$ indices $(-2)$ distincts $g, h, i, j, k, l$
tels que $a_{gh}, a_{hi}, a_{ij}, a_{jk}, a_{il}$ soient non nuls.
Voir la figure du lemme \ref{lemma-E6}. On peut alors appliquer le
lemme \ref{lemma-D5} pour voir que l'on se trouve nécessairement dans la situation
du lemme \ref{lemma-E6}. Comme le déterminant vaut $3w^6 > 0$,
on conclut que, dans ce cas, on n'a jamais $n = 6$ !

\begin{lemma}
\label{lemma-E6}
Classification des sous-graphes propres de la forme
$$
\xymatrix{
\bullet \ar@{-}[r] & \bullet \ar@{-}[r] & \bullet \ar@{-}[r] \ar@{-}[d] &
\bullet \ar@{-}[r] & \bullet \\
& & \bullet
}
$$
Soit $n > 6$. Alors, pour $6$ indices $(-2)$ distincts
$i_1, \ldots, i_6$ tels que
$a_{12}, a_{23}, a_{34},
a_{45}, a_{36}$ soient non nuls, on a
les possibilités suivantes pour les $m$, les $a$ et les $w$ :
\begin{enumerate}
\item
\label{item-E6}
ils sont donnés par
$$
\left(
\begin{matrix}
m_1 \\
m_2 \\
m_3 \\
m_4 \\
m_5 \\
m_6
\end{matrix}
\right),
\quad
\left(
\begin{matrix}
-2w & w & 0 & 0 & 0 & 0 \\
w & -2w & w & 0 & 0 & 0 \\
0 & w & -2w & w & 0 & w \\
0 & 0 & w & -2w & w & 0 \\
0 & 0 & 0 & w & -2w & 0 \\
0 & 0 & w & 0 & 0 & -2w
\end{matrix}
\right),
\quad
\left(
\begin{matrix}
w \\
w \\
w \\
w \\
w \\
w
\end{matrix}
\right)
$$
avec $2m_1 \geq m_2$, $2m_2 \geq m_1 + m_3$, $2m_3 \geq m_2 + m_4 + m_6$,
$2m_4 \geq m_3 + m_5$, $2m_5 \geq m_3$ et $2m_6 \geq m_3$.
\end{enumerate}
\end{lemma}

\begin{proof}
Voir la discussion ci-dessus.
\end{proof}

\noindent
Supposons que $t \geq 4$ et que $i_0, \ldots, i_{t + 1}$ soient
$t + 2$ indices $(-2)$ distincts tels que
$a_{i_ji_{j + 1}} > 0$ pour $j = 1, \ldots, t - 1$,
$a_{i_0i_2} > 0$ et $a_{i_{t - 1}i_{t + 1}} > 0$.
Voir la figure du lemme \ref{lemma-double-triple}.
On peut alors appliquer les lemmes \ref{lemma-D5} et \ref{lemma-Dn}
pour voir que tous les autres $a_{i_ji_k}$ pour $j < k$ sont nuls,
que $w_{i_0} = \ldots = w_{i_{t + 1}} = w$ pour un certain
entier $w$ et que les coefficients hors diagonale requis
non nuls de $A$ sont égaux à $w$.
Un calcul montre que le déterminant de la
matrice correspondante est nul. Ainsi, $n = t + 2$
et l'on se trouve dans le cas (\ref{item-double-triple}) du
lemme \ref{lemma-genus-one}.

\begin{lemma}
\label{lemma-double-triple}
Non-existence de sous-graphes propres de la forme
$$
\xymatrix{
\bullet \ar@{-}[r] & \bullet \ar@{..}[r] \ar@{-}[d] &
\bullet \ar@{-}[d] \ar@{-}[r] & \bullet \\
& \bullet & \bullet
}
$$
Supposons $t \geq 4$ et $n > t + 2$.
Il n'existe {\bf pas} $t + 2$ indices
$(-2)$ distincts $i_0, \ldots, i_{t + 1}$ tels que
$a_{i_ji_{j + 1}} > 0$ pour $j = 1, \ldots, t - 1$,
$a_{i_0i_2} > 0$ et $a_{i_{t - 1}i_{t + 1}} > 0$.
\end{lemma}

\begin{proof}
Voir la discussion ci-dessus.
\end{proof}

\noindent
Supposons donnés $7$ indices $(-2)$ distincts
$f, g, h, i,
j, k, l$
tels que les nombres
$a_{fg}, a_{gh}, a_{ij},
a_{jh}, a_{kl}, a_{lh}$ soient non nuls.
Voir la figure du lemme \ref{lemma-E6-completed}. On peut alors appliquer le
lemme \ref{lemma-D5} pour voir que la matrice correspondante est
$$
\left(
\begin{matrix}
-2w & w & 0 & 0 & 0 & 0 & 0 \\
w & -2w & w & 0 & 0 & 0 & 0 \\
0 & w & -2w & 0 & w & 0 & w \\
0 & 0 & 0 & -2w & w & 0 & 0 \\
0 & 0 & w & w & -2w & 0 & 0 \\
0 & 0 & 0 & 0 & 0 & -2w & w \\
0 & 0 & w & 0 & 0 & w & -2w
\end{matrix}
\right)
$$
Comme le déterminant vaut $0$, on conclut que nécessairement $n = 7$
et $g = 1$, et l'on obtient le cas (\ref{item-E6-completed})
du lemme \ref{lemma-genus-one}.

\begin{lemma}
\label{lemma-E6-completed}
Non-existence de sous-graphes propres de la forme
$$
\xymatrix{
\bullet \ar@{-}[r] & \bullet \ar@{-}[r] & \bullet \ar@{-}[r] \ar@{-}[d] &
\bullet \ar@{-}[r] & \bullet \\
& & \bullet \ar@{-}[d] \\
& & \bullet
}
$$
Supposons $n > 7$. Il n'existe {\bf pas} $7$ indices
$(-2)$ distincts $f, g, h, i, j, k, l$
tels que $a_{fg}, a_{gh}, a_{ij}, a_{jh}, a_{kl}, a_{lh}$ soient non nuls.
\end{lemma}

\begin{proof}
Voir la discussion ci-dessus.
\end{proof}

\noindent
Supposons donnés $7$ indices $(-2)$ distincts
$f, g, h, i,
j, k, l$
tels que les nombres
$a_{fg}, a_{gh}, a_{hi},
a_{ij}, a_{jk}, a_{il}$ soient non nuls.
Voir la figure du lemme \ref{lemma-E7}. On peut alors appliquer les
lemmes \ref{lemma-D5} et \ref{lemma-Dn}
pour voir que l'on se trouve nécessairement dans la situation
du lemme \ref{lemma-E7}. Comme le déterminant vaut $-8w^7 > 0$,
on conclut que, dans ce cas, on n'a jamais $n = 7$ !

\begin{lemma}
\label{lemma-E7}
Classification des sous-graphes propres de la forme
$$
\xymatrix{
\bullet \ar@{-}[r] & \bullet \ar@{-}[r] & \bullet \ar@{-}[r] &
\bullet \ar@{-}[r] \ar@{-}[d] & \bullet \ar@{-}[r] & \bullet \\
& & & \bullet
}
$$
Soit $n > 7$. Alors, pour $7$ indices $(-2)$ distincts
$i_1, \ldots, i_7$ tels que
$a_{12}, a_{23}, a_{34},
a_{45}, a_{56}, a_{47}$ soient non nuls,
on a les possibilités suivantes pour les $m$, les $a$ et les $w$ :
\begin{enumerate}
\item
\label{item-E7}
ils sont donnés par
$$
\left(
\begin{matrix}
m_1 \\
m_2 \\
m_3 \\
m_4 \\
m_5 \\
m_6 \\
m_7
\end{matrix}
\right),
\quad
\left(
\begin{matrix}
-2w & w & 0 & 0 & 0 & 0 & 0 \\
w & -2w & w & 0 & 0 & 0 & 0 \\
0 & w & -2w & w & 0 & 0 & 0 \\
0 & 0 & w & -2w & w & 0 & w \\
0 & 0 & 0 & w & -2w & w & 0 \\
0 & 0 & 0 & 0 & w & -2w & 0 \\
0 & 0 & 0 & w & 0 & 0 & -2w
\end{matrix}
\right),
\quad
\left(
\begin{matrix}
w \\
w \\
w \\
w \\
w \\
w \\
w
\end{matrix}
\right)
$$
avec $2m_1 \geq m_2$, $2m_2 \geq m_1 + m_3$, $2m_3 \geq m_2 + m_4$,
$2m_4 \geq m_3 + m_5 + m_7$, $2m_5 \geq m_4 + m_6$, $2m_6 \geq m_5$
et $2m_7 \geq m_4$.
\end{enumerate}
\end{lemma}

\begin{proof}
Voir la discussion ci-dessus.
\end{proof}

\noindent
Supposons donnés $8$ indices $(-2)$ distincts dont la configuration
des coefficients non nuls $a_{ij}$ de la matrice $A$ est de la forme
$$
\xymatrix{
\bullet \ar@{-}[r] & \bullet \ar@{-}[r] & \bullet \ar@{-}[r] &
\bullet \ar@{-}[r] & \bullet \ar@{-}[r] \ar@{-}[d] &
\bullet \ar@{-}[r] & \bullet \\
& & & & \bullet
}
$$
ou de la forme
$$
\xymatrix{
\bullet \ar@{-}[r] & \bullet \ar@{-}[r] & \bullet \ar@{-}[r] &
\bullet \ar@{-}[r] \ar@{-}[d] & \bullet \ar@{-}[r] &
\bullet \ar@{-}[r] & \bullet \\
& & & \bullet
}
$$
En raisonnant exactement comme dans la démonstration du lemme \ref{lemma-E7},
on voit que la première configuration conduit au cas
(\ref{item-E8}) du lemme \ref{lemma-E8}
et ne conduit pas à un nouveau cas dans le lemme \ref{lemma-genus-one}.
En raisonnant exactement comme dans la démonstration du lemme \ref{lemma-E6-completed},
on voit que la seconde configuration ne se présente pas si
$n > 8$, mais qu'elle conduit au cas (\ref{item-E7-completed})
du lemme \ref{lemma-genus-one} lorsque $n = 8$.

\begin{lemma}
\label{lemma-E8}
Classification des sous-graphes propres de la forme
$$
\xymatrix{
\bullet \ar@{-}[r] & \bullet \ar@{-}[r] & \bullet \ar@{-}[r] &
\bullet \ar@{-}[r] & \bullet \ar@{-}[r] \ar@{-}[d] &
\bullet \ar@{-}[r] & \bullet \\
& & & & \bullet
}
$$
Soit $n > 8$. Alors, pour $8$ indices $(-2)$ distincts
$i_1, \ldots, i_8$ tels que
$a_{12}, a_{23}, a_{34},
a_{45}, a_{56}, a_{65}, a_{57}$
soient non nuls, on a les possibilités suivantes pour les $m$, les $a$ et les $w$ :
\begin{enumerate}
\item
\label{item-E8}
ils sont donnés par
$$
\left(
\begin{matrix}
m_1 \\
m_2 \\
m_3 \\
m_4 \\
m_5 \\
m_6 \\
m_7 \\
m_8
\end{matrix}
\right),
\quad
\left(
\begin{matrix}
-2w & w & 0 & 0 & 0 & 0 & 0 & 0 \\
w & -2w & w & 0 & 0 & 0 & 0 & 0 \\
0 & w & -2w & w & 0 & 0 & 0 & 0 \\
0 & 0 & w & -2w & w & 0 & 0 & 0 \\
0 & 0 & 0 & w & -2w & w & 0 & w \\
0 & 0 & 0 & 0 & w & -2w & w & 0 \\
0 & 0 & 0 & 0 & 0 & w & -2w & 0 \\
0 & 0 & 0 & 0 & w & 0 & 0 & -2w
\end{matrix}
\right),
\quad
\left(
\begin{matrix}
w \\
w \\
w \\
w \\
w \\
w \\
w \\
w
\end{matrix}
\right)
$$
avec $2m_1 \geq m_2$, $2m_2 \geq m_1 + m_3$, $2m_3 \geq m_2 + m_4$,
$2m_4 \geq m_3 + m_5$, $2m_5 \geq m_4 + m_6 + m_8$, $2m_6 \geq m_5 + m_7$,
$2m_7 \geq m_6$ et $2m_8 \geq m_5$.
\end{enumerate}
\end{lemma}

\begin{proof}
Voir la discussion ci-dessus.
\end{proof}

\begin{lemma}
\label{lemma-E7-completed}
Non-existence de sous-graphes propres de la forme
$$
\xymatrix{
\bullet \ar@{-}[r] &
\bullet \ar@{-}[r] &
\bullet \ar@{-}[r] & \bullet \ar@{-}[r] \ar@{-}[d] &
\bullet \ar@{-}[r] & \bullet \ar@{-}[r] & \bullet \\
& & & \bullet
}
$$
Supposons $n > 8$. Il n'existe {\bf pas} $8$ indices
$(-2)$ distincts $e, f, g, h, i, j, k, l$
tels que $a_{ef}, a_{fg}, a_{gh}, a_{hi}, a_{ij}, a_{jk}, a_{lh}$
soient non nuls.
\end{lemma}

\begin{proof}
Voir la discussion ci-dessus.
\end{proof}

\noindent
Supposons donnés $9$ indices $(-2)$ distincts dont la configuration
des coefficients non nuls $a_{ij}$ de la matrice $A$ est de la forme
$$
\xymatrix{
\bullet \ar@{-}[r] & \bullet \ar@{-}[r] &
\bullet \ar@{-}[r] & \bullet \ar@{-}[r] &
\bullet \ar@{-}[r] & \bullet \ar@{-}[r] \ar@{-}[d] &
\bullet \ar@{-}[r] & \bullet \\
& & & & & \bullet
}
$$
En raisonnant exactement comme dans la démonstration du lemme \ref{lemma-E6-completed},
on voit que cette configuration ne se présente pas si
$n > 9$, mais qu'elle conduit au cas (\ref{item-E8-completed})
du lemme \ref{lemma-genus-one} lorsque $n = 9$.

\begin{lemma}
\label{lemma-E8-completed}
Non-existence de sous-graphes propres de la forme
$$
\xymatrix{
\bullet \ar@{-}[r] & \bullet \ar@{-}[r] &
\bullet \ar@{-}[r] & \bullet \ar@{-}[r] &
\bullet \ar@{-}[r] & \bullet \ar@{-}[r] \ar@{-}[d] &
\bullet \ar@{-}[r] & \bullet \\
& & & & & \bullet
}
$$
Supposons $n > 9$. Il n'existe {\bf pas} $9$ indices
$(-2)$ distincts $d, e, f, g, h, i, j, k, l$
tels que $a_{de}, a_{ef}, a_{fg}, a_{gh}, a_{hi}, a_{ij}, a_{jk}, a_{lh}$
soient non nuls.
\end{lemma}

\begin{proof}
Voir la discussion ci-dessus.
\end{proof}

\noindent
En réunissant toutes les informations, on obtient le résultat suivant.

\begin{proposition}
\label{proposition-classify-subgraphs}
Soient $n, m_i, a_{ij}, w_i, g_i$ un type numérique de genre $g$.
Soit $I \subset \{1, \ldots, n\}$ un sous-ensemble strict de cardinal $\geq 2$
constitué d'indices $(-2)$ et tel qu'il
n'existe aucun sous-ensemble strict non vide $I' \subset I$
avec $a_{i'i} = 0$ pour $i' \in I$, $i \in I \setminus I'$.
Alors, à réindexation près, les $m_i$, les $a_{ij}$ et les $w_i$
pour $i, j \in I$ sont ceux qui figurent dans les
lemmes \ref{lemma-two-by-two},
\ref{lemma-three-by-three},
\ref{lemma-four-by-four},
\ref{lemma-D4},
\ref{lemma-five-by-five},
\ref{lemma-D5},
\ref{lemma-long},
\ref{lemma-Dn},
\ref{lemma-E6},
\ref{lemma-E7} ou
\ref{lemma-E8}.
\end{proposition}

\begin{proof}
Cela résulte de la discussion ci-dessus ; voir la discussion au
début de la section \ref{section-classify-proper-subgraphs}.
\end{proof}













\section{Classification des types minimaux de genres zéro et un}
\label{section-classification-genus-one}

\noindent
Le titre de la section dit tout.

\begin{lemma}[Genre zéro]
\label{lemma-genus-zero}
Le seul type numérique minimal de genre zéro est
$n = 1$, $m_1 = 1$, $a_{11} = 0$, $w_1 = 1$, $g_1 = 0$.
\end{lemma}

\begin{proof}
Cela résulte des lemmes \ref{lemma-non-irreducible-minimal-type-genus-at-least-one}
et \ref{lemma-irreducible}.
\end{proof}

\begin{lemma}[Genre un]
\label{lemma-genus-one}
À équivalence près, les types numériques minimaux de genre un sont les suivants :
\begin{enumerate}
\item
\label{item-one}
$n = 1$, $a_{11} = 0$, $g_1 = 1$, et $m_1, w_1 \geq 1$ sont arbitraires ;
\item
\label{item-two-cycle}
$n = 2$, et $m_i, a_{ij}, w_i, g_i$ sont donnés par
$$
\left(
\begin{matrix}
m \\
m
\end{matrix}
\right),
\quad
\left(
\begin{matrix}
-2w & 2w \\
2w & -2w
\end{matrix}
\right),
\quad
\left(
\begin{matrix}
w \\
w
\end{matrix}
\right),
\quad
\left(
\begin{matrix}
0 \\
0
\end{matrix}
\right)
$$
où $w$ et $m$ sont arbitraires ;
\item
\label{item-up4}
$n = 2$, et $m_i, a_{ij}, w_i, g_i$ sont donnés par
$$
\left(
\begin{matrix}
2m \\
m
\end{matrix}
\right),
\quad
\left(
\begin{matrix}
-2w & 4w \\
4w & -8w
\end{matrix}
\right),
\quad
\left(
\begin{matrix}
w \\
4w
\end{matrix}
\right),
\quad
\left(
\begin{matrix}
0 \\
0
\end{matrix}
\right)
$$
où $w$ et $m$ sont arbitraires ;
\item
\label{item-three-cycle}
$n = 3$, et $m_i, a_{ij}, w_i, g_i$ sont donnés par
$$
\left(
\begin{matrix}
m \\
m \\
m
\end{matrix}
\right),
\quad
\left(
\begin{matrix}
-2w & w & w \\
w & -2w & w \\
w & w & -2w
\end{matrix}
\right),
\quad
\left(
\begin{matrix}
w \\
w \\
w
\end{matrix}
\right),
\quad
\left(
\begin{matrix}
0 \\
0 \\
0
\end{matrix}
\right)
$$
où $w$ et $m$ sont arbitraires ;
\item
\label{item-equal-up3}
$n = 3$, et $m_i, a_{ij}, w_i, g_i$ sont donnés par
$$
\left(
\begin{matrix}
m \\
2m \\
m
\end{matrix}
\right),
\quad
\left(
\begin{matrix}
-2w & w & 0 \\
w & -2w & 3w \\
0 & 3w & -6w
\end{matrix}
\right),
\quad
\left(
\begin{matrix}
w \\
w \\
3w
\end{matrix}
\right),
\quad
\left(
\begin{matrix}
0 \\
0 \\
0
\end{matrix}
\right)
$$
où $w$ et $m$ sont arbitraires ;
\item
\label{item-equal-down3}
$n = 3$, et $m_i, a_{ij}, w_i, g_i$ sont donnés par
$$
\left(
\begin{matrix}
m \\
2m \\
3m
\end{matrix}
\right),
\quad
\left(
\begin{matrix}
-6w & 3w & 0 \\
3w & -6w & 3w \\
0 & 3w & -2w
\end{matrix}
\right),
\quad
\left(
\begin{matrix}
3w \\
3w \\
w
\end{matrix}
\right),
\quad
\left(
\begin{matrix}
0 \\
0 \\
0
\end{matrix}
\right)
$$
où $w$ et $m$ sont arbitraires ;
\item
\label{item-up-up}
$n = 3$, et $m_i, a_{ij}, w_i, g_i$ sont donnés par
$$
\left(
\begin{matrix}
2m \\
2m \\
m
\end{matrix}
\right),
\quad
\left(
\begin{matrix}
-2w & 2w & 0 \\
2w & -4w & 4w \\
0 & 4w & -8w
\end{matrix}
\right),
\quad
\left(
\begin{matrix}
w \\
2w \\
4w
\end{matrix}
\right),
\quad
\left(
\begin{matrix}
0 \\
0 \\
0
\end{matrix}
\right)
$$
où $w$ et $m$ sont arbitraires ;
\item
\label{item-up-down}
$n = 3$, et $m_i, a_{ij}, w_i, g_i$ sont donnés par
$$
\left(
\begin{matrix}
m \\
m \\
m
\end{matrix}
\right),
\quad
\left(
\begin{matrix}
-2w & 2w & 0 \\
2w & -4w & 2w \\
0 & 2w & -2w
\end{matrix}
\right),
\quad
\left(
\begin{matrix}
w \\
2w \\
w
\end{matrix}
\right),
\quad
\left(
\begin{matrix}
0 \\
0 \\
0
\end{matrix}
\right)
$$
où $w$ et $m$ sont arbitraires ;
\item
\label{item-down-up}
$n = 3$, et $m_i, a_{ij}, w_i, g_i$ sont donnés par
$$
\left(
\begin{matrix}
m \\
2m \\
m
\end{matrix}
\right),
\quad
\left(
\begin{matrix}
-4w & 2w & 0 \\
2w & -2w & 2w \\
0 & 2w & -4w
\end{matrix}
\right),
\quad
\left(
\begin{matrix}
2w \\
w \\
2w
\end{matrix}
\right),
\quad
\left(
\begin{matrix}
0 \\
0 \\
0
\end{matrix}
\right)
$$
où $w$ et $m$ sont arbitraires ;
\item
\label{item-four-cycle}
$n = 4$, et $m_i, a_{ij}, w_i, g_i$ sont donnés par
$$
\left(
\begin{matrix}
m \\
m \\
m \\
m
\end{matrix}
\right),
\quad
\left(
\begin{matrix}
-2w & w & 0 & w \\
w & -2w & w & 0 \\
0 & w & -2w & w \\
w & 0 & w & -2w
\end{matrix}
\right),
\quad
\left(
\begin{matrix}
w \\
w \\
w \\
w
\end{matrix}
\right),
\quad
\left(
\begin{matrix}
0 \\
0 \\
0 \\
0
\end{matrix}
\right)
$$
où $w$ et $m$ sont arbitraires ;
\item
\label{item-up-equal-up}
$n = 4$, et $m_i, a_{ij}, w_i, g_i$ sont donnés par
$$
\left(
\begin{matrix}
2m \\
2m \\
2m \\
m
\end{matrix}
\right),
\quad
\left(
\begin{matrix}
-2w & 2w & 0 & 0 \\
2w & -4w & 2w & 0 \\
0 & 2w & -4w & 4w \\
0 & 0 & 4w & -8w
\end{matrix}
\right),
\quad
\left(
\begin{matrix}
w \\
2w \\
2w \\
4w
\end{matrix}
\right),
\quad
\left(
\begin{matrix}
0 \\
0 \\
0 \\
0
\end{matrix}
\right)
$$
où $w$ et $m$ sont arbitraires ;
\item
\label{item-up-equal-down}
$n = 4$, et $m_i, a_{ij}, w_i, g_i$ sont donnés par
$$
\left(
\begin{matrix}
m \\
m \\
m \\
m
\end{matrix}
\right),
\quad
\left(
\begin{matrix}
-2w & 2w & 0 & 0 \\
2w & -4w & 2w & 0 \\
0 & 2w & -4w & 2w \\
0 & 0 & 2w & -2w
\end{matrix}
\right),
\quad
\left(
\begin{matrix}
w \\
2w \\
2w \\
w
\end{matrix}
\right),
\quad
\left(
\begin{matrix}
0 \\
0 \\
0 \\
0
\end{matrix}
\right)
$$
où $w$ et $m$ sont arbitraires ;
\item
\label{item-down-equal-up}
$n = 4$, et $m_i, a_{ij}, w_i, g_i$ sont donnés par
$$
\left(
\begin{matrix}
m \\
2m \\
2m \\
m
\end{matrix}
\right),
\quad
\left(
\begin{matrix}
-4w & 2w & 0 & 0 \\
2w & -2w & w & 0 \\
0 & w & -2w & 2w \\
0 & 0 & 2w & -4w
\end{matrix}
\right),
\quad
\left(
\begin{matrix}
2w \\
w \\
w \\
2w
\end{matrix}
\right),
\quad
\left(
\begin{matrix}
0 \\
0 \\
0 \\
0
\end{matrix}
\right)
$$
où $w$ et $m$ sont arbitraires ;
\item
\label{item-triple-with-up}
$n = 4$, et $m_i, a_{ij}, w_i, g_i$ sont donnés par
$$
\left(
\begin{matrix}
2m \\
m \\
m \\
m
\end{matrix}
\right),
\quad
\left(
\begin{matrix}
-2w & w & w & 2w \\
w & -2w & 0 & 0 \\
w & 0 & -2w & 0 \\
2w & 0 & 0 & -4w
\end{matrix}
\right),
\quad
\left(
\begin{matrix}
w \\
w \\
w \\
2w
\end{matrix}
\right),
\quad
\left(
\begin{matrix}
0 \\
0 \\
0 \\
0
\end{matrix}
\right)
$$
où $w$ et $m$ sont arbitraires ;
\item
\label{item-triple-with-down}
$n = 4$, et $m_i, a_{ij}, w_i, g_i$ sont donnés par
$$
\left(
\begin{matrix}
2m \\
m \\
m \\
2m
\end{matrix}
\right),
\quad
\left(
\begin{matrix}
-4w & 2w & 2w & 2w \\
2w & -4w & 0 & 0 \\
2w & 0 & -4w & 0 \\
2w & 0 & 0 & -2w
\end{matrix}
\right),
\quad
\left(
\begin{matrix}
2w \\
2w \\
2w \\
w
\end{matrix}
\right),
\quad
\left(
\begin{matrix}
0 \\
0 \\
0 \\
0
\end{matrix}
\right)
$$
où $w$ et $m$ sont arbitraires ;
\item
\label{item-five-cycle}
$n = 5$, et $m_i, a_{ij}, w_i, g_i$ sont donnés par
$$
\left(
\begin{matrix}
m \\
m \\
m \\
m \\
m
\end{matrix}
\right),
\quad
\left(
\begin{matrix}
-2w & w & 0 & 0 & w \\
w & -2w & w & 0 & 0 \\
0 & w & -2w & w & 0 \\
0 & 0 & w & -2w & w \\
w & 0 & 0 & w & -2w \\
\end{matrix}
\right),
\quad
\left(
\begin{matrix}
w \\
w \\
w \\
w \\
w
\end{matrix}
\right),
\quad
\left(
\begin{matrix}
0 \\
0 \\
0 \\
0 \\
0
\end{matrix}
\right)
$$
où $w$ et $m$ sont arbitraires ;
\item
\label{item-equal-equal-up-equal}
$n = 5$, et $m_i, a_{ij}, w_i, g_i$ sont donnés par
$$
\left(
\begin{matrix}
m \\
2m \\
3m \\
2m \\
m
\end{matrix}
\right),
\quad
\left(
\begin{matrix}
-2w & w & 0 & 0 & 0 \\
w & -2w & w & 0 & 0 \\
0 & w & -2w & 2w & 0 \\
0 & 0 & 2w & -4w & 2w \\
0 & 0 & 0 & 2w & -4w \\
\end{matrix}
\right),
\quad
\left(
\begin{matrix}
w \\
w \\
w \\
2w \\
2w
\end{matrix}
\right),
\quad
\left(
\begin{matrix}
0 \\
0 \\
0 \\
0 \\
0
\end{matrix}
\right)
$$
où $w$ et $m$ sont arbitraires ;
\item
\label{item-equal-equal-down-equal}
$n = 5$, et $m_i, a_{ij}, w_i, g_i$ sont donnés par
$$
\left(
\begin{matrix}
m \\
2m \\
3m \\
4m \\
2m
\end{matrix}
\right),
\quad
\left(
\begin{matrix}
-4w & 2w & 0 & 0 & 0 \\
2w & -4w & 2w & 0 & 0 \\
0 & 2w & -4w & 2w & 0 \\
0 & 0 & 2w & -2w & w \\
0 & 0 & 0 & w & -2w \\
\end{matrix}
\right),
\quad
\left(
\begin{matrix}
2w \\
2w \\
2w \\
w \\
w
\end{matrix}
\right),
\quad
\left(
\begin{matrix}
0 \\
0 \\
0 \\
0 \\
0
\end{matrix}
\right)
$$
où $w$ et $m$ sont arbitraires ;
\item
\label{item-up-equal-equal-up}
$n = 5$, et $m_i, a_{ij}, w_i, g_i$ sont donnés par
$$
\left(
\begin{matrix}
2m \\
2m \\
2m \\
2m \\
m
\end{matrix}
\right),
\quad
\left(
\begin{matrix}
-2w & 2w & 0 & 0 & 0 \\
2w & -4w & 2w & 0 & 0 \\
0 & 2w & -4w & 2w & 0 \\
0 & 0 & 2w & -4w & 4w \\
0 & 0 & 0 & 4w & -8w \\
\end{matrix}
\right),
\quad
\left(
\begin{matrix}
w \\
2w \\
2w \\
2w \\
4w
\end{matrix}
\right),
\quad
\left(
\begin{matrix}
0 \\
0 \\
0 \\
0 \\
0
\end{matrix}
\right)
$$
où $w$ et $m$ sont arbitraires ;
\item
\label{item-up-equal-equal-down}
$n = 5$, et $m_i, a_{ij}, w_i, g_i$ sont donnés par
$$
\left(
\begin{matrix}
m \\
m \\
m \\
m \\
m
\end{matrix}
\right),
\quad
\left(
\begin{matrix}
-2w & 2w & 0 & 0 & 0 \\
2w & -4w & 2w & 0 & 0 \\
0 & 2w & -4w & 2w & 0 \\
0 & 0 & 2w & -4w & 2w \\
0 & 0 & 0 & 2w & -2w \\
\end{matrix}
\right),
\quad
\left(
\begin{matrix}
w \\
2w \\
2w \\
2w \\
w
\end{matrix}
\right),
\quad
\left(
\begin{matrix}
0 \\
0 \\
0 \\
0 \\
0
\end{matrix}
\right)
$$
où $w$ et $m$ sont arbitraires ;
\item
\label{item-down-equal-equal-up}
$n = 5$, et $m_i, a_{ij}, w_i, g_i$ sont donnés par
$$
\left(
\begin{matrix}
m \\
2m \\
2m \\
2m \\
m
\end{matrix}
\right),
\quad
\left(
\begin{matrix}
-4w & 2w & 0 & 0 & 0 \\
2w & -2w & w & 0 & 0 \\
0 & w & -2w & w & 0 \\
0 & 0 & w & -2w & 2w \\
0 & 0 & 0 & 2w & -4w \\
\end{matrix}
\right),
\quad
\left(
\begin{matrix}
2w \\
w \\
w \\
w \\
2w
\end{matrix}
\right),
\quad
\left(
\begin{matrix}
0 \\
0 \\
0 \\
0 \\
0
\end{matrix}
\right)
$$
où $w$ et $m$ sont arbitraires ;
\item
\label{item-quadruple}
$n = 5$, et $m_i, a_{ij}, w_i, g_i$ sont donnés par
$$
\left(
\begin{matrix}
2m \\
m \\
m \\
m \\
m
\end{matrix}
\right),
\quad
\left(
\begin{matrix}
-2w & w & w & w & w \\
w & -2w & 0 & 0 & 0 \\
w & 0 & -2w & 0 & 0 \\
w & 0 & 0 & -2w & 0 \\
w & 0 & 0 & 0 & -2w \\
\end{matrix}
\right),
\quad
\left(
\begin{matrix}
w \\
w \\
w \\
w \\
w
\end{matrix}
\right),
\quad
\left(
\begin{matrix}
0 \\
0 \\
0 \\
0 \\
0
\end{matrix}
\right)
$$
où $w$ et $m$ sont arbitraires ;
\item
\label{item-triple-extended-up}
$n = 5$, et $m_i, a_{ij}, w_i, g_i$ sont donnés par
$$
\left(
\begin{matrix}
m \\
2m \\
2m \\
m \\
m
\end{matrix}
\right),
\quad
\left(
\begin{matrix}
-4w & 2w & 0 & 0 & 0 \\
2w & -2w & w & 0 & 0 \\
0 & w & -2w & w & w \\
0 & 0 & w & -2w & 0 \\
0 & 0 & w & 0 & -2w \\
\end{matrix}
\right),
\quad
\left(
\begin{matrix}
2w \\
w \\
w \\
w \\
w
\end{matrix}
\right),
\quad
\left(
\begin{matrix}
0 \\
0 \\
0 \\
0 \\
0
\end{matrix}
\right)
$$
où $w$ et $m$ sont arbitraires ;
\item
\label{item-triple-extended-down}
$n = 5$, et $m_i, a_{ij}, w_i, g_i$ sont donnés par
$$
\left(
\begin{matrix}
2m \\
2m \\
2m \\
m \\
m
\end{matrix}
\right),
\quad
\left(
\begin{matrix}
-2w & 2w & 0 & 0 & 0 \\
2w & -4w & 2w & 0 & 0 \\
0 & 2w & -4w & 2w & 2w \\
0 & 0 & 2w & -4w & 0 \\
0 & 0 & 2w & 0 & -4w \\
\end{matrix}
\right),
\quad
\left(
\begin{matrix}
w \\
2w \\
2w \\
2w \\
2w
\end{matrix}
\right),
\quad
\left(
\begin{matrix}
0 \\
0 \\
0 \\
0 \\
0
\end{matrix}
\right)
$$
où $w$ et $m$ sont arbitraires ;
\item
\label{item-n-cycle}
$n \geq 6$ et l'on a un $n$-cycle qui généralise (\ref{item-five-cycle}) :
\begin{enumerate}
\item $m_1 = \ldots = m_n = m$,
\item $a_{12} = \ldots = a_{(n - 1) n} = w$, $a_{1n} = w$,
et, pour les autres $i < j$, on a $a_{ij} = 0$,
\item $w_1 = \ldots = w_n = w$
\end{enumerate}
où $w$ et $m$ sont arbitraires ;
\item
\label{item-up-chain-equal-up}
$n \geq 6$ et l'on a une chaîne qui généralise (\ref{item-up-equal-equal-up}) :
\begin{enumerate}
\item $m_1 = \ldots = m_{n - 1} = 2m$, $m_n = m$,
\item $a_{12} = \ldots = a_{(n - 2) (n - 1)} = 2w$, $a_{(n - 1) n} = 4w$,
et, pour les autres $i < j$, on a $a_{ij} = 0$,
\item $w_1 = w$, $w_2 = \ldots = w_{n - 1} = 2w$, $w_n = 4w$
\end{enumerate}
où $w$ et $m$ sont arbitraires ;
\item
\label{item-up-chain-equal-down}
$n \geq 6$ et l'on a une chaîne qui généralise (\ref{item-up-equal-equal-down}) :
\begin{enumerate}
\item $m_1 = \ldots = m_n = m$,
\item $a_{12} = \ldots = a_{(n - 1) n} = w$,
et, pour les autres $i < j$, on a $a_{ij} = 0$,
\item $w_1 = w$, $w_2 = \ldots = w_{n - 1} = 2w$, $w_n = w$
\end{enumerate}
où $w$ et $m$ sont arbitraires ;
\item
\label{item-down-chain-equal-up}
$n \geq 6$ et l'on a une chaîne qui généralise (\ref{item-down-equal-equal-up}) :
\begin{enumerate}
\item $m_1 = w$, $w_2 = \ldots = m_{n - 1} = 2m$, $m_n = m$,
\item $a_{12} = 2w$, $a_{23} = \ldots = a_{(n - 2) (n - 1)} = w$,
$a_{(n - 1) n} = 2w$ et, pour les autres $i < j$, on a $a_{ij} = 0$,
\item $w_1 = 2w$, $w_2 = \ldots = w_{n - 1} = w$, $w_n = 2w$
\end{enumerate}
où $w$ et $m$ sont arbitraires ;
\item
\label{item-Dn-extended-up}
$n \geq 6$ et l'on a un type qui généralise (\ref{item-triple-extended-up}) :
\begin{enumerate}
\item $m_1 = m$, $m_2 = \ldots = m_{n - 3} = 2m$, $m_{n - 1} = m_n = m$,
\item $a_{12} = 2w$, $a_{23} = \ldots = a_{(n - 2) (n - 1)} = w$,
$a_{(n - 2) n} = w$ et, pour les autres $i < j$, on a $a_{ij} = 0$,
\item $w_1 = 2w$, $w_2 = \ldots = w_n = w$
\end{enumerate}
où $w$ et $m$ sont arbitraires ;
\item
\label{item-Dn-extended-down}
$n \geq 6$ et l'on a un type qui généralise (\ref{item-triple-extended-down}) :
\begin{enumerate}
\item $m_1 = \ldots = m_{n - 3} = 2m$, $m_{n - 1} = m_n = m$,
\item $a_{12} = \ldots = a_{(n - 2) (n - 1)} = 2w$,
$a_{(n - 2) n} = 2w$ et, pour les autres $i < j$, on a $a_{ij} = 0$,
\item $w_1 = w$, $w_2 = \ldots = w_n = 2w$
\end{enumerate}
où $w$ et $m$ sont arbitraires ;
\item
\label{item-double-triple}
$n \geq 6$ et l'on a un type qui généralise (\ref{item-quadruple}) :
\begin{enumerate}
\item $m_1 = m_2 = m$, $m_3 = \ldots = m_{n - 2} = 2m$, $m_{n - 1} = m_n = m$,
\item $a_{13} = w$, $a_{23} = \ldots = a_{(n - 2) (n - 1)} = w$,
$a_{(n - 2) n} = w$ et, pour les autres $i < j$, on a $a_{ij} = 0$,
\item $w_1 = \ldots = w_n = w$,
\end{enumerate}
où $w$ et $m$ sont arbitraires ;
\item
\label{item-E6-completed}
$n = 7$, et $m_i, a_{ij}, w_i, g_i$ sont donnés par
$$
\left(
\begin{matrix}
m \\
2m \\
3m \\
m \\
2m \\
m \\
2m
\end{matrix}
\right),
\quad
\left(
\begin{matrix}
-2w & w & 0 & 0 & 0 & 0 & 0 \\
w & -2w & w & 0 & 0 & 0 & 0 \\
0 & w & -2w & 0 & w & 0 & w \\
0 & 0 & 0 & -2w & w & 0 & 0 \\
0 & 0 & w & w & -2w & 0 & 0 \\
0 & 0 & 0 & 0 & 0 & -2w & w \\
0 & 0 & w & 0 & 0 & w & -2w
\end{matrix}
\right),
\quad
\left(
\begin{matrix}
w \\
w \\
w \\
w \\
w \\
w \\
w
\end{matrix}
\right),
\quad
\left(
\begin{matrix}
0 \\
0 \\
0 \\
0 \\
0 \\
0 \\
0
\end{matrix}
\right)
$$
où $w$ et $m$ sont arbitraires ;
\item
\label{item-E7-completed}
$n = 8$, et $m_i, a_{ij}, w_i, g_i$ sont donnés par
$$
\left(
\begin{matrix}
m \\
2m \\
3m \\
4m \\
3m \\
2m \\
m \\
2m
\end{matrix}
\right),
\quad
\left(
\begin{matrix}
-2w & w & 0 & 0 & 0 & 0 & 0 & 0 \\
w & -2w & w & 0 & 0 & 0 & 0 & 0 \\
0 & w & -2w & w & 0 & 0 & 0 & 0 \\
0 & 0 & w & -2w & w & 0 & 0 & w \\
0 & 0 & 0 & w & -2w & w & 0 & 0 \\
0 & 0 & 0 & 0 & w & -2w & w & 0 \\
0 & 0 & 0 & 0 & 0 & w & -2w & 0 \\
0 & 0 & 0 & w & 0 & 0 & 0 & -2w \\
\end{matrix}
\right),
\quad
\left(
\begin{matrix}
w \\
w \\
w \\
w \\
w \\
w \\
w \\
w
\end{matrix}
\right),
\quad
\left(
\begin{matrix}
0 \\
0 \\
0 \\
0 \\
0 \\
0 \\
0 \\
0
\end{matrix}
\right)
$$
où $w$ et $m$ sont arbitraires ;
\item
\label{item-E8-completed}
$n = 9$, et $m_i, a_{ij}, w_i, g_i$ sont donnés par
$$
\left(
\begin{matrix}
m \\
2m \\
3m \\
4m \\
5m \\
6m \\
4m \\
2m \\
3m
\end{matrix}
\right),
\quad
\left(
\begin{matrix}
-2w & w & 0 & 0 & 0 & 0 & 0 & 0 & 0 \\
w & -2w & w & 0 & 0 & 0 & 0 & 0 & 0 \\
0 & w & -2w & w & 0 & 0 & 0 & 0 & 0 \\
0 & 0 & w & -2w & w & 0 & 0 & 0 & 0 \\
0 & 0 & 0 & w & -2w & w & 0 & 0 & 0 \\
0 & 0 & 0 & 0 & w & -2w & w & 0 & w \\
0 & 0 & 0 & 0 & 0 & w & -2w & w & 0 \\
0 & 0 & 0 & 0 & 0 & 0 & w & -2w & 0 \\
0 & 0 & 0 & 0 & 0 & w & 0 & 0 & -2w \\
\end{matrix}
\right),
\quad
\left(
\begin{matrix}
w \\
w \\
w \\
w \\
w \\
w \\
w \\
w \\
w
\end{matrix}
\right),
\quad
\left(
\begin{matrix}
0 \\
0 \\
0 \\
0 \\
0 \\
0 \\
0 \\
0 \\
0
\end{matrix}
\right)
$$
où $w$ et $m$ sont arbitraires.
\end{enumerate}
\end{lemma}

\begin{proof}
Cela est démontré dans la section \ref{section-classify-proper-subgraphs}.
Voir la discussion au
début de la section \ref{section-classify-proper-subgraphs}.
\end{proof}






\section{Bornes pour les invariants des types numériques}
\label{section-towards-classification}

\noindent
Dans notre démonstration de la réduction semi-stable pour les courbes, nous utiliserons une borne
sur les groupes de Picard des types numériques de genre $g$, que nous établirons
dans cette section.

\begin{lemma}
\label{lemma-bound-neighbours}
Soient $n, m_i, a_{ij}, w_i, g_i$ un type numérique de genre $g$.
Pour $i, j$ tels que $a_{ij} > 0$, on a
$m_ia_{ij} \leq m_j|a_{jj}|$ et $m_iw_i \leq m_j|a_{jj}|$.
\end{lemma}

\begin{proof}
Pour tout indice $j$, on a $m_j a_{jj} + \sum_{i \not = j} m_ia_{ij} = 0$.
Ainsi, si l'on dispose d'une borne supérieure pour $|a_{jj}|$ et $m_j$, on obtient aussi une
borne supérieure pour les $a_{ij}$ non nuls (donc positifs), ainsi que pour
$m_i$. En rappelant que $w_i$ divise $a_{ij}$, le lecteur vérifiera aisément
le lemme.
\end{proof}

\begin{lemma}
\label{lemma-bound-heart}
Fixons $g \geq 2$. Pour tout type numérique minimal $n, m_i, a_{ij}, w_i, g_i$
de genre $g$ avec $n > 1$, on a :
\begin{enumerate}
\item l'ensemble $J \subset \{1, \ldots, n\}$ des indices non $(-2)$
possède au plus $2g - 2$ éléments ;
\item pour $j \in J$, on a $g_j < g$ ;
\item pour $j \in J$, on a $m_j|a_{jj}| \leq 6g - 6$ ; et
\item pour $j \in J$ et $i \in \{1, \ldots, n\}$,
on a $m_ia_{ij} \leq 6g - 6$.
\end{enumerate}
\end{lemma}

\begin{proof}
Rappelons que $g = 1 + \sum m_j(w_j(g_j - 1) - \frac{1}{2} a_{jj})$.
Pour $j \in J$, la contribution $m_j(w_j(g_j - 1) - \frac{1}{2} a_{jj})$
au genre $g$ est $> 0$, donc $\geq 1/2$. On utilise ici le
lemme \ref{lemma-minus-one},
la définition \ref{definition-type-minus-one},
la définition \ref{definition-type-minimal},
le lemme \ref{lemma-minus-two} et
la définition \ref{definition-type-minus-two} ; dans la suite, nous utiliserons ces
résultats sans le mentionner davantage.
Ainsi, $J$ possède au plus $2(g - 1)$ éléments.
Cela démontre (1).

\medskip\noindent
Rappelons que $-a_{ii} > 0$ pour tout $i$ d'après le lemme \ref{lemma-diagonal-negative}.
Par conséquent, pour $j \in J$, la contribution $m_j(w_j(g_j - 1) - \frac{1}{2} a_{jj})$
au genre $g$ est $> m_jw_j(g_j - 1)$. Ainsi,
$$
g - 1 > m_jw_j(g_j - 1) \Rightarrow g_j < (g - 1)/m_jw_j + 1
$$
Cela implique bien $g_j < g$, ce qui démontre (2).

\medskip\noindent
Pour $j \in J$, si $g_j > 0$, alors la contribution
$m_j(w_j(g_j - 1) - \frac{1}{2} a_{jj})$ au genre $g$
est $\geq -\frac{1}{2}m_ja_{jj}$, et l'on conclut immédiatement
que $m_j|a_{jj}| \leq 2(g - 1)$. Sinon, $a_{jj} = -kw_j$
pour un entier $k \geq 3$ (car $j \in J$), et l'on obtient
$$
m_jw_j(-1 + \frac{k}{2}) \leq g - 1
\Rightarrow
m_jw_j \leq \frac{2(g - 1)}{k - 2}
$$
En reportant ceci dans $a_{jj} = -km_jw_j$, on obtient
$$
m_j|a_{jj}| \leq 2(g - 1) \frac{k}{k - 2} \leq 6(g - 1)
$$
Cela démontre (3).

\medskip\noindent
L'assertion (4) résulte du lemme \ref{lemma-bound-neighbours} et de (3).
\end{proof}

\begin{lemma}
\label{lemma-bound-wm}
Fixons $g \geq 2$. Pour tout type numérique minimal $n, m_i, a_{ij}, w_i, g_i$
de genre $g$, on a $m_i|a_{ij}| \leq 768g$.
\end{lemma}

\begin{proof}
D'après le lemme \ref{lemma-bound-neighbours}, il suffit de montrer que
$m_i|a_{ii}| \leq 768g$ pour tout $i$.
Soit $J \subset \{1, \ldots, n\}$ l'ensemble des indices non $(-2)$,
comme dans le lemme \ref{lemma-bound-heart}. Observons que $J$ est
non vide puisque $g \geq 2$. De plus, $m_j|a_{jj}| \leq 6g$ pour $j \in J$
d'après ce lemme.

\medskip\noindent
Supposons que l'on ait $j \in J$ et une suite $i_1, \ldots, i_7$
d'indices $(-2)$ telle que $a_{ji_1}$ et $a_{i_1i_2}$,
$a_{i_2i_3}$, $a_{i_3i_4}$, $a_{i_4i_5}$, $a_{i_5i_6}$ et
$a_{i_6i_7}$ soient non nuls. Le
lemme \ref{lemma-bound-neighbours} montre alors
que $m_{i_1}w_{i_1} \leq 6g$ et $m_{i_1}a_{ji_1} \leq 6g$.
Comme $i_1$ est un indice $(-2)$, on a $a_{i_1i_1} = -2w_{i_1}$,
et l'on conclut que $m_{i_1}|a_{i_1i_1}| \leq 12g$.
En répétant l'argument, on conclut que
$m_{i_2}w_{i_2} \leq 12g$ et $m_{i_2}a_{i_1i_2} \leq 12g$.
Puis $m_{i_2}|a_{i_2i_2}| \leq 24g$, et ainsi de suite.
Finalement, on conclut que
$m_{i_k}|a_{i_ki_k}| \leq 2^k(6g) \leq 768g$ pour $k = 1, \ldots, 7$.

\medskip\noindent
Soit $I \subset \{1, \ldots, n\} \setminus J$ un sous-ensemble connexe maximal.
Autrement dit, il n'existe aucun sous-ensemble strict non vide
$I' \subset I$ tel que
$a_{i'i} = 0$ pour $i' \in I'$ et $i \in I \setminus I'$,
et $I$ est maximal pour cette propriété. En particulier, puisqu'un
type numérique est connexe par définition, on voit qu'il
existe $j \in J$ et $i \in I$ tels que $a_{ij} > 0$.
En considérant la classification
de tels $I$ dans la proposition \ref{proposition-classify-subgraphs}
et en utilisant le résultat du paragraphe précédent, on voit que
$w_i|a_{ii}| \leq 768g$ pour tout $i \in I$, sauf si $I$ est décrit dans le
lemme \ref{lemma-long} ou le
lemme \ref{lemma-Dn}.
On peut donc supposer que le lieu de non-annulation des $a_{ii'}$, $i, i' \in I$,
a l'une des deux formes suivantes :
$$
\xymatrix{
\bullet \ar@{-}[r] &
\bullet \ar@{-}[r] &
\bullet \ar@{..}[r] &
\bullet \ar@{-}[r] &
\bullet \ar@{-}[r] &
\bullet
}
$$
(laquelle possède 3 sous-cas détaillés dans le lemme \ref{lemma-long}),
ou bien
$$
\xymatrix{
\bullet \ar@{-}[r] &
\bullet \ar@{-}[r] &
\bullet \ar@{..}[r] &
\bullet \ar@{-}[r] &
\bullet \ar@{-}[r] \ar@{-}[d] &
\bullet \\
& & & & \bullet
}
$$
Nous allons démontrer la borne pour le premier sous-cas du
lemme \ref{lemma-long} et laisser les autres cas au lecteur (l'argument
y est presque exactement le même).

\medskip\noindent
Après réindexation, on peut supposer $I = \{1, \ldots, t\} \subset \{1, \ldots, n\}$
et qu'il existe un entier $w$ tel que
$$
w = w_1 = \ldots = w_t =
a_{12} = \ldots = a_{(t - 1) t} =
-\frac{1}{2} a_{i_1i_2} = \ldots = -\frac{1}{2} a_{(t - 1) t}
$$
Les égalités $a_{ii}m_i + \sum_{j \not = i} a_{ij}m_j = 0$ impliquent
que l'on a
$$
2m_2 \geq m_1 + m_3, \ldots, 2m_{t - 1} \geq m_{t - 2} + m_t
$$
On a égalité dans $2m_i \geq m_{i - 1} + m_{i + 1}$
si et seulement si $i$ ne « rencontre » aucun indice autre que
$i - 1$ et $i + 1$. Et si $i$ rencontre effectivement
un autre indice, alors cet indice appartient à $J$ (par maximalité de $I$).
En particulier, l'application
$\{1, \ldots, t\} \to \mathbf{Z}$, $i \mapsto m_i$, est concave.

\medskip\noindent
Soit $m = \max(m_i, i \in \{1, \ldots, t\})$. Alors
$m_i|a_{ii}| \leq 2mw$ pour $i \leq t$, et
notre but est de montrer que $2mw \leq 768g$.
Soient $s$, respectivement $s'$, dans $\{1, \ldots, t\}$, le
plus petit, respectivement le plus grand, indice tel que $m_s = m = m_{s'}$.
La concavité montre que $m_i = m$ pour $s \leq i \leq s'$.
Si $s > 1$, alors on n'a pas égalité dans
$2m_s \geq m_{s - 1} + m_{s + 1}$,
et l'on voit que $s$ rencontre un indice de $J$.
Dans ce cas, $2mw \leq 12g$ d'après le résultat du deuxième paragraphe
de la démonstration.
De même, si $s' < t$, alors $s'$ rencontre un indice de $J$
et l'on obtient également $2mw \leq 12g$.
Mais si $s = 1$ et $s' = t$, alors on conclut
que $a_{ij} = 0$ pour tout $j \in J$ et $i \in \{2, \ldots, t - 1\}$.
Or, puisque nous avons vu qu'il doit exister une paire $(i, j) \in I \times J$
telle que $a_{ij} > 0$, on conclut que cela se produit soit pour
$i = 1$, soit pour $i = t$, et l'on conclut $2mw \leq 12g$
de la même manière qu'auparavant (car $m_1 = m = m_t$ dans ce cas).
\end{proof}

\begin{proposition}
\label{proposition-bound-picard-group}
Soit $g \geq 2$. Pour tout type numérique $T$ de genre $g$
et tout nombre premier $\ell > 768g$, on a
$$
\dim_{\mathbf{F}_\ell} \Pic(T)[\ell] \leq g
$$
où $\Pic(T)$ est défini dans la définition \ref{definition-picard-group}.
Si $T$ est minimal, on a même
$$
\dim_{\mathbf{F}_\ell} \Pic(T)[\ell] \leq g_{top} \leq g
$$
où $g_{top}$ est défini dans la définition \ref{definition-top-genus}.
\end{proposition}

\begin{proof}
Supposons que $T$ soit donné par $n, m_i, a_{ij}, w_i, g_i$.
Si $T$ n'est pas minimal, alors il existe un indice $(-1)$.
Après avoir remplacé $T$ par un type équivalent, on peut supposer que
$n$ est un indice $(-1)$. En appliquant le lemme \ref{lemma-contract-picard-group},
on obtient $\Pic(T) \subset \Pic(T')$, où $T'$
est un type numérique de genre $g$ (lemme \ref{lemma-contract})
à $n - 1$ indices. On conclut donc par récurrence sur $n$,
pourvu que le lemme soit démontré pour les types numériques minimaux.

\medskip\noindent
Supposons que $T$ soit un type numérique minimal de genre $\geq 2$.
Observons que $g_{top} \leq g$ d'après le lemme \ref{lemma-genus-nonnegative}.
Si $A = (a_{ij})$, on a $\Pic(T) \subset \Coker(A)$ d'après le
lemme \ref{lemma-picard-T-and-A}. Il suffit donc de démontrer le lemme
pour $\Coker(A)$.
Le lemme \ref{lemma-bound-wm} montre que $m_i|a_{ij}| \leq 768g$ pour
tous $i, j$.
Le résultat découle donc du lemme \ref{lemma-recurring-symmetric-integer}.
\end{proof}







\section{Modèles}
\label{section-models}

\noindent
Dans ce chapitre, $R$ désignera un anneau de valuation discrète et $K$ son
corps des fractions. Si nécessaire, nous noterons $\pi \in R$ une
uniformisante et $k = R/(\pi)$ son corps résiduel.

\medskip\noindent
Soit $V$ un $K$-schéma algébrique
(Variétés, définition \ref{varieties-definition-algebraic-scheme}).
Par {\it modèle} de $V$, on entendra un morphisme plat de type
fini\footnote{Il est parfois utile d'autoriser les modèles qui sont
localement de type fini sur $R$, mais nous nous en préoccuperons
le moment venu.}
$X \to \Spec(R)$ muni d'un
isomorphisme $V \to X_K = X \times_{\Spec(R)} \Spec(K)$. Souvent, nous
identifierons $V$ et la fibre générique $X_K$ de $X$, et
écrirons simplement $V = X_K$.
La fibre spéciale est $X_k = X \times_{\Spec(R)} \Spec(k)$.
Un {\it morphisme de modèles $X \to X'$ de $V$}
est un morphisme $X \to X'$ de schémas sur $R$ qui induit
l'identité sur $V$.

\medskip\noindent
On dira que {\it $X$ est un modèle propre de $V$} si $X$ est un modèle
de $V$ et si le morphisme structural $X \to \Spec(R)$ est propre.
De même pour les modèles séparés, les modèles lisses et d'autres encore,
à ajouter ici.
On dira que {\it $X$ est un modèle régulier de $V$} si $X$ est un modèle
de $V$ et si $X$ est un schéma régulier.
De même pour les modèles normaux, les modèles réduits et d'autres encore,
à ajouter ici.

\medskip\noindent
Soit $R \subset R'$ une extension d'anneaux de valuation discrète
(Compléments d'algèbre, définition
\ref{more-algebra-definition-extension-discrete-valuation-rings}).
Celle-ci induit une extension $K'/K$ des corps des fractions.
Étant donné un schéma algébrique $V$ sur $K$, notons $V'$ le
changement de base $V \times_{\Spec(K)} \Spec(K')$. On a alors
un foncteur
$$
\text{modèles de }V\text{ sur }R
\longrightarrow
\text{modèles de }V'\text{ sur }R'
$$
qui envoie $X$ sur $X \times_{\Spec(R)} \Spec(R')$.

\begin{lemma}
\label{lemma-closure-is-model}
Soit $V_1 \to V_2$ une immersion fermée de schémas algébriques sur $K$.
Si $X_2$ est un modèle de $V_2$, alors l'image schématique
de $V_1 \to X_2$ est un modèle de $V_1$.
\end{lemma}

\begin{proof}
En utilisant
Morphismes, lemme \ref{morphisms-lemma-quasi-compact-scheme-theoretic-image} et
l'exemple \ref{morphisms-example-scheme-theoretic-image},
on se ramène à l'énoncé algébrique suivant.
Soit $A_1$ une $R$-algèbre de type fini, plate sur $R$.
Soit $A_1 \otimes_R K \to B_2$ une surjection. Alors
$A_2 = A_1 / \Ker(A_1 \to B_2)$ est une $R$-algèbre de type fini,
plate sur $R$, telle que $B_2 = A_2 \otimes_R K$.
Nous omettons la démonstration détaillée ; utiliser
Compléments d'algèbre, lemme \ref{more-algebra-lemma-dedekind-torsion-free-flat},
pour montrer que $A_2$ est plate.
\end{proof}

\begin{lemma}
\label{lemma-normalization}
Soit $X$ un modèle d'une variété géométriquement normale $V$ sur $K$.
Alors la normalisation $\nu : X^\nu \to X$ est finie et
le changement de base de $X^\nu$ à la complétion $R^\wedge$
est la normalisation du changement de base de $X$. De plus, pour
tout $x \in X^\nu$, la complétion de $\mathcal{O}_{X^\nu, x}$
est normale.
\end{lemma}

\begin{proof}
Observons que $R^\wedge$ est un anneau de valuation discrète
(Compléments d'algèbre, lemme \ref{more-algebra-lemma-completion-dvr}).
Posons $Y = X \times_{\Spec(R)} \Spec(R^\wedge)$.
Puisque $R^\wedge$ est un anneau de valuation discrète, on voit que
$$
Y \setminus Y_k =
Y \times_{\Spec(R^\wedge)} \Spec(K^\wedge) =
V \times_{\Spec(K)} \Spec(K^\wedge)
$$
où $K^\wedge$ est le corps des fractions de $R^\wedge$.
Comme $V$ est géométriquement normale, on trouve qu'il s'agit
d'un schéma normal. La première partie du lemme résulte donc de
Résolution des surfaces, lemme \ref{resolve-lemma-normalization-completion}.

\medskip\noindent
Pour démontrer la seconde partie, on peut supposer $X$ et $Y$ normaux
(par la première partie). Si $x$ appartient à la fibre générique, alors
$\mathcal{O}_{X, x} = \mathcal{O}_{V, x}$ est un anneau local normal
essentiellement de type fini sur un corps. Un tel anneau est
excellent (Compléments d'algèbre, proposition
\ref{more-algebra-proposition-ubiquity-excellent}).
Si $x$ est un point de la fibre spéciale d'image $y \in Y$, alors
$\mathcal{O}_{X, x}^\wedge = \mathcal{O}_{Y, y}^\wedge$
par Résolution des surfaces, lemme \ref{resolve-lemma-iso-completions}.
Dans ce cas, $\mathcal{O}_{Y, y}$ est un domaine local normal excellent,
par la même référence que précédemment, puisque $R^\wedge$ est excellent.
Si $B$ est un domaine local normal excellent, alors la complétion
$B^\wedge$ est normale (car $B \to B^\wedge$ est régulier et le
lemme \ref{more-algebra-lemma-normal-goes-up} de Compléments d'algèbre s'applique).
Ceci achève la démonstration.
\end{proof}

\begin{lemma}
\label{lemma-regular}
Soit $X$ un modèle d'une courbe lisse $C$ sur $K$. Alors
il existe une résolution des singularités de $X$,
et toute résolution est un modèle de $C$.
\end{lemma}

\begin{proof}
Vérifions que la condition (4) du théorème de Lipman
(Résolution des surfaces, théorème \ref{resolve-theorem-resolve}) est satisfaite.
Cela découle clairement du lemme \ref{lemma-normalization},
à l'exception de l'assertion selon laquelle $X^\nu$ ne possède qu'un nombre fini de
points singuliers. Pour le voir, on peut utiliser le fait que $R$ est J-2 d'après
Compléments d'algèbre, proposition \ref{more-algebra-proposition-ubiquity-J-2},
de sorte que le lieu non singulier est ouvert dans $X^\nu$.
Comme $X^\nu$ est normal de dimension $\leq 2$, les points singuliers
sont fermés ; le fait que le lieu singulier soit fermé signifie donc
qu'ils sont en nombre fini (puisque $X$ est quasi-compact).
Observons que toute résolution de $X$ est une modification de $X$
(Résolution des surfaces, définition \ref{resolve-definition-resolution}).
Celle-ci est un isomorphisme au-dessus du lieu normal de $X$ d'après Variétés, lemme
\ref{varieties-lemma-modification-normal-iso-over-codimension-1}.
Puisque l'ensemble des points normaux contient
$C = X_K$, on conclut que toute résolution est un modèle de $C$.
\end{proof}

\begin{definition}
\label{definition-minimal-model}
Soit $C$ une courbe projective lisse sur $K$ telle que
$H^0(C, \mathcal{O}_C) = K$. Un {\it modèle minimal}
est un modèle propre régulier $X$ de $C$ tel que
$X$ ne contienne aucune courbe exceptionnelle de première espèce
(Résolution des surfaces, section \ref{resolve-section-minus-one}).
\end{definition}

\noindent
En toute rigueur, un tel objet devrait être appelé modèle propre régulier minimal,
voire modèle projectif régulier relativement minimal. Mais tant que
nous nous limitons aux modèles sur des anneaux de valuation discrète (comme nous le ferons
dans ce chapitre), aucune confusion ne devrait en résulter.

\medskip\noindent
Les modèles minimaux existent toujours
(proposition \ref{proposition-exists-minimal-model}) et sont uniques
lorsque le genre est $> 0$ (lemme \ref{lemma-minimal-model-unique}).

\begin{lemma}
\label{lemma-pre-exists-minimal-model}
\begin{slogan}
Un modèle propre régulier d'une courbe s'obtient par éclatements successifs
à partir d'un modèle minimal
\end{slogan}
Soit $C$ une courbe projective lisse sur $K$ telle que
$H^0(C, \mathcal{O}_C) = K$. Si $X$ est un modèle propre régulier
de $C$, alors il existe une suite de morphismes
$$
X = X_m \to X_{m - 1} \to \ldots \to X_1 \to X_0
$$
de modèles propres réguliers de $C$, telle que chaque morphisme soit une
contraction d'une courbe exceptionnelle de première espèce et que
$X_0$ soit un modèle minimal.
\end{lemma}

\begin{proof}
D'après Résolution des surfaces, lemme \ref{resolve-lemma-regular-dim-2-projective},
$X$ est projectif sur $R$. Par conséquent, $X$ possède un faisceau inversible
ample d'après
Compléments sur les morphismes, lemme \ref{more-morphisms-lemma-projective}
(nous l'utiliserons ci-dessous).
Soit $E \subset X$ une courbe exceptionnelle de première espèce.
Voir Résolution des surfaces, section \ref{resolve-section-minus-one}.
D'après Résolution des surfaces, lemme \ref{resolve-lemma-contract-ample},
on peut contracter $E$ par un morphisme $X \to X'$ tel que $X'$ soit
régulier et projectif sur $R$. Il est clair que le nombre de
composantes irréductibles de $X'_k$ est inférieur d'une unité au
nombre de composantes irréductibles de $X_k$. On ne peut donc
effectuer qu'un nombre fini de ces contractions avant
d'obtenir un modèle minimal.
\end{proof}

\begin{proposition}
\label{proposition-exists-minimal-model}
Soit $C$ une courbe projective lisse sur $K$ telle que
$H^0(C, \mathcal{O}_C) = K$. Il existe un modèle minimal.
\end{proposition}

\begin{proof}
Choisissons une immersion fermée $C \to \mathbf{P}^n_K$. Soit
$X$ l'image schématique de $C \to \mathbf{P}^n_R$.
Alors $X \to \Spec(R)$ est un modèle projectif de $C$ d'après le
lemme \ref{lemma-closure-is-model}.
D'après le lemme \ref{lemma-regular}, il existe une résolution
des singularités $X' \to X$, et $X'$ est un modèle de $C$.
Alors $X' \to \Spec(R)$ est propre comme composé de morphismes propres.
On peut alors appliquer le lemme \ref{lemma-pre-exists-minimal-model}
pour obtenir un modèle minimal.
\end{proof}





\section{La géométrie d'un modèle régulier}
\label{section-special-fibre}

\noindent
Dans cette section, nous décrivons la géométrie d'un modèle propre régulier $X$ d'une
courbe projective lisse $C$ sur $K$ telle que $H^0(C, \mathcal{O}_C) = K$.

\begin{lemma}
\label{lemma-divisor-special-fiber}
Soit $X$ un modèle régulier d'une courbe lisse $C$ sur $K$.
\begin{enumerate}
\item la fibre spéciale $X_k$ est un diviseur de Cartier effectif sur $X$,
\item chaque composante irréductible $C_i$ de $X_k$ est un diviseur de
Cartier effectif sur $X$,
\item $X_k = \sum m_i C_i$ (somme de diviseurs de Cartier effectifs),
où $m_i$ est la multiplicité de $C_i$ dans $X_k$,
\item $\mathcal{O}_X(X_k) \cong \mathcal{O}_X$.
\end{enumerate}
\end{lemma}

\begin{proof}
Rappelons que $R$ est un anneau de valuation discrète d'uniformisante $\pi$
et de corps résiduel $k = R/(\pi)$. Puisque $X \to \Spec(R)$ est plat,
l'élément $\pi$ est un non-diviseur de zéro localement sur les ouverts affines de $X$
(voir Compléments d'algèbre, lemme
\ref{more-algebra-lemma-dedekind-torsion-free-flat}). Ainsi,
si $U = \Spec(A) \subset X$ est un ouvert affine, alors
$$
X_k \cap U = U_k = \Spec(A \otimes_R k) = \Spec(A/\pi A)
$$
et $\pi$ est un non-diviseur de zéro dans $A$.
Par conséquent, $X_k = V(\pi)$ est un diviseur de Cartier effectif d'après
Diviseurs, lemme \ref{divisors-lemma-characterize-effective-Cartier-divisor}.
Cela démontre (1).

\medskip\noindent
La discussion précédente montre que la paire $(\mathcal{O}_X(X_k), 1)$
est isomorphe à la paire $(\mathcal{O}_X, \pi)$, ce qui démontre (4).

\medskip\noindent
D'après Diviseurs, lemme \ref{divisors-lemma-effective-Cartier-divisor-is-a-sum},
il existe des diviseurs de Cartier effectifs intègres deux à deux distincts
$D_i \subset X$ et des entiers $a_i \geq 0$ tels que $X_k = \sum a_i D_i$.
On peut éliminer les diviseurs $D_i$ tels que $a_i = 0$. Il est alors
clair (d'après la définition de l'addition des diviseurs de Cartier
effectifs) que $X_k = \bigcup D_i$ ensemblistement. Ainsi, les $C_i = D_i$
sont les composantes irréductibles de $X_k$, ce qui démontre (2).
Soit $\xi_i$ le point générique de $C_i$.
Alors $\mathcal{O}_{X, \xi_i}$ est un anneau de valuation discrète
(Diviseurs, lemme \ref{divisors-lemma-integral-effective-Cartier-divisor-dvr}).
L'uniformisante $\pi_i \in \mathcal{O}_{X, \xi_i}$ est une équation locale
de $C_i$, et l'image de $\pi$ est une équation locale de $X_k$.
Comme $X_k = \sum a_i C_i$, on voit que $\pi$ et $\pi_i^{a_i}$
engendrent le même idéal dans $\mathcal{O}_{X, \xi_i}$.
D'autre part, la multiplicité de $C_i$ dans $X_k$ est
$$
m_i =
\text{longueur}_{\mathcal{O}_{C_i, \xi_i}} \mathcal{O}_{X_k, \xi_i} =
\text{longueur}_{\mathcal{O}_{C_i, \xi_i}} \mathcal{O}_{X, \xi_i}/(\pi) =
\text{longueur}_{\mathcal{O}_{C_i, \xi_i}} \mathcal{O}_{X, \xi_i}/(\pi_i^{a_i}) =
a_i
$$
Voir Homologie de Chow, définition
\ref{chow-definition-cycle-associated-to-closed-subscheme}.
Ainsi $a_i = m_i$, ce qui démontre (3).
\end{proof}

\begin{lemma}
\label{lemma-gorenstein}
Soit $X$ un modèle régulier d'une courbe lisse $C$ sur $K$. Alors
\begin{enumerate}
\item $X \to \Spec(R)$ est un morphisme de Gorenstein de dimension relative $1$,
\item chacune des composantes irréductibles $C_i$ de $X_k$ est de Gorenstein.
\end{enumerate}
\end{lemma}

\begin{proof}
Comme $X \to \Spec(R)$ est plat, pour démontrer (1),
il suffit de montrer que les fibres sont de Gorenstein
(Dualité pour les schémas, lemme \ref{duality-lemma-gorenstein-morphism}).
La fibre générique est une courbe lisse, qui est régulière et donc de Gorenstein
(Dualité pour les schémas, lemme \ref{duality-lemma-regular-gorenstein}).
Pour la fibre spéciale $X_k$, on utilise le fait qu'elle est un diviseur de
Cartier effectif sur un schéma régulier (donc de Gorenstein), et qu'elle est donc
de Gorenstein, par exemple d'après Complexes dualisants, lemme
\ref{dualizing-lemma-gorenstein-divide-by-nonzerodivisor}.
Les courbes $C_i$ sont de Gorenstein par le même argument.
\end{proof}

\begin{situation}
\label{situation-regular-model}
Soit $R$ un anneau de valuation discrète de corps des fractions $K$,
de corps résiduel $k$ et d'uniformisante $\pi$.
Soit $C$ une courbe projective lisse sur $K$ telle que $H^0(C, \mathcal{O}_C) = K$.
Soit $X$ un modèle propre régulier de $C$.
Soient $C_1, \ldots, C_n$ les composantes irréductibles de la fibre
spéciale $X_k$. Écrivons $X_k = \sum m_i C_i$ comme dans le
lemme \ref{lemma-divisor-special-fiber}.
\end{situation}

\begin{lemma}
\label{lemma-regular-model-connected}
Dans la situation \ref{situation-regular-model}, la fibre spéciale $X_k$ est connexe.
\end{lemma}

\begin{proof}
C'est une conséquence de Compléments sur les morphismes, lemme
\ref{more-morphisms-lemma-geometrically-connected-fibres-towards-normal}.
\end{proof}

\begin{lemma}
\label{lemma-regular-model-pic}
Dans la situation \ref{situation-regular-model}, il existe une suite exacte
$$
0 \to \mathbf{Z} \to \mathbf{Z}^{\oplus n} \to
\Pic(X) \to \Pic(C) \to 0
$$
où la première application envoie $1$ sur $(m_1, \ldots, m_n)$ et la seconde
envoie le $i$-ème vecteur de base sur $\mathcal{O}_X(C_i)$.
\end{lemma}

\begin{proof}
Observons que $C \subset X$ est un sous-schéma ouvert. L'application de restriction
$\Pic(X) \to \Pic(C)$ est surjective d'après
Diviseurs, lemme \ref{divisors-lemma-extend-invertible-module}.
Soit $\mathcal{L}$ un $\mathcal{O}_X$-module inversible
tel qu'il existe un isomorphisme $s : \mathcal{O}_C \to \mathcal{L}|_C$.
Alors $s$ est une section méromorphe régulière de $\mathcal{L}$,
et l'on voit que $\text{div}_\mathcal{L}(s) = \sum a_i C_i$
pour certains $a_i \in \mathbf{Z}$
(Diviseurs, définition \ref{divisors-definition-divisor-invertible-sheaf}).
D'après Diviseurs, lemme \ref{divisors-lemma-normal-c1-injective}
(et le fait que $X$ est normal),
on conclut que $\mathcal{L} = \mathcal{O}_X(\sum a_iC_i)$.
Enfin, supposons que $\mathcal{O}_X(\sum a_i C_i) \cong \mathcal{O}_X$.
Il existe alors un élément $g$ du corps des fonctions de $X$
tel que $\text{div}_X(g) = \sum a_i C_i$. En particulier, la fonction rationnelle
$g$ n'a ni zéro ni pôle sur la fibre générique $C$ de $X$.
Comme $C$ est un schéma normal, ceci implique $g \in H^0(C, \mathcal{O}_C) = K$.
Ainsi $g = \pi^a u$ pour un certain $a \in \mathbf{Z}$ et un certain $u \in R^*$.
On conclut que $\text{div}_X(g) = a \sum m_i C_i$, ce qui achève
la démonstration.
\end{proof}

\noindent
Dans la situation \ref{situation-regular-model}, pour tout
$\mathcal{O}_X$-module inversible $\mathcal{L}$ et tout $i$, on obtient un entier
$$
\deg(\mathcal{L}|_{C_i}) =
\chi(C_i, \mathcal{L}|_{C_i}) - \chi(C_i, \mathcal{O}_{C_i})
$$
en prenant le degré de la restriction de $\mathcal{L}$ à $C_i$
relativement au corps de base $k$\footnote{Observons qu'il peut arriver
que le corps $\kappa_i = H^0(C_i, \mathcal{O}_{C_i})$ soit strictement plus grand
que $k$. Dans ce cas, tout module inversible sur $C_i$ a un
degré (au sens défini ci-dessus) divisible par $[\kappa_i : k]$.},
comme dans Variétés, section \ref{varieties-section-divisors-curves}.

\begin{lemma}
\label{lemma-intersection-pairing}
Dans la situation \ref{situation-regular-model}, étant donnés un
$\mathcal{O}_X$-module inversible $\mathcal{L}$ et
$a = (a_1, \ldots, a_n) \in \mathbf{Z}^{\oplus n}$, on définit
$$
\langle a, \mathcal{L} \rangle = \sum a_i\deg(\mathcal{L}|_{C_i})
$$
Alors $\langle , \rangle$ est bilinéaire et, pour
$b = (b_1, \ldots, b_n) \in \mathbf{Z}^{\oplus n}$, on a
$$
\left\langle a, \mathcal{O}_X(\sum b_i C_i) \right\rangle =
\left\langle b, \mathcal{O}_X(\sum a_i C_i) \right\rangle
$$
\end{lemma}

\begin{proof}
La bilinéarité résulte immédiatement de la définition et de
Variétés, lemme \ref{varieties-lemma-degree-tensor-product}.
Pour démontrer la symétrie, il suffit de supposer que $a$ et $b$ sont des
vecteurs de la base canonique de $\mathbf{Z}^{\oplus n}$.
Il suffit donc de démontrer que
$$
\deg(\mathcal{O}_X(C_j)|_{C_i}) = \deg(\mathcal{O}_X(C_i)|_{C_j})
$$
pour tous $1 \leq i, j \leq n$. Si $i = j$, il n'y a rien à démontrer.
Si $i \not = j$, alors la section canonique $1$ de $\mathcal{O}_X(C_j)$
se restreint en une section non nulle (donc régulière) de $\mathcal{O}_X(C_j)|_{C_i}$,
dont le schéma des zéros est exactement $C_i \cap C_j$ (intersection schématique).
Autrement dit, $C_i \cap C_j$ est un diviseur de Cartier effectif sur $C_i$,
et
$$
\deg(\mathcal{O}_X(C_j)|_{C_i}) = \deg(C_i \cap C_j)
$$
d'après Variétés, lemme \ref{varieties-lemma-degree-effective-Cartier-divisor}.
Par symétrie, on obtient la même (!) formule pour l'autre membre,
ce qui achève la démonstration.
\end{proof}

\noindent
Dans la situation \ref{situation-regular-model}, il est souvent commode de considérer
$\mathbf{Z}^{\oplus n}$ comme le groupe abélien libre sur l'ensemble
$\{C_1, \ldots, C_n\}$. Nous noterons un élément de ce groupe
sous la forme $\sum a_i C_i$ ; nous le considérons ici comme une somme formelle, bien que,
de manière équivalente, on puisse (et nous le ferons parfois)
considérer une telle somme comme un diviseur de Weil sur $X$
à support dans la fibre spéciale $X_k$. Le
lemme \ref{lemma-intersection-pairing}
nous permet alors de définir une forme bilinéaire symétrique $(\ \cdot\ )$
sur ce groupe abélien libre par la règle
\begin{equation}
\label{equation-form}
\left(\sum a_i C_i \cdot \sum b_j C_j\right) =
\left\langle a, \mathcal{O}_X(\sum b_j C_j) \right\rangle =
\left\langle b, \mathcal{O}_X(\sum a_i C_i) \right\rangle
\end{equation}
Nous allons établir quelques propriétés de cette forme bilinéaire.

\begin{lemma}
\label{lemma-properties-form}
Dans la situation \ref{situation-regular-model}, la forme bilinéaire symétrique
(\ref{equation-form}) possède les propriétés suivantes :
\begin{enumerate}
\item $(C_i \cdot C_j) \geq 0$ si $i \not = j$, avec égalité si et seulement
si $C_i \cap C_j = \emptyset$,
\item $(\sum m_i C_i \cdot C_j) = 0$,
\item il n'existe aucun sous-ensemble propre non vide $I \subset \{1, \ldots, n\}$
tel que $(C_i \cdot C_j) = 0$ pour $i \in I$, $j \not \in I$,
\item $(\sum a_i C_i \cdot \sum a_i C_i) \leq 0$, avec égalité si et
seulement s'il existe $q \in \mathbf{Q}$ tel que $a_i = qm_i$
pour $i = 1, \ldots, n$.
\end{enumerate}
\end{lemma}

\begin{proof}
Dans la démonstration du lemme \ref{lemma-intersection-pairing}, nous avons vu que
$(C_i \cdot C_j) = \deg(C_i \cap C_j)$ si $i \not = j$. Ce nombre est
$\geq 0$ et il est $> 0 $ si et seulement si $C_i \cap C_j \not = \emptyset$.
Ceci démontre (1).

\medskip\noindent
Démonstration de (2). Cela résulte du fait que, d'après le
lemme \ref{lemma-divisor-special-fiber}, le faisceau inversible associé à
$\sum m_i C_i$ est trivial et que le faisceau trivial est de degré zéro.

\medskip\noindent
Démonstration de (3). Cette assertion exprime le fait que $X_k$ est connexe
(lemme \ref{lemma-regular-model-connected})
au moyen de la description des produits d'intersection donnée dans la démonstration de (1).

\medskip\noindent
La partie (4) résulte de (1), (2) et (3) par le
lemme \ref{lemma-recurring-symmetric-real}.
\end{proof}

\begin{lemma}
\label{lemma-multiple-fibre-normal-bundle}
Dans la situation \ref{situation-regular-model}, posons $d = \gcd(m_1, \ldots, m_n)$
et soit $D = \sum (m_i/d)C_i$ considéré comme un diviseur de Cartier effectif.
Alors $\mathcal{O}_X(D)$ est d'ordre divisant $d$ dans $\Pic(X)$,
et $\mathcal{C}_{D/X}$ est un $\mathcal{O}_D$-module inversible
d'ordre divisant $d$ dans $\Pic(D)$.
\end{lemma}

\begin{proof}
On a
$$
\mathcal{O}_X(D)^{\otimes d} = \mathcal{O}_X(dD) =
\mathcal{O}_X(X_k) = \mathcal{O}_X
$$
d'après le lemme \ref{lemma-divisor-special-fiber}.
On conclut puisque $\mathcal{C}_{D/X}$ est l'image inverse de
$\mathcal{O}_X(-D)$.
\end{proof}

\begin{lemma}
\label{lemma-regular-model-field}
\begin{reference}
\cite[Lemma 2.6]{Artin-Winters}
\end{reference}
Dans la situation \ref{situation-regular-model}, posons $d = \gcd(m_1, \ldots, m_n)$.
Soit $D = \sum (m_i/d) C_i$ considéré comme un diviseur de Cartier effectif. Alors il existe
une suite de diviseurs de Cartier effectifs
$$
(X_k)_{red} = Z_0 \subset Z_1 \subset \ldots \subset Z_m = D
$$
telle que $Z_j = Z_{j - 1} + C_{i_j}$ pour un certain $i_j \in \{1, \ldots, n\}$
pour $j = 1, \ldots, m$, et telle que $H^0(Z_j, \mathcal{O}_{Z_j})$
soit un corps fini sur $k$ pour $j = 0, \ldots m$.
\end{lemma}

\begin{proof}
Le réduit $D_{red} = (X_k)_{red} = \sum C_i$ est connexe
(lemme \ref{lemma-regular-model-connected}) et propre sur $k$. Par conséquent,
$H^0(D_{red}, \mathcal{O})$ est un corps et une extension finie de
$k$ d'après Variétés, lemme
\ref{varieties-lemma-proper-geometrically-reduced-global-sections}.
Le résultat est donc vrai pour $Z_0 = D_{red} = (X_k)_{red}$.
Supposons que nous ayons déjà construit
$$
(X_k)_{red} = Z_0 \subset Z_1 \subset \ldots \subset Z_t \subset D
$$
avec $Z_j = Z_{j - 1} + C_{i_j}$ pour un certain $i_j \in \{1, \ldots, n\}$
pour $j = 1, \ldots, t$, et de sorte que $H^0(Z_j, \mathcal{O}_{Z_j})$
soit un corps fini sur $k$ pour $j = 0, \ldots, t$.
Écrivons $Z_t = \sum a_i C_i$ avec $1 \leq a_i \leq m_i/d$.
Si $a_i = m_i/d$ pour tout $i$, alors $Z_t = D$ et le lemme est démontré.
Sinon, $a_i < m_i/d$ pour un certain $i$, et il s'ensuit que
$(Z_t \cdot Z_t) < 0$ d'après le lemme \ref{lemma-properties-form}. Cela signifie que
$(D - Z_t \cdot Z_t) > 0$, puisque $(D \cdot Z_t) = 0$ d'après ce lemme.
On peut donc trouver un $i$ tel que $a_i < m_i/d$ et
$(C_i \cdot Z_t) > 0$. Posons $Z_{t + 1} = Z_t + C_i$ et $i_{t + 1} = i$.
Considérons la suite exacte courte
$$
0 \to \mathcal{O}_X(-Z_t)|_{C_i} \to \mathcal{O}_{Z_{t + 1}} \to
\mathcal{O}_{Z_t} \to 0
$$
de Diviseurs, lemme \ref{divisors-lemma-ses-add-divisor}.
Par notre choix de $i$, on voit que
$\mathcal{O}_X(-Z_t)|_{C_i}$ est un faisceau inversible de degré négatif
sur la courbe propre $C_i$ ; il ne possède donc aucune section globale non nulle
(Variétés, lemme \ref{varieties-lemma-check-invertible-sheaf-trivial}).
On conclut que $H^0(\mathcal{O}_{Z_{t + 1}}) \subset H^0(\mathcal{O}_{Z_t})$
est un corps (cela est clair, mais résulte également de
Algèbre, lemme \ref{algebra-lemma-integral-under-field})
et une extension finie de $k$. Nous avons ainsi prolongé la suite.
Comme le procédé doit s'arrêter, par exemple parce que $t \leq \sum (m_i/d - 1)$,
ceci achève la démonstration.
\end{proof}

\begin{lemma}
\label{lemma-regular-model-genus}
\begin{reference}
\cite[Lemma 2.6]{Artin-Winters}
\end{reference}
Dans la situation \ref{situation-regular-model}, posons $d = \gcd(m_1, \ldots, m_n)$.
Soit $D = \sum (m_i/d) C_i$ considéré comme un diviseur de Cartier effectif sur $X$. Alors
$$
1 - g_C = d [\kappa : k] (1 - g_D)
$$
où $g_C$ est le genre de $C$, $g_D$ le genre de $D$, et
$\kappa = H^0(D, \mathcal{O}_D)$.
\end{lemma}

\begin{proof}
D'après le lemme \ref{lemma-regular-model-field}, $\kappa$ est un corps
et une extension finie de $k$. Comme on a aussi $H^0(C, \mathcal{O}_C) = K$,
on voit que les genres de $C$ et de $D$ sont définis (voir
Courbes algébriques, définition \ref{curves-definition-genus}) et que
$g_C = \dim_K H^1(C, \mathcal{O}_C)$ et
$g_D = \dim_\kappa H^1(D, \mathcal{O}_D)$.
D'après Catégories dérivées de schémas, lemme
\ref{perfect-lemma-chi-locally-constant-geometric},
on a
$$
1 - g_C = \chi(C, \mathcal{O}_C) =
\chi(X_k, \mathcal{O}_{X_k}) = \dim_k H^0(X_k, \mathcal{O}_{X_k})
- \dim_k H^1(X_k, \mathcal{O}_{X_k})
$$
Nous affirmons que
$$
\chi(X_k, \mathcal{O}_{X_k}) = d \chi(D, \mathcal{O}_D)
$$
Cela démontrera le lemme, car
$$
\chi(D, \mathcal{O}_D) =
\dim_k H^0(D, \mathcal{O}_D) - \dim_k H^1(D, \mathcal{O}_D) =
[\kappa : k](1 - g_D)
$$
Observons que $X_k = dD$ comme diviseur de Cartier effectif.
Pour démontrer l'affirmation, montrons par récurrence sur $1 \leq r \leq d$ que
$\chi(rD, \mathcal{O}_{rD}) = r \chi(D, \mathcal{O}_D)$.
Le cas initial $r = 1$ est trivial. Si $1 \leq r < d$, considérons
la suite exacte courte
$$
0 \to \mathcal{O}_X(rD)|_D \to \mathcal{O}_{(r + 1)D} \to
\mathcal{O}_{rD} \to 0
$$
de Diviseurs, lemme \ref{divisors-lemma-ses-add-divisor}. Par additivité
des caractéristiques d'Euler
(Variétés, lemme \ref{varieties-lemma-euler-characteristic-additive}),
il suffit de démontrer que
$\chi(D, \mathcal{O}_X(rD)|_D) = \chi(D, \mathcal{O}_D)$.
C'est vrai parce que $\mathcal{O}_X(rD)|_D$ est un élément de torsion
de $\Pic(D)$ (lemme \ref{lemma-multiple-fibre-normal-bundle})
et parce que le degré d'un fibré en droites est additif
(Variétés, lemme \ref{varieties-lemma-degree-tensor-product}),
donc nul pour les faisceaux inversibles de torsion.
\end{proof}

\begin{lemma}
\label{lemma-exceptional-curves-dont-meet}
Dans la situation \ref{situation-regular-model}, soient deux indices $i, j$
tels que $C_i$ et $C_j$ soient des courbes exceptionnelles de première espèce
et que $C_i \cap C_j \not = \emptyset$. Alors
$n = 2$, $m_1 = m_2 = 1$, $C_1 \cong \mathbf{P}^1_k$,
$C_2 \cong \mathbf{P}^1_k$, $C_1$ et $C_2$ se rencontrent en un point $k$-rationnel,
et $C$ est de genre $0$.
\end{lemma}

\begin{proof}
Choisissons des isomorphismes $C_i = \mathbf{P}^1_{\kappa_i}$ et
$C_j = \mathbf{P}^1_{\kappa_j}$. Le schéma $C_i \cap C_j$
est un diviseur de Cartier effectif non vide tant dans $C_i$ que dans $C_j$.
Ainsi,
$$
(C_i \cdot C_j) = \deg(C_i \cap C_j) \geq \max([\kappa_i: k], [\kappa_j : k])
$$
La première égalité a été démontrée dans la preuve du
lemme \ref{lemma-intersection-pairing}.
D'autre part, l'auto-intersection $(C_i \cdot C_i)$ est égale
au degré de $\mathcal{O}_X(C_i)$ sur $C_i$, qui vaut $-[\kappa_i : k]$
puisque $C_i$ est une courbe exceptionnelle de première espèce. De même pour
$C_j$. D'après le lemme \ref{lemma-properties-form},
$$
0 \geq (C_i + C_j)^2 = -[\kappa_i : k] + 2(C_i \cdot C_j) - [\kappa_j : k]
$$
Ceci implique que $[\kappa_i : k] = \deg(C_i \cap C_j) = [\kappa_j : k]$
et que $(C_i + C_j)^2 = 0$. En appliquant de nouveau le lemme,
on conclut que $n = 2$, $\{1, 2\} = \{i, j\}$ et $m_1 = m_2$.
De plus, l'intersection schématique $C_i \cap C_j$ consiste en
un unique point $p$ de corps résiduel $\kappa$, et
$\kappa_i \to \kappa \leftarrow \kappa_j$ sont des isomorphismes.
Soit $D = C_1 + C_2$ considéré comme un diviseur de Cartier effectif sur $X$.
Observons que $D$ est la réunion schématique de $C_1$ et $C_2$
(Diviseurs, lemme \ref{divisors-lemma-sum-effective-Cartier-divisors-union}) ;
on a donc une suite exacte courte
$$
0 \to \mathcal{O}_D \to \mathcal{O}_{C_1} \oplus \mathcal{O}_{C_2} \to
\mathcal{O}_p \to 0
$$
d'après Morphismes, lemme \ref{morphisms-lemma-scheme-theoretic-union}.
Puisque nous connaissons la cohomologie de $C_i \cong \mathbf{P}^1_\kappa$
(Cohomologie des schémas, lemme
\ref{coherent-lemma-cohomology-projective-space-over-ring}),
la suite exacte longue de cohomologie donne
$H^0(D, \mathcal{O}_D) = \kappa$ et
$H^1(D, \mathcal{O}_D) = 0$. D'après le lemme \ref{lemma-regular-model-genus},
on conclut que
$$
1 - g_C = d[\kappa : k](1 - 0)
$$
où $d = m_1 = m_2$. Il s'ensuit que $g_C = 0$, $d = m_1 = m_2 = 1$
et $\kappa = k$.
\end{proof}





\section{Unicité du modèle minimal}
\label{section-uniqueness}

\noindent
Si le genre de la fibre générique est positif, les modèles minimaux sont uniques
(lemme \ref{lemma-minimal-model-unique}) et possèdent par conséquent une propriété
universelle convenable (lemme \ref{lemma-minimal-model-mapping-property}).

\begin{lemma}
\label{lemma-minimal-model-unique}
Soit $C$ une courbe projective lisse sur $K$ telle que
$H^0(C, \mathcal{O}_C) = K$ et de genre $> 0$.
Il existe un unique modèle minimal de $C$.
\end{lemma}

\begin{proof}
Nous avons déjà démontré la partie difficile du lemme, à savoir l'existence
d'un modèle minimal (dont la preuve repose sur la
résolution des singularités de surfaces) ; voir la
proposition \ref{proposition-exists-minimal-model}.
Pour démontrer l'unicité, supposons que $X$ et $Y$ soient deux
modèles minimaux. D'après
Résolution des surfaces, lemme \ref{resolve-lemma-birational-regular-surfaces},
il existe un diagramme de $S$-morphismes
$$
X = X_0 \leftarrow X_1 \leftarrow \ldots \leftarrow X_n = Y_m
\to \ldots \to Y_1 \to Y_0 = Y
$$
où chaque morphisme est un éclatement en un point fermé. La
fibre exceptionnelle du morphisme $X_n \to X_{n - 1}$ est une
courbe exceptionnelle de première espèce $E$. Nous affirmons que $E$ est
contractée en un point par le morphisme $X_n = Y_m \to Y$.
Si tel est le cas, alors $X_n \to Y$ se factorise par $X_{n - 1}$ d'après
Résolution des surfaces, lemme \ref{resolve-lemma-factor-through-contraction}.
Dans ce cas, le morphisme $X_{n - 1} \to Y$ est encore une suite de
contractions de courbes exceptionnelles d'après
Résolution des surfaces, lemme
\ref{resolve-lemma-proper-birational-regular-surfaces}.
On conclut donc par récurrence sur $n$. (Le cas initial $n = 0$ signifie
qu'il existe une suite de contractions
$X = Y_m \to \ldots \to Y_1 \to Y_0 = Y$
aboutissant à $Y$. Toutefois, comme $X$ est un modèle minimal, il ne contient
aucune courbe exceptionnelle de première espèce ; ainsi $m = 0$ et $X = Y$.)

\medskip\noindent
Démonstration de l'affirmation. Nous allons montrer par récurrence sur $m$ que toute courbe
exceptionnelle de première espèce $E \subset Y_m$ est envoyée sur un point
par le morphisme $Y_m \to Y$. Si $m = 0$, c'est clair puisque
$Y$ est un modèle minimal. Si $m > 0$, alors soit
$Y_m \to Y_{m - 1}$ contracte $E$ (et nous avons terminé), soit
la fibre exceptionnelle $E' \subset Y_m$ de $Y_m \to Y_{m - 1}$
est une seconde courbe exceptionnelle de première espèce.
Puisque $E$ et $E'$ sont toutes deux des composantes irréductibles de la fibre
spéciale, et puisque $g_C > 0$ par hypothèse, on conclut que
$E \cap E' = \emptyset$ d'après le
lemme \ref{lemma-exceptional-curves-dont-meet}.
L'image de $E$ dans $Y_{m - 1}$ est alors une courbe
exceptionnelle de première espèce (cela est clair, car le morphisme
$Y_m \to Y_{m - 1}$ est un isomorphisme au voisinage de $E$).
Par récurrence, on voit que $Y_{m - 1} \to Y$ contracte cette courbe,
ce qui achève la démonstration.
\end{proof}

\begin{lemma}
\label{lemma-minimal-model-mapping-property}
Soit $C$ une courbe projective lisse sur $K$ telle que $H^0(C, \mathcal{O}_C) = K$
et de genre $> 0$. Soit $X$ le modèle minimal de $C$
(lemme \ref{lemma-minimal-model-unique}).
Soit $Y$ un modèle propre régulier de $C$. Il existe alors un unique
morphisme de modèles $Y \to X$ qui soit une suite de contractions de
courbes exceptionnelles de première espèce.
\end{lemma}

\begin{proof}
L'existence et les propriétés du morphisme $X \to Y$
résultent immédiatement du lemme \ref{lemma-pre-exists-minimal-model}
et de l'unicité du modèle minimal.
Le morphisme $Y \to X$ est unique parce que
$C \subset Y$ est schématiquement dense et que $X$ est séparé
(voir Morphismes, lemme \ref{morphisms-lemma-equality-of-morphisms}).
\end{proof}

\begin{example}
\label{example-nonunique-in-genus-zero}
Si le genre de $C$ est $0$, les modèles minimaux ne sont effectivement pas uniques.
Considérons en effet le sous-schéma fermé
$$
X \subset \mathbf{P}^2_R
$$
défini par $T_1T_2 - \pi T_0^2 = 0$. Plus précisément, $X$ est défini
comme $\text{Proj}(R[T_0, T_1, T_2]/(T_1T_2 - \pi T_0^2))$. La
fibre spéciale $X_k$ est alors la réunion de deux courbes exceptionnelles $C_1$, $C_2$, toutes deux
isomorphes à $\mathbf{P}^1_k$
(exactement comme dans le lemme \ref{lemma-exceptional-curves-dont-meet}).
La projection depuis $(0 : 1 : 0)$ définit un morphisme $X \to \mathbf{P}^1_R$
qui contracte $C_2$ et induit un isomorphisme de $C_1$ avec la fibre spéciale
de $\mathbf{P}^1_R$. La projection depuis $(0 : 0 : 1)$ définit un
morphisme $X \to \mathbf{P}^1_R$ qui contracte $C_1$ et induit un
isomorphisme de $C_2$ avec la fibre spéciale de $\mathbf{P}^1_R$.
Plus précisément, ces morphismes correspondent aux homomorphismes de $R$-algèbres
graduées
$$
R[T_0, T_1] \longrightarrow
R[T_0, T_1, T_2]/(T_1T_2 - \pi T_0^2) \longleftarrow
R[T_0, T_2]
$$
Nous étudierons ce phénomène dans le lemme \ref{lemma-nonuniqueness}.
\end{example}









\section{Une formule pour le genre}
\label{section-genus-formula}

\noindent
La structure combinatoire issue d'un modèle propre régulier
est soumise à une restriction supplémentaire.

\begin{lemma}
\label{lemma-add-component}
Dans la situation \ref{situation-regular-model}, supposons donnés des
diviseurs de Cartier effectifs $D, D' \subset X$ tels que
$D' = D + C_i$ pour un certain $i \in \{1, \ldots, n\}$ et $D' \subset X_k$.
Alors
$$
\chi(X_k, \mathcal{O}_{D'}) - \chi(X_k, \mathcal{O}_D) =
\chi(X_k, \mathcal{O}_X(-D)|_{C_i}) =
-(D \cdot C_i) + \chi(C_i, \mathcal{O}_{C_i})
$$
\end{lemma}

\begin{proof}
La seconde égalité résulte de la définition de la forme bilinéaire
$(\ \cdot\ )$ dans (\ref{equation-form}) et du
lemme \ref{lemma-intersection-pairing}. Pour voir la première
égalité, distinguons deux cas.
En effet, si $C_i \not \subset D$, alors $D'$ est la réunion
schématique de $D$ et $C_i$ (d'après
Diviseurs, lemme \ref{divisors-lemma-sum-effective-Cartier-divisors-union}),
et l'on obtient une suite exacte courte
$$
0 \to \mathcal{O}_{D'} \to
\mathcal{O}_D \times \mathcal{O}_{C_i} \to
\mathcal{O}_{D \cap C_i} \to 0
$$
d'après Morphismes, lemme \ref{morphisms-lemma-scheme-theoretic-union}.
Comme on a aussi une suite exacte
$$
0 \to \mathcal{O}_X(-D)|_{C_i} \to
\mathcal{O}_{C_i} \to \mathcal{O}_{D \cap C_i} \to 0
$$
(Diviseurs, remarque \ref{divisors-remark-ses-regular-section}),
on conclut que l'affirmation est vraie
par additivité des caractéristiques d'Euler
(Variétés, lemme \ref{varieties-lemma-euler-characteristic-additive}).
D'autre part, si $C_i \subset D$, alors on obtient une
suite exacte
$$
0 \to \mathcal{O}_X(-D)|_{C_i} \to \mathcal{O}_{D'} \to \mathcal{O}_D \to 0
$$
d'après Diviseurs, lemme \ref{divisors-lemma-ses-add-divisor},
et l'on voit immédiatement que le lemme est vrai.
\end{proof}

\begin{lemma}
\label{lemma-genus-formula}
Dans la situation \ref{situation-regular-model}, on a
$$
g_C = 1 + \sum\nolimits_{i = 1, \ldots, n}
m_i\left([\kappa_i : k] (g_i - 1) - \frac{1}{2}(C_i \cdot C_i)\right)
$$
où $\kappa_i = H^0(C_i, \mathcal{O}_{C_i})$,
$g_i$ est le genre de $C_i$, et $g_C$ le genre de $C$.
\end{lemma}

\begin{proof}
Notre outil fondamental sera Catégories dérivées de schémas, lemme
\ref{perfect-lemma-chi-locally-constant-geometric},
qui montre que
$$
1 - g_C = \chi(C, \mathcal{O}_C) =
\chi(X_k, \mathcal{O}_{X_k})
$$
Choisissons une suite de diviseurs de Cartier effectifs
$$
X_k = D_m \supset D_{m - 1} \supset \ldots \supset D_1 \supset D_0 = \emptyset
$$
telle que $D_{j + 1} = D_j + C_{i_j}$ pour tout $j$. (Il est clair que
l'on peut choisir une telle suite en diminuant d'une unité, à chaque étape, une multiplicité
non nulle de $D_{j + 1}$.) En appliquant le lemme \ref{lemma-add-component},
en partant de $\chi(\mathcal{O}_{D_0}) = 0$, on obtient
\begin{align*}
1 - g_C
& =
\chi(X_k, \mathcal{O}_{X_k}) \\
& =
\sum\nolimits_j
\left(-(D_j \cdot C_{i_j}) +  \chi(C_{i_j}, \mathcal{O}_{C_{i_j}})\right) \\
& =
- \sum\nolimits_j
(C_{i_1} + C_{i_2} + \ldots + C_{i_{j - 1}} \cdot C_{i_j}) +
\sum\nolimits_j \chi(C_{i_j}, \mathcal{O}_{C_{i_j}}) \\
& =
-\frac{1}{2}\sum\nolimits_{j \not = j'} (C_{i_{j'}} \cdot C_{i_j}) +
\sum m_i \chi(C_i, \mathcal{O}_{C_i}) \\
& =
\frac{1}{2} \sum m_i(C_i \cdot C_i) + \sum m_i \chi(C_i, \mathcal{O}_{C_i})
\end{align*}
La dernière égalité mérite peut-être une explication. En effet, puisque
$\sum_j C_{i_j} = \sum m_i C_i$, on a
$(\sum_j C_{i_j} \cdot \sum_j C_{i_j}) = 0$ d'après le
lemme \ref{lemma-properties-form}. On voit donc que
$$
0 = \sum\nolimits_{j \not = j'} (C_{i_{j'}} \cdot C_{i_j}) +
\sum m_i(C_i \cdot C_i)
$$
en décomposant ce produit en termes « non diagonaux » et « diagonaux ».
Notons que $\kappa_i$ est un corps fini sur $k$ d'après
Variétés, lemme \ref{varieties-lemma-regular-functions-proper-variety}.
Le genre de $C_i$ est donc défini et l'on a
$\chi(C_i, \mathcal{O}_{C_i}) = [\kappa_i : k](1 - g_i)$.
En réunissant le tout et en réarrangeant les termes, on obtient
$$
g_C = - \frac{1}{2}\sum m_i(C_i \cdot C_i) +
\sum m_i[\kappa_i : k](g_i - 1) + 1
$$
ce qui est aussi l'assertion du lemme.
\end{proof}

\begin{lemma}
\label{lemma-numerical-type-of-model}
Dans la situation \ref{situation-regular-model}, avec
$\kappa_i = H^0(C_i, \mathcal{O}_{C_i})$ et $g_i$ le genre de $C_i$,
les données
$$
n, m_i, (C_i \cdot C_j), [\kappa_i : k], g_i
$$
forment un type numérique de genre égal au genre de $C$.
\end{lemma}

\begin{proof}
(Dans la démonstration du lemme \ref{lemma-genus-formula},
nous avons vu que les quantités
utilisées dans l'énoncé du lemme sont bien définies.)
Nous devons vérifier les conditions (1) -- (5) de la
définition \ref{definition-type}.

\medskip\noindent
La condition (1) est immédiate.

\medskip\noindent
Condition (2). La symétrie de la matrice $(C_i \cdot C_j)$ résulte de
l'équation (\ref{equation-form}) et du
lemme \ref{lemma-intersection-pairing}.
La positivité de $(C_i \cdot C_j)$ pour $i \not = j$
est la partie (1) du lemme \ref{lemma-properties-form}.

\medskip\noindent
La condition (3) est la partie (3) du lemme \ref{lemma-properties-form}.

\medskip\noindent
La condition (4) est la partie (2) du lemme \ref{lemma-properties-form}.

\medskip\noindent
La condition (5) résulte du fait que $(C_i \cdot C_j)$ est
le degré d'un module inversible sur $C_i$, lequel est divisible
par $[\kappa_i : k]$ ; voir Variétés, lemme \ref{varieties-lemma-divisible}.

\medskip\noindent
La formule du genre démontrée dans le lemme \ref{lemma-genus-formula}
nous dit que le type numérique possède le genre annoncé ; voir la
définition \ref{definition-genus}.
\end{proof}

\begin{definition}
\label{definition-numerical-type-model}
Dans la situation \ref{situation-regular-model}, le
{\it type numérique associé à $X$} est le type numérique
décrit dans le lemme \ref{lemma-numerical-type-of-model}.
\end{definition}

\noindent
Nous allons maintenant faire correspondre la minimalité du modèle à celle du type.

\begin{lemma}
\label{lemma-numerical-type-minimal-model}
Dans la situation \ref{situation-regular-model}, les assertions suivantes
sont équivalentes :
\begin{enumerate}
\item $X$ est un modèle minimal, et
\item le type numérique associé à $X$ est minimal.
\end{enumerate}
\end{lemma}

\begin{proof}
Si le type numérique est minimal, il n'existe aucun $i$ tel que
$g_i = 0$ et $(C_i \cdot C_i) = -[\kappa_i: k]$ ; voir la
définition \ref{definition-type-minimal}.
Ceci implique assurément qu'aucune des courbes $C_i$
n'est une courbe exceptionnelle de première espèce.

\medskip\noindent
Réciproquement, supposons que le type numérique ne soit pas minimal.
Il existe alors un $i$ tel que $g_i = 0$ et
$(C_i \cdot C_i) = -[\kappa_i: k]$.
Nous affirmons que ceci implique que $C_i$ est une courbe exceptionnelle
de première espèce. En effet, le faisceau inversible
$\mathcal{O}_X(-C_i)|_{C_i}$ est de degré $-(C_i \cdot C_i) = [\kappa_i : k]$
lorsque $C_i$ est vue comme une courbe propre sur $k$, et est donc
de degré $1$ lorsque $C_i$ est vue comme une courbe propre sur $\kappa_i$.
En appliquant
Courbes algébriques, proposition \ref{curves-proposition-projective-line},
on conclut que $C_i \cong \mathbf{P}^1_{\kappa_i}$ comme schémas
sur $\kappa_i$. Comme le groupe de Picard de $\mathbf{P}^1$
sur un corps est $\mathbf{Z}$, on voit que le faisceau normal
de $C_i$ dans $X$ est isomorphe à $\mathcal{O}_{\mathbf{P}_{\kappa_i}}(-1)$,
ce qui achève la démonstration.
\end{proof}

\begin{remark}
\label{remark-numerical-type-not-from-model}
Tout type numérique ne provient pas d'un modèle, pour la raison triviale
qu'il existe des types numériques dont le genre est négatif.
Il existe des types numériques minimaux de genre positif qui
ne sont pas le type numérique associé à un modèle (sur un anneau de valuation discrète)
d'une courbe projective lisse géométriquement irréductible (sur le
corps des fractions de cet anneau). Un exemple simple est
$n = 1$, $m_1 = 1$, $a_{11} = 0$, $w_1 = 6$, $g_1 = 1$.
En effet, dans ce cas, la fibre spéciale $X_k$ ne serait pas
géométriquement connexe, car elle serait définie sur une extension
$\kappa$ de $k$ de degré $6$. Cela contredit le
fait que la fibre générique est géométriquement connexe (voir
Compléments sur les morphismes, lemme
\ref{more-morphisms-lemma-geometrically-connected-fibres-towards-normal}).
De même, $n = 2$, $m_1 = m_2 = 1$, $-a_{11} = -a_{22} = a_{12} = a_{21} = 6$,
$w_1 = w_2 = 6$, $g_1 = g_2 = 1$ fournirait un exemple
pour la même raison (détails omis). Mais si le pgcd des
$w_i$ vaut $1$, nous ne disposons d'aucun exemple.
\end{remark}

\begin{lemma}
\label{lemma-numerical-type-rational-point}
Dans la situation \ref{situation-regular-model}, supposons que $C$ ait un point $K$-rationnel.
Alors
\begin{enumerate}
\item $X_k$ possède un point $k$-rationnel $x$ qui est un point lisse de $X_k$
sur $k$,
\item si $x \in C_i$, alors $H^0(C_i, \mathcal{O}_{C_i}) = k$ et
$m_i = 1$, et
\item $H^0(X_k, \mathcal{O}_{X_k}) = k$, et $X_k$ est de genre égal au
genre de $C$.
\end{enumerate}
\end{lemma}

\begin{proof}
Puisque $X \to \Spec(R)$ est propre, le point $K$-rationnel se prolonge en
un morphisme $a : \Spec(R) \to X$ par le critère valuatif de propreté
(Morphismes, lemme \ref{morphisms-lemma-characterize-proper}).
Soit $x \in X$ l'image par $a$ du point fermé de $\Spec(R)$.
Alors $a$ correspond à un homomorphisme de $R$-algèbres
$\psi : \mathcal{O}_{X, x} \to R$
(voir Schémas, section \ref{schemes-section-points}).
Il s'ensuit que $\pi \not \in \mathfrak m_x^2$ (puisque l'image
de $\pi$ dans $R$ n'appartient pas à $\mathfrak m_R^2$).
Par conséquent, $\mathcal{O}_{X_k, x} = \mathcal{O}_{X, x}/\pi \mathcal{O}_{X, x}$
est régulier (Algèbre, lemme \ref{algebra-lemma-regular-ring-CM}).
Alors $X_k \to \Spec(k)$ est lisse en $x$ d'après
Algèbre, lemme \ref{algebra-lemma-separable-smooth}.
Il s'ensuit que $x$ appartient à une unique composante irréductible
$C_i$ de $X_k$, que $\mathcal{O}_{C_i, x} = \mathcal{O}_{X_k, x}$
et que $m_i = 1$. Le fait que $C_i$ possède un
point $k$-rationnel implique que le corps
$\kappa_i = H^0(C_i, \mathcal{O}_{C_i})$
(Variétés, lemme \ref{varieties-lemma-regular-functions-proper-variety})
est égal à $k$. Ceci démontre (1). On a
$H^0(X_k, \mathcal{O}_{X_k}) = k$
parce que $H^0(X_k, \mathcal{O}_{X_k})$ est une extension de corps
de $k$ (lemme \ref{lemma-regular-model-field})
qui s'envoie dans $H^0(C_i, \mathcal{O}_{C_i}) = k$.
L'égalité des genres résulte du lemme \ref{lemma-regular-model-genus}.
\end{proof}

\begin{lemma}
\label{lemma-genus-reduction-smaller}
Dans la situation \ref{situation-regular-model}, supposons que $X$ soit un modèle minimal,
que $\gcd(m_1, \ldots, m_n) = 1$ et que $H^0((X_k)_{red}, \mathcal{O}) = k$. Alors
l'application
$$
H^1(X_k, \mathcal{O}_{X_k}) \to H^1((X_k)_{red}, \mathcal{O}_{(X_k)_{red}})
$$
est surjective et possède un noyau non trivial dès que $(X_k)_{red} \not = X_k$.
\end{lemma}

\begin{proof}
Par annulation de la cohomologie en degrés $\geq 2$ sur $X_k$
(Cohomologie, proposition \ref{cohomology-proposition-vanishing-Noetherian}),
toute surjection de faisceaux abéliens sur $X_k$ induit une surjection sur $H^1$.
Considérons la suite
$$
(X_k)_{red} = Z_0 \subset Z_1 \subset \ldots \subset Z_m = X_k
$$
du lemme \ref{lemma-regular-model-field}. Puisque les homomorphismes de corps
$H^0(Z_j, \mathcal{O}_{Z_j}) \to
H^0((X_k)_{red}, \mathcal{O}_{(X_k)_{red}}) = k$
sont injectifs, on conclut que $H^0(Z_j, \mathcal{O}_{Z_j}) = k$ pour
$j = 0, \ldots, m$. Il s'ensuit que
$H^0(X_k, \mathcal{O}_{X_k}) \to H^0(Z_{m - 1}, \mathcal{O}_{Z_{m - 1}})$
est surjectif. Soit $C = C_{i_m}$. Alors $X_k = Z_{m - 1} + C$.
Soit $\mathcal{L} = \mathcal{O}_X(-Z_{m - 1})|_C$.
Alors $\mathcal{L}$ est un $\mathcal{O}_C$-module inversible.
Comme dans la démonstration du lemme \ref{lemma-regular-model-field},
il existe une suite exacte
$$
0 \to \mathcal{L} \to \mathcal{O}_{X_k} \to \mathcal{O}_{Z_{m - 1}} \to 0
$$
de faisceaux cohérents sur $X_k$. On en déduit
une suite exacte courte
$$
0 \to
H^1(C, \mathcal{L}) \to H^1(X_k, \mathcal{O}_{X_k}) \to
H^1(Z_{m - 1}, \mathcal{O}_{Z_{m - 1}}) \to 0
$$
Le degré de $\mathcal{L}$ sur $C$ relativement à $k$ est
$$
(C \cdot -Z_{m - 1}) = (C \cdot C - X_k) = (C \cdot C)
$$
Posons $\kappa = H^0(C, \mathcal{O}_C)$ et $w = [\kappa : k]$.
Par définition du degré d'un faisceau inversible, on voit que
$$
\chi(C, \mathcal{L}) =
\chi(C, \mathcal{O}_C) + (C \cdot C) =
w(1 - g_C) + (C \cdot C)
$$
où $g_C$ est le genre de $C$. Cette expression est $< 0$ puisque $X$ est
minimal et que, par conséquent, $C$ n'est pas une courbe exceptionnelle de
première espèce (voir la preuve du lemme
\ref{lemma-numerical-type-minimal-model}). Ainsi
$\dim_k H^1(C, \mathcal{L}) > 0$, ce qui achève la preuve.
\end{proof}

\begin{lemma}
\label{lemma-genus-reduction-bigger-than}
Dans la situation \ref{situation-regular-model}, supposons que $X_k$ possède
un point $k$-rationnel $x$ qui est un point lisse de
$X_k \to \Spec(k)$. Alors
$$
\dim_k H^1((X_k)_{red}, \mathcal{O}_{(X_k)_{red}}) \geq
g_{top} + g_{geom}(X_k/k)
$$
où $g_{geom}$ est défini dans Courbes algébriques, section
\ref{curves-section-genus-geometric-genus}, et où $g_{top}$ est le genre
topologique (définition \ref{definition-top-genus}) du type numérique
associé à $X_k$ (définition \ref{definition-numerical-type-model}).
\end{lemma}

\begin{proof}
Nous allons démontrer l'inégalité
$$
\dim_k H^1(D, \mathcal{O}_D) \geq g_{top}(D) + g_{geom}(D/k)
$$
pour tout diviseur de Cartier effectif réduit connexe
$D \subset (X_k)_{red}$ contenant $x$, par récurrence sur le nombre de
composantes irréductibles de $D$.
Ici $g_{top}(D) = 1 - m + e$, où $m$ est le nombre de composantes
irréductibles de $D$ et $e$ le nombre de paires non ordonnées de composantes
de $D$ qui se rencontrent.

\medskip\noindent
Cas initial : $D$ possède une seule composante irréductible. Alors
$D = C_i$ est l'unique composante irréductible contenant $x$.
Dans ce cas, $\dim_k H^1(D, \mathcal{O}_D) = g_i$ et
$g_{top}(D) = 0$. Puisque $C_i$ possède un point lisse $k$-rationnel,
elle est géométriquement intègre (Variétés, lemme
\ref{varieties-lemma-variety-with-smooth-rational-point}). Il s'ensuit que
$g_i$ est le genre de $C_{i, \overline{k}}$ (Courbes algébriques, lemme
\ref{curves-lemma-genus-base-change}). Il s'ensuit aussi que
$g_{geom}(D/k)$ est le genre de la normalisation
$C_{i, \overline{k}}^\nu$ de $C_{i, \overline{k}}$. En appliquant Courbes
algébriques, lemme \ref{curves-lemma-genus-normalization}, au morphisme de
normalisation $C_{i, \overline{k}}^\nu \to C_{i, \overline{k}}$, on obtient
\begin{equation}
\label{equation-genus-change-special-component}
\text{genre de }C_{i, \overline{k}} \geq
\text{genre de }C_{i, \overline{k}}^\nu
\end{equation}
En combinant ce qui précède, on conclut que
$\dim_k H^1(D, \mathcal{O}_D) \geq g_{top}(D) + g_{geom}(D/k)$
dans ce cas.

\medskip\noindent
Étape de récurrence. Supposons que $D$ possède plus de $1$ composante
irréductible. On peut alors écrire $D = C_i + D'$, avec $x \in D'$ et
$D'$ encore connexe. C'est un exercice de théorie des graphes que nous
laissons au lecteur (indication : prendre pour $C_i$ la composante de $D$
la plus éloignée de $x$). Calculons la variation des invariants.
Comme $x \in D'$, on a
$H^0(D, \mathcal{O}_D) = H^0(D', \mathcal{O}_{D'}) = k$.
En considérant la suite exacte courte de faisceaux
$$
0 \to \mathcal{O}_D \to \mathcal{O}_{C_i} \oplus \mathcal{O}_{D'}
\to \mathcal{O}_{C_i \cap D'} \to 0
$$
(Morphismes, lemme \ref{morphisms-lemma-scheme-theoretic-union}) et en
utilisant l'additivité des caractéristiques d'Euler, on trouve
\begin{align*}
\dim_k H^1(D, \mathcal{O}_D) - \dim_k H^1(D', \mathcal{O}_{D'})
& =
-\chi(\mathcal{O}_{C_i}) + \chi(\mathcal{O}_{C_i \cap D'}) \\
& =
w_i(g_i - 1) + \sum\nolimits_{C_j \subset D'} a_{ij}
\end{align*}
Comme dans le lemme \ref{lemma-numerical-type-of-model}, on pose ici
$w_i = [\kappa_i : k]$, $\kappa_i = H^0(C_i, \mathcal{O}_{C_i})$ ;
$g_i$ est le genre de $C_i$, et $a_{ij} = (C_i \cdot C_j)$.
On a
$$
g_{top}(D) - g_{top}(D') = -1 +
\sum\nolimits_{C_j \subset D'\text{ rencontrant }C_i} 1
$$
On a
$$
g_{geom}(D/k) - g_{geom}(D'/k) = g_{geom}(C_i/k)
$$
d'après Courbes algébriques, lemme
\ref{curves-lemma-bound-geometric-genus}. En combinant ces égalités avec
l'hypothèse de récurrence, on conclut qu'il suffit de montrer que
$$
w_i g_i - g_{geom}(C_i/k) +
\sum\nolimits_{C_j \subset D'\text{ rencontrant } C_i} (a_{ij} - 1) - (w_i - 1)
$$
est positif ou nul. En effet, on a
\begin{equation}
\label{equation-genus-change}
w_i g_i \geq [\kappa_i : k]_s g_i \geq g_{geom}(C_i/k)
\end{equation}
La seconde inégalité résulte de Courbes algébriques, lemme
\ref{curves-lemma-bound-geometric-genus-curve}. D'autre part, puisque
$w_i$ divise $a_{ij}$ (Variétés, lemme
\ref{varieties-lemma-divisible}), il est clair que
\begin{equation}
\label{equation-change-intersections}
\sum\nolimits_{C_j \subset D'\text{ rencontrant } C_i} (a_{ij} - 1) - (w_i - 1)
\geq 0
\end{equation}
car il existe au moins une composante $C_j \subset D'$ rencontrant $C_i$.
\end{proof}

\begin{lemma}
\label{lemma-equality-genus-reduction-bigger-than}
Si l'égalité a lieu dans le lemme
\ref{lemma-genus-reduction-bigger-than}, alors
\begin{enumerate}
\item l'unique composante irréductible de $X_k$ contenant $x$ est une
courbe projective lisse géométriquement irréductible sur $k$ ;
\item si $C \subset X_k$ est une autre composante irréductible, alors
$\kappa = H^0(C, \mathcal{O}_C)$ est une extension finie séparable de $k$,
$C$ possède un point $\kappa$-rationnel et $C$ est lisse sur $\kappa$.
\end{enumerate}
\end{lemma}

\begin{proof}
En reprenant la preuve du lemme
\ref{lemma-genus-reduction-bigger-than}, on voit que, pour obtenir
l'égalité, les inégalités
(\ref{equation-genus-change-special-component}),
(\ref{equation-genus-change}) et
(\ref{equation-change-intersections})
doivent être des égalités.

\medskip\noindent
Soit $C_i$ la composante irréductible contenant $x$.
L'égalité dans (\ref{equation-genus-change-special-component}) montre, par
Courbes algébriques, lemme \ref{curves-lemma-genus-normalization}, que
$C_{i, \overline{k}}^\nu \to C_{i, \overline{k}}$ est un isomorphisme.
Par conséquent, $C_{i, \overline{k}}$ est lisse et l'assertion (1) est vraie.

\medskip\noindent
Soit ensuite $C_i \subset X_k$ une autre composante irréductible.
On peut alors supposer que $D = D' + C_i$ comme dans l'étape de récurrence
de la preuve du lemme \ref{lemma-genus-reduction-bigger-than}.
L'égalité dans (\ref{equation-genus-change}) implique immédiatement que
$\kappa_i/k$ est finie séparable. L'égalité dans
(\ref{equation-change-intersections}) implique soit que $a_{ij} = 1$ pour
un certain $j$, soit qu'il existe une unique composante
$C_j \subset D'$ rencontrant $C_i$ et que $a_{ij} = w_i$.
Dans les deux cas, on trouve que $C_i$ possède un point
$\kappa_i$-rationnel $c$ et que $c = C_i \cap C_j$ au sens schématique.
Puisque $\mathcal{O}_{X,c}$ est un anneau local régulier, cela implique que
les équations locales de $C_i$ et de $C_j$ forment un système régulier de
paramètres de l'anneau local $\mathcal{O}_{X,c}$. Alors
$\mathcal{O}_{C_i,c}$ est régulier d'après Algèbre, lemme
\ref{algebra-lemma-regular-ring-CM}. On conclut que
$C_i \to \Spec(\kappa_i)$ est lisse en $c$ (Algèbre, lemme
\ref{algebra-lemma-separable-smooth}). Il s'ensuit que $C_i$ est
géométriquement intègre sur $\kappa_i$ (Variétés, lemme
\ref{varieties-lemma-variety-with-smooth-rational-point}). Pour conclure,
il reste à montrer que $C_i$ est lisse sur $\kappa_i$. Observons que
$$
C_{i, \overline{k}} = C_i \times_{\Spec(k)} \Spec(\overline{k})
= \coprod\nolimits_{\kappa_i \to \overline{k}}
C_i \times_{\Spec(\kappa_i)} \Spec(\overline{k})
$$
où il y a $[\kappa_i : k]$ composantes. Ainsi, si $C_i$ n'est pas lisse
sur $\kappa_i$, chacune de ces courbes n'est pas lisse ; elles ne sont donc
pas normales, et le morphisme de normalisation fait baisser le genre
(Courbes algébriques, lemme \ref{curves-lemma-genus-normalization}). Cela
est impossible, car il ferait baisser le genre géométrique de $C_i/k$, en
contradiction avec $[\kappa_i : k]g_i = g_{geom}(C_i/k)$.
\end{proof}








\section{Contraction des courbes exceptionnelles}
\label{section-contract}

\noindent
Le lemme suivant décrit la variation des nombres d'intersection lorsque
l'on contracte une courbe exceptionnelle de première espèce dans un modèle
propre régulier. Nous le plaçons ici surtout pour le comparer aux
contractions numériques introduites dans le lemme \ref{lemma-contract}.
Nous comparerons les contractions géométriques et numériques dans la
remarque \ref{remark-compare-contractions}.

\begin{lemma}
\label{lemma-blowdown-regular-model}
Dans la situation \ref{situation-regular-model}, supposons que $C_n$ soit
une courbe exceptionnelle de première espèce. Soit $f : X \to X'$ la
contraction de $C_n$. Posons $C'_i = f(C_i)$ et écrivons
$X'_k = \sum m'_iC'_i$. Alors $X'$, $C'_i$, $i = 1, \ldots, n' = n - 1$,
et $m'_i = m_i$ sont comme dans la situation
\ref{situation-regular-model}, et l'on a
\begin{enumerate}
\item pour $i,j<n$,
$(C'_i \cdot C'_j) =
(C_i \cdot C_j) - (C_i \cdot C_n)(C_j \cdot C_n)/(C_n \cdot C_n)$ ;
\item pour $i<n$, si $C_i \cap C_n \not = \emptyset$, il existe des
applications $\kappa_i \leftarrow \kappa'_i \rightarrow \kappa_n$.
\end{enumerate}
Ici $\kappa_i = H^0(C_i, \mathcal{O}_{C_i})$ et
$\kappa'_i = H^0(C'_i, \mathcal{O}_{C'_i})$.
\end{lemma}

\begin{proof}
D'après Résolution des surfaces, lemme
\ref{resolve-lemma-contract-ample}, on peut contracter $C_n$ par un
morphisme $f : X \to X'$ tel que $X'$ soit régulier et projectif sur $R$.
On voit donc que $X'$ est comme dans la situation
\ref{situation-regular-model}. Soit $x \in X'$ l'image de $C_n$.
Puisque $f$ induit un isomorphisme
$X \setminus C_n \to X' \setminus \{x\}$, il est clair que
$m'_i = m_i$ pour $i<n$.

\medskip\noindent
L'assertion (2) du lemme résulte immédiatement de l'existence des
morphismes $C_i \to C'_i$ et $C_n \to x \to C'_i$.

\medskip\noindent
D'après Diviseurs, lemme
\ref{divisors-lemma-blow-up-pullback-effective-Cartier}, l'image réciproque
$f^{-1}C'_i$ est définie. D'après Diviseurs, lemme
\ref{divisors-lemma-effective-Cartier-divisor-is-a-sum}, on voit que
$f^{-1}C'_i = C_i + e_iC_n$ pour un certain $e_i \geq 0$. Puisque
$\mathcal{O}_X(C_i + e_iC_n) = \mathcal{O}_X(f^{-1}C'_i) =
f^*\mathcal{O}_{X'}(C'_i)$
(Diviseurs, lemme \ref{divisors-lemma-pullback-effective-Cartier-divisors})
et que l'image réciproque d'un faisceau inversible se restreint au faisceau
inversible trivial sur $C_n$, on voit que
$$
0 = \deg_{C_n}(\mathcal{O}_X(C_i + e_iC_n)) =
(C_i + e_iC_n \cdot C_n) = (C_i \cdot C_n) + e_i(C_n \cdot C_n)
$$
Comme $f_j = f|_{C_j} : C_j \to C_j$ est un morphisme propre birationnel
de courbes propres sur $k$, on voit que
$\deg_{C'_j}(\mathcal{O}_{X'}(C'_i)|_{C'_j})$ est égal à
$\deg_{C_j}(f_j^*\mathcal{O}_{X'}(C'_i)|_{C'_j})$
(Variétés, lemme \ref{varieties-lemma-degree-birational-pullback}).
En considérant le diagramme commutatif
$$
\xymatrix{
C_j \ar[r] \ar[d]_{f_j} & X \ar[d]^f \\
C'_j \ar[r] & X'
}
$$
et en utilisant Diviseurs, lemme
\ref{divisors-lemma-pullback-effective-Cartier-divisors}, on voit que
$$
(C'_i \cdot C'_j) = \deg_{C'_j}(\mathcal{O}_{X'}(C'_i)|_{C'_j}) =
\deg_{C_j}(\mathcal{O}_X(C_i + e_iC_n)) = (C_i + e_iC_n \cdot C_j)
$$
En y substituant la formule de $e_i$ obtenue ci-dessus, on voit que (1)
est vraie.
\end{proof}

\begin{remark}
\label{remark-genus-change}
Dans la situation du lemme \ref{lemma-blowdown-regular-model}, on peut aussi
décrire exactement la relation entre le genre $g_i$ de $C_i$ et le genre
$g'_i$ de $C'_i$. La formule est
$$
g'_i = \frac{w_i}{w'_i}(g_i - 1) + 1 +
\frac{(C_i \cdot C_n)^2 - w_n(C_i \cdot C_n)}{2w'_iw_n}
$$
où $w_i = [\kappa_i : k]$, $w_n = [\kappa_n : k]$ et
$w'_i = [\kappa'_i : k]$.
Pour le démontrer, considérons la suite exacte courte
$$
0 \to \mathcal{O}_{X'}(-C'_i) \to \mathcal{O}_{X'} \to
\mathcal{O}_{C'_i} \to 0
$$
et son image réciproque sur $X$, qui s'écrit
$$
0 \to \mathcal{O}_X(-C'_i - e_iC_n) \to \mathcal{O}_X \to
\mathcal{O}_{C_i + e_iC_n} \to 0
$$
où $e_i$ est comme dans la preuve du lemme
\ref{lemma-blowdown-regular-model}. Comme
$Rf_*f^*\mathcal{L} = \mathcal{L}$ pour tout module inversible
$\mathcal{L}$ sur $X'$ (détails omis), on conclut que
$$
Rf_*\mathcal{O}_{C_i + e_iC_n} = \mathcal{O}_{C'_i}
$$
en tant que complexes de faisceaux cohérents sur $X'_k$.
Les deux membres ont donc la même caractéristique d'Euler, laquelle coïncide
avec la caractéristique d'Euler de $\mathcal{O}_{C_i + e_iC_n}$ sur $X_k$.
En utilisant la suite exacte
$$
0 \to \mathcal{O}_{C_i + e_iC_n} \to
\mathcal{O}_{C_i} \oplus \mathcal{O}_{e_iC_n} \to
\mathcal{O}_{C_i \cap e_iC_n} \to 0
$$
et en filtrant encore $\mathcal{O}_{e_iC_n}$ (détails omis), on trouve
$$
\chi(\mathcal{O}_{C'_i}) =
\chi(\mathcal{O}_{C_i}) - {e_i + 1 \choose 2}(C_n \cdot C_n)
- e_i(C_i \cdot C_n)
$$
Comme $e_i = -(C_i \cdot C_n)/(C_n \cdot C_n)$ et
$(C_n \cdot C_n) = -w_n$, on en déduit la formule énoncée au début de
cette remarque. Si nous en avons un jour besoin, nous la formulerons comme
un lemme et en donnerons une preuve détaillée.
\end{remark}

\begin{remark}
\label{remark-compare-contractions}
Soit $f : X \to X'$ comme dans le lemme
\ref{lemma-blowdown-regular-model}. Soient $n, m_i, a_{ij}, w_i, g_i$ le type
numérique associé à $X$, et $n', m'_i, a'_{ij}, w'_i, g'_i$ celui associé à
$X'$. Il résulte clairement du lemme \ref{lemma-blowdown-regular-model} et
de la remarque \ref{remark-genus-change} que cela coïncide avec la
contraction des types numériques du lemme \ref{lemma-contract}, à
l'exception de la valeur de $w'_i$. Dans la situation géométrique, $w'_i$
est un entier positif divisant à la fois $w_i$ et $w_n$. Dans le cas
numérique, nous avons choisi pour $w'_i$ le plus grand entier possible
divisant $w_i$ et tel que $g'_i$ (donné par la formule) soit entier.
Ce choix convient bien dans le cadre numérique, en ce qu'il aide à comparer
les groupes de Picard des types numériques ; voir le lemme
\ref{lemma-contract-picard-group} (bien que seule l'injectivité soit
utilisée dans la suite, et que celle-ci reste vraie pour des $w'_i$ plus
petits).
\end{remark}

\begin{lemma}
\label{lemma-nonuniqueness}
Soit $C$ une courbe projective lisse sur $K$, avec
$H^0(C, \mathcal{O}_C) = K$ et de genre $0$. S'il existe plus d'un modèle
minimal de $C$, alors la fibre spéciale de tout modèle minimal est
isomorphe à $\mathbf{P}^1_k$.
\end{lemma}

\noindent
On peut renforcer ce lemme en affirmant que la transformation birationnelle
entre deux modèles minimaux non isomorphes se factorise en une suite de
transformations élémentaires comme dans l'exemple
\ref{example-nonunique-in-genus-zero}. Si nous en avons un jour besoin,
nous en donnerons ici un énoncé précis et une preuve.

\begin{proof}
Soit $X$ un modèle minimal de $C$. Le type numérique associé à $X$ est de
genre $0$ et minimal (définition \ref{definition-numerical-type-model} et
lemme \ref{lemma-numerical-type-minimal-model}). Le lemme
\ref{lemma-genus-zero} montre donc que $X_k$ est réduit, irréductible,
vérifie $H^0(X_k, \mathcal{O}_{X_k}) = k$ et est de genre $0$.
Soit $Y$ un second modèle minimal de $C$, non isomorphe à $X$.
D'après Résolution des surfaces, lemme
\ref{resolve-lemma-birational-regular-surfaces}, il existe un diagramme de
$S$-morphismes
$$
X = X_0 \leftarrow X_1 \leftarrow \ldots \leftarrow X_n = Y_m
\to \ldots \to Y_1 \to Y_0 = Y
$$
où chaque morphisme est un éclatement en un point fermé. Nous démontrerons
le lemme par récurrence sur $m$. Le cas initial est $m=0$ ; l'assertion y
est vraie puisque nous avons supposé $Y$ minimal : cela impliquerait
$n=0$, mais $X$ n'est pas isomorphe à $Y$, de sorte que ce cas ne se
présente pas et qu'il n'y a rien à vérifier.

\medskip\noindent
Avant de poursuivre, remarquons que $n+1=m+1$ est égal au nombre de
composantes irréductibles de la fibre spéciale de $X_n=Y_m$, puisque
$X_k$ et $Y_k$ sont toutes deux irréductibles. Nous utiliserons également
l'observation suivante : si $X' \to X''$ est un morphisme de modèles
propres réguliers de $C$, alors $X' \to X''$ est un isomorphisme au-dessus
d'un ouvert de $X''$ dont le complément est un ensemble fini de points
fermés de la fibre spéciale de $X''$ ; voir Variétés, lemme
\ref{varieties-lemma-modification-normal-iso-over-codimension-1}.
En fait, tout morphisme $X' \to X''$ de ce type est une suite d'éclatements
en des points fermés (Résolution des surfaces, lemme
\ref{resolve-lemma-proper-birational-regular-surfaces}), et le nombre
d'éclatements est la différence entre les nombres de composantes
irréductibles des fibres spéciales de $X'$ et de $X''$.

\medskip\noindent
Soit $E_i \subset Y_i$, $m \geq i \geq 1$, la courbe contractée par le
morphisme $Y_i \to Y_{i-1}$. Soit $i$ le plus grand indice tel que $E_i$
soit de multiplicité $>1$ dans la fibre spéciale de $Y_i$. Alors les
éclatements suivants $Y_m \to \ldots \to Y_{i+1} \to Y_i$ sont des
isomorphismes au-dessus de $E_i$ ; sinon, un $E_j$ avec $j>i$ aurait une
multiplicité $>1$. Soit $E \subset Y_m$ l'image réciproque de $E_i$.
D'après ce qui précède, $E \subset Y_m$ est une courbe exceptionnelle de
première espèce. Soit $Y_m \to Y'$ la contraction de $E$ (elle existe par
Résolution des surfaces, lemme
\ref{resolve-lemma-contract-when-quasi-projective}). Le morphisme
$Y_m \to X$ doit contracter $E$, puisque $X_k$ est réduit. Il existe donc
des morphismes $Y' \to Y$ et $Y' \to X$ (Résolution des surfaces, lemme
\ref{resolve-lemma-factor-through-contraction}) qui sont des composés d'au
plus $n-1=m-1$ contractions de courbes exceptionnelles (voir la discussion
ci-dessus). La récurrence sur $m$ permet de conclure. En définitive, on peut
supposer que les fibres spéciales de tous les $X_i$ et $Y_i$ sont réduites.

\medskip\noindent
Puisque les fibres de $X_i$ et $Y_i$ sont réduites, les éclatements
$X_i \to X_{i-1}$ et $Y_i \to Y_{i-1}$ ont nécessairement lieu en des
points fermés qui sont des points réguliers des fibres spéciales.
En effet, si $X''$ est un modèle régulier de $C$, si $x \in X''$ est un
point fermé de la fibre spéciale et si $\pi \in \mathfrak m_x^2$, alors la
fibre exceptionnelle $E$ de l'éclatement $X' \to X''$ en $x$ a une
multiplicité au moins $2$ dans la fibre spéciale de $X'$ (calcul local
omis). Ainsi
$\mathcal{O}_{X''_k,x}=\mathcal{O}_{X'',x}/\pi$ est régulier (Algèbre,
lemme \ref{algebra-lemma-regular-ring-CM}), comme annoncé. En particulier,
$x$ est un diviseur de Cartier sur l'unique composante irréductible
$Z'$ de $X''_k$ qui le contient (Variétés, lemme
\ref{varieties-lemma-regular-point-on-curve}). Il s'ensuit que la
transformée stricte $Z \subset X'$ de $Z'$ s'envoie isomorphiquement sur
$Z'$ (utiliser Diviseurs, lemmes \ref{divisors-lemma-strict-transform} et
\ref{divisors-lemma-blow-up-effective-Cartier-divisor}). Autrement dit, si
une composante irréductible $Z$ de $X_i$ n'est pas contractée par
$X_i \to X_j$ ($i>j$), elle s'envoie isomorphiquement sur son image.

\medskip\noindent
Nous pouvons maintenant démontrer le lemme. Soit $E \subset Y_m$ la courbe
exceptionnelle de première espèce contractée par $Y_m \to Y_{m-1}$.
Si $E$ est contractée par le morphisme $Y_m=X_n \to X$, il existe une
factorisation $Y_{m-1} \to X$ (Résolution des surfaces, lemme
\ref{resolve-lemma-factor-through-contraction}) ; de plus,
$Y_{m-1} \to X$ est une suite d'éclatements en des points fermés
(Résolution des surfaces, lemme
\ref{resolve-lemma-proper-birational-regular-surfaces}). Dans ce cas, on
diminue $m$ et l'on conclut par récurrence. Supposons enfin que $E$ ne soit
pas contractée par $Y_m \to X$. Alors $E \to X_k$ est surjective puisque
$X_k$ est irréductible ; ce qui précède montre qu'il s'agit d'un
isomorphisme. Ainsi $X_k$ est isomorphe à une droite projective, comme
voulu.
\end{proof}





\section{Groupes de Picard des modèles}
\label{section-picard-groups-models}

\noindent
Supposons que $R,K,k,\pi,C,X,n,C_1,\ldots,C_n,m_1,\ldots,m_n$ soient comme
dans la situation \ref{situation-regular-model}. Dans le lemme
\ref{lemma-regular-model-pic}, nous avons obtenu une suite exacte
$$
0 \to \mathbf{Z} \to \mathbf{Z}^{\oplus n} \to
\Pic(X) \to \Pic(C) \to 0
$$
Nous voulons utiliser cette suite pour étudier la $\ell$-torsion des
groupes de Picard pour des nombres premiers $\ell$ convenables.

\begin{lemma}
\label{lemma-characterize-trivial}
Dans la situation \ref{situation-regular-model}, posons
$d=\gcd(m_1,\ldots,m_n)$. Si $\mathcal{L}$ est un
$\mathcal{O}_X$-module inversible qui
\begin{enumerate}
\item se restreint au module inversible trivial sur $C$, et
\item est de degré $0$ sur chaque $C_i$,
\end{enumerate}
alors $\mathcal{L}^{\otimes d} \cong \mathcal{O}_X$.
\end{lemma}

\begin{proof}
D'après le lemme \ref{lemma-regular-model-pic}, on a
$\mathcal{L} \cong \mathcal{O}_X(\sum a_iC_i)$ pour certains
$a_i \in \mathbf{Z}$. Le degré de $\mathcal{L}|_{C_j}$ est
$\sum_j a_j(C_i \cdot C_j)$. En particulier,
$(\sum a_iC_i \cdot \sum a_iC_i)=0$. Le lemme
\ref{lemma-properties-form} montre donc que
$(a_1,\ldots,a_n)=q(m_1,\ldots,m_n)$ pour un certain $q\in\mathbf{Q}$.
Ainsi $\mathcal{L}=\mathcal{O}_X(lD)$ pour un certain
$l\in\mathbf{Z}$, où $D=\sum(m_i/d)C_i$ est comme dans le lemme
\ref{lemma-multiple-fibre-normal-bundle}, et la conclusion en résulte.
\end{proof}

\begin{lemma}
\label{lemma-canonical-map-of-pic}
Dans la situation \ref{situation-regular-model}, soit $T$ le type numérique
associé à $X$. Il existe une application canonique
$$
\Pic(C) \to \Pic(T)
$$
dont le noyau est exactement constitué des modules inversibles sur $C$ qui
sont les restrictions de modules inversibles $\mathcal{L}$ sur $X$ tels
que $\deg_{C_i}(\mathcal{L}|_{C_i})=0$ pour $i=1,\ldots,n$.
\end{lemma}

\begin{proof}
Rappelons que $w_i=[\kappa_i:k]$, où
$\kappa_i=H^0(C_i,\mathcal{O}_{C_i)})$, et que le degré de tout module
inversible sur $C_i$ est divisible par $w_i$ (Variétés, lemme
\ref{varieties-lemma-divisible}). On peut donc considérer l'application
$$
\frac{\deg}{w} : \Pic(X) \to \mathbf{Z}^{\oplus n}, \quad
\mathcal{L} \mapsto
(\frac{\deg(\mathcal{L}|_{C_1})}{w_1}, \ldots,
\frac{\deg(\mathcal{L}|_{C_n})}{w_n})
$$
L'image de $\mathcal{O}_X(C_j)$ par cette application est
$$
((C_j \cdot C_1)/w_1, \ldots, (C_j \cdot C_n)/w_n) =
(a_{1j}/w_1, \ldots, a_{nj}/w_n)
$$
qui est exactement l'image du $j$-ième vecteur de la base canonique par
l'application $(a_{ij}/w_i):\mathbf{Z}^{\oplus n}\to\mathbf{Z}^{\oplus n}$
définissant le groupe de Picard de $T$ ; voir la définition
\ref{definition-picard-group}. Ainsi l'application canonique du lemme
provient du diagramme commutatif
$$
\xymatrix{
\mathbf{Z}^{\oplus n} \ar[r] \ar[d]_{\text{id}} &
\Pic(X) \ar[r] \ar[d]^{\frac{\deg}{w}} &
\Pic(C) \ar[r] \ar[d] & 0 \\
\mathbf{Z}^{\oplus n} \ar[r]^{(a_{ij}/w_i)} &
\mathbf{Z}^{\oplus n} \ar[r] &
\Pic(T) \ar[r] & 0
}
$$
à lignes exactes (la ligne supérieure l'est par le lemme
\ref{lemma-regular-model-pic}). La description du noyau est claire.
\end{proof}

\begin{lemma}
\label{lemma-sequence-torsion}
Dans la situation \ref{situation-regular-model}, posons
$d=\gcd(m_1,\ldots,m_n)$ et soit $T$ le type numérique associé à $X$.
Soit $h\geq1$ un entier premier à $d$. Il existe une suite exacte
$$
0 \to \Pic(X)[h] \to \Pic(C)[h] \to \Pic(T)[h]
$$
\end{lemma}

\begin{proof}
En prenant la $h$-torsion dans la suite exacte du lemme
\ref{lemma-regular-model-pic}, on obtient l'exactitude de
$0\to\Pic(X)[h]\to\Pic(C)[h]$, car $h$ est premier à $d$.
L'application du lemme \ref{lemma-canonical-map-of-pic} donne une
application $\Pic(C)[h]\to\Pic(T)[h]$ qui annule les éléments de
$\Pic(X)[h]$. Réciproquement, si $\xi\in\Pic(C)[h]$ s'envoie sur zéro
dans $\Pic(T)[h]$, on peut trouver un $\mathcal{O}_X$-module inversible
$\mathcal{L}$ tel que $\deg(\mathcal{L}|_{C_i})=0$ pour tout $i$ et dont
la restriction à $C$ est $\xi$. Alors $\mathcal{L}^{\otimes h}$ est de
$d$-torsion d'après le lemme \ref{lemma-characterize-trivial}.
Soit $d'$ un entier tel que $dd'\equiv1\bmod h$. Un tel entier existe car
$h$ et $d$ sont premiers entre eux. Alors
$\mathcal{L}^{\otimes dd'}$ est un faisceau inversible de $h$-torsion sur
$X$ dont la restriction à $C$ est $\xi$.
\end{proof}

\begin{lemma}
\label{lemma-torsion-embeds}
Dans la situation \ref{situation-regular-model}, soit $h$ un entier premier
à la caractéristique de $k$. Alors l'application
$$
\Pic(X)[h] \longrightarrow \Pic((X_k)_{red})[h]
$$
est injective.
\end{lemma}

\begin{proof}
Observons que $X\times_{\Spec(R)}\Spec(R/\pi^n)$ est un épaississement
d'ordre fini de $(X_k)_{red}$ (cela résulte par exemple de Cohomologie des
schémas, lemme \ref{coherent-lemma-power-ideal-kills-sheaf}). Ainsi,
l'application canonique
$\Pic(X\times_{\Spec(R)}\Spec(R/\pi^n))\to\Pic((X_k)_{red})$ identifie les
sous-groupes de $h$-torsion d'après Compléments sur les morphismes, lemme
\ref{more-morphisms-lemma-torsion-pic-thickening}, et notre hypothèse sur
$h$. Par conséquent, si $\mathcal{L}$ est un faisceau inversible de
$h$-torsion sur $X$ dont la restriction à $(X_k)_{red}$ est triviale, alors
la restriction de $\mathcal{L}$ à
$X\times_{\Spec(R)}\Spec(R/\pi^n)$ est triviale pour tout $n$.
On obtient
\begin{align*}
H^0(X, \mathcal{L})^\wedge
& =
\lim H^0(X \times_{\Spec(R)} \Spec(R/\pi^n),
\mathcal{L}|_{X \times_{\Spec(R)} \Spec(R/\pi^n)}) \\
& \cong
\lim H^0(X \times_{\Spec(R)} \Spec(R/\pi^n),
\mathcal{O}_{X \times_{\Spec(R)} \Spec(R/\pi^n)}) \\
& =
R^\wedge
\end{align*}
On utilise ici le théorème des fonctions formelles (Cohomologie des schémas,
théorème \ref{coherent-theorem-formal-functions}) pour la première et la
dernière égalité, et, par exemple, Compléments d'algèbre, lemme
\ref{more-algebra-lemma-isomorphic-completions}, pour l'isomorphisme du
milieu. Puisque $H^0(X,\mathcal{L})$ est un $R$-module fini et que $R$ est
un anneau de valuation discrète, cela signifie que
$H^0(X,\mathcal{L})$ est libre de rang $1$ comme $R$-module.
Soit $s\in H^0(X,\mathcal{L})$ un élément de base. En remontant les
isomorphismes ci-dessus, on voit alors que
$s|_{X\times_{\Spec(R)}\Spec(R/\pi^n)}$ est une trivialisation pour tout
$n$. Comme le lieu des zéros de $s$ est fermé dans $X$ et que
$X\to\Spec(R)$ est propre, on conclut que le lieu des zéros de $s$ est vide,
comme voulu.
\end{proof}






\section{Réduction semi-stable}
\label{section-semistable-reduction}

\noindent
Dans cette section, nous définissons avec soin ce que nous entendons par
réduction semi-stable.

\begin{example}
\label{example-blowup}
Soit $R$ un anneau de valuation discrète d'uniformisante $\pi$.
Pour $n\geq0$, considérons le morphisme d'anneaux
$$
R \longrightarrow A = R[x, y]/(xy - \pi^n)
$$
Posons $X=\Spec(A)$ et $S=\Spec(R)$. Si $n=0$, le morphisme
$X\to S$ est lisse. Pour tout $n$, le morphisme $X\to S$ est à
singularités au plus nodales de dimension relative $1$, au sens de Courbes
algébriques, section \ref{curves-section-families-nodal}. Si $n=1$, alors
$X$ est régulier ; mais si $n>1$, $X$ n'est pas régulier, car $(x,y)$
n'engendre plus l'idéal maximal $\mathfrak m=(\pi,x,y)$. Pour améliorer la
situation lorsque $n>1$, considérons l'éclatement
$b:X'\to X$ de $X$ en $\mathfrak m$. Voir Diviseurs, section
\ref{divisors-section-blowing-up}. Par construction, $X'$ est recouvert par
trois ouverts affines correspondant aux algèbres d'éclatement
$A[\frac{\mathfrak m}{\pi}]$, $A[\frac{\mathfrak m}{x}]$ et
$A[\frac{\mathfrak m}{y}]$.

\medskip\noindent
L'algèbre $A[\frac{\mathfrak m}{\pi}]$ possède les générateurs
$x'=x/\pi$ et $y'=y/\pi$, et $x'y'=\pi^{n-2}$. Cette partie de $X'$ est
donc le spectre de $R[x',y'](x'y'-\pi^{n-2})$.

\medskip\noindent
L'algèbre $A[\frac{\mathfrak m}{x}]$ possède les générateurs $x$ et
$u=\pi/x$, soumis à la relation $xu-\pi$. Notons que cet anneau contient
$y/x=\pi^n/x^2=u^2\pi^{n-2}$. Cette partie de $X'$ est donc régulière.

\medskip\noindent
Par symétrie, le cas de l'algèbre $A[\frac{\mathfrak m}{y}]$ est identique
à celui de $A[\frac{\mathfrak m}{x}]$.

\medskip\noindent
On voit ainsi que $X'\to S$ est à singularités au plus nodales de dimension
relative $1$ et que $X'$ est régulier, sauf en un point qui possède un
voisinage ouvert affine exactement comme ci-dessus, mais avec $n$ remplacé
par $n-2$. Par récurrence sur $n$, on conclut qu'il existe une suite
d'éclatements en des points fermés
$$
X_{\lfloor n/2 \rfloor} \to \ldots \to X_1 \to X_0 = X
$$
telle que $X_{\lfloor n/2 \rfloor}\to S$ soit à singularités au plus
nodales de dimension relative $1$ et que
$X_{\lfloor n/2 \rfloor}$ soit régulier.
\end{example}

\begin{lemma}
\label{lemma-etale-local-at-worst-nodal}
Soit $R$ un anneau de valuation discrète. Soit $X$ un schéma à singularités
au plus nodales de dimension relative $1$ sur $R$. Soit $x\in X$ un point
de la fibre spéciale de $X$ sur $R$. Il existe alors un diagramme commutatif
$$
\xymatrix{
X \ar[d] &
U \ar[r] \ar[d] \ar[l] &
\Spec(A) \ar[dl] \\
\Spec(R) &
\Spec(R') \ar[l]
}
$$
où $R\subset R'$ est une extension étale d'anneaux de valuation discrète,
le morphisme $U\to X$ est étale, le morphisme $U\to\Spec(A)$ est étale,
il existe un point $x'\in U$ d'image $x$, et
$$
A = R'[u, v]/(uv)
\quad\text{ou}\quad
A = R'[u, v]/(uv - \pi^n)
$$
où $n\geq0$ et $\pi\in R'$ est une uniformisante.
\end{lemma}

\begin{proof}
Nous avons déjà démontré ce lemme dans une généralité bien plus grande ;
voir Courbes algébriques, lemme
\ref{curves-lemma-etale-local-structure-nodal-family}. Il suffit ici de
traduire l'énoncé qui y est donné sous la forme de celui ci-dessus.

\medskip\noindent
Tout d'abord, si $X\to\Spec(R)$ est lisse en $x$, on peut trouver un
morphisme étale $U\to\mathbf{A}^1_R=\Spec(R[u])$ pour un voisinage ouvert
affine $U\subset X$ de $x$. C'est le lemme
\ref{morphisms-lemma-smooth-etale-over-affine-space} de Morphismes.
Après avoir, si nécessaire, remplacé la coordonnée $u$ par $u+1$, on peut
supposer que $x$ s'envoie sur un point de l'ouvert standard
$D(u)\subset\mathbf{A}^1_R$. Alors $D(u)=\Spec(A)$, avec
$A=R[u,v]/(uv-1)$, et le résultat est vrai dans ce cas.

\medskip\noindent
Supposons ensuite que $x$ soit un point singulier de la fibre. On peut
alors appliquer Courbes algébriques, lemme
\ref{curves-lemma-etale-local-structure-nodal-family}, pour obtenir un
diagramme
$$
\xymatrix{
X \ar[d] &
U \ar[rr] \ar[l] \ar[rd] & &
W \ar[r] \ar[ld] &
\Spec(\mathbf{Z}[u, v, a]/(uv - a)) \ar[d] \\
\Spec(R) & &
V \ar[ll] \ar[rr] & & \Spec(\mathbf{Z}[a])
}
$$
ayant toutes les propriétés mentionnées dans l'énoncé du lemme cité.
Soit $x'\in U$ le point d'image $x$ fourni par ce lemme. Rétrécissons
d'abord $V$ en un voisinage affine de l'image de $x'$. Écrivons
$V=\Spec(R')$. Alors $R\to R'$ est étale. Puisque $R$ est un anneau de
valuation discrète, $R'$ est un produit fini de domaines de Dedekind
quasi-locaux (utiliser Compléments d'algèbre, lemme
\ref{more-algebra-lemma-Dedekind-etale-extension}). On trouve donc, par
exemple à l'aide de l'évitement des idéaux premiers, un ouvert standard
$D(f)\subset V=\Spec(R')$ contenant l'image de $x'$ et tel que $R'_f$
soit un anneau de valuation discrète. En remplaçant $R'$ par $R'_f$, on se
ramène à la situation où $V=\Spec(R')$, avec $R\subset R'$ une extension
étale d'anneaux de valuation discrète (les extensions d'anneaux de valuation
discrète sont définies dans Compléments d'algèbre, définition
\ref{more-algebra-definition-extension-discrete-valuation-rings}).

\medskip\noindent
Le morphisme $V\to\Spec(\mathbf{Z}[a])$ est déterminé par l'image $h$ de
$a$ dans $R'$. Alors $W=\Spec(R'[u,v]/(uv-h))$. Ainsi le lemme est vrai
avec $A=R'[u,v]/(uv-h)$. Si $h=0$, on obtient clairement le premier cas
du lemme. Si $h\not = 0$, on peut écrire $h=\epsilon\pi^n$ pour un certain
$n\geq0$, où $\epsilon$ est une unité de $R'$. En effectuant le changement
de coordonnées $u_{new}=\epsilon u$ et $v_{new}=v$, on obtient le second
type d'isomorphisme de $A$ indiqué dans le lemme.
\end{proof}

\begin{lemma}
\label{lemma-blowup-at-worst-nodal}
Soit $R$ un anneau de valuation discrète. Soit $X$ un schéma quasi-compact
à singularités au plus nodales de dimension relative $1$ sur $R$, dont la
fibre générique est lisse. Il existe alors $m\geq0$ et une suite
$$
X_m \to \ldots \to X_1 \to X_0 = X
$$
tels que
\begin{enumerate}
\item $X_{i+1}\to X_i$ soit l'éclatement d'un point fermé $x_i$ où $X_i$
est singulier ;
\item $X_i\to\Spec(R)$ soit à singularités au plus nodales de dimension
relative $1$ ;
\item $X_m$ soit régulier.
\end{enumerate}
\end{lemma}

\noindent
Un énoncé légèrement plus fort, également vrai, serait que, quelle que soit
la manière dont on éclate les points singuliers, on aboutit finalement à
une résolution et que tous les éclatements intermédiaires sont à
singularités au plus nodales de dimension relative $1$ sur $R$.

\begin{proof}
Comme $X$ est quasi-compact, la fibre spéciale $X_k$ est quasi-compacte.
Puisque les singularités de $X_k$ sont au plus nodales, $X_k$ possède un nombre fini
de nœuds et est lisse sur $k$ ailleurs. Comme $X\to\Spec(R)$ est plat à
fibre générique lisse, il s'ensuit que $X$ est lisse sur $R$ hors du nombre
fini de nœuds de $X_k$ (utiliser Morphismes, lemme
\ref{morphisms-lemma-smooth-at-point}). Il s'ensuit que $X$ est régulier en
tout point, sauf éventuellement aux nœuds de sa fibre spéciale (voir
Algèbre, lemme \ref{algebra-lemma-regular-goes-up}). Soit $x\in X$ un tel
nœud. Choisissons un diagramme
$$
\xymatrix{
X \ar[d] &
U \ar[r] \ar[d] \ar[l] &
\Spec(A) \ar[dl] \\
\Spec(R) &
\Spec(R') \ar[l]
}
$$
comme dans le lemme \ref{lemma-etale-local-at-worst-nodal}.
Le cas $A=R'[u,v]/(uv)$ ne peut pas se produire, car il impliquerait que la
fibre générique de $X/R$ est singulière (petit détail omis). Ainsi
$A=R'[u,v]/(uv-\pi^n)$ pour un certain $n\geq0$. Puisque $x$ est un point
singulier, on a $n\geq2$ ; voir la discussion de l'exemple
\ref{example-blowup}.

\medskip\noindent
Après avoir rétréci $U$, on peut supposer qu'il existe un unique point
$u\in U$ d'image $x$. Soit $w\in\Spec(A)$ l'image de $u$. On peut aussi
supposer que $u$ est l'unique point de $U$ s'envoyant sur $w$. Puisque les
deux flèches horizontales sont étales, on voit que $u$, considéré comme
sous-schéma fermé de $U$, est l'image réciproque schématique de $x\in X$
et celle de $w\in\Spec(A)$. Comme l'éclatement commute au changement de
base plat (Diviseurs, lemme
\ref{divisors-lemma-flat-base-change-blowing-up}), on obtient un diagramme
commutatif
$$
\xymatrix{
X' \ar[d] &
U' \ar[l] \ar[d] \ar[r] &
W' \ar[d] \\
X & U \ar[l] \ar[r] & \Spec(A)
}
$$
à carrés cartésiens, où les flèches verticales sont les éclatements de
$x,u,w$ dans $X,U,\Spec(A)$. Le schéma $W'$ a été décrit dans l'exemple
\ref{example-blowup}. Nous y avons vu que $W'$ est à singularités au plus
nodales de dimension relative $1$ sur $R'$. Ainsi, $W'$ est à singularités
au plus nodales de dimension relative $1$ sur $R$ (Courbes algébriques, lemme
\ref{curves-lemma-nodal-family-postcompose-etale}). Par conséquent, $U'$ est
à singularités au plus nodales de dimension relative $1$ sur $R$ (voir
Courbes algébriques, lemme
\ref{curves-lemma-nodal-family-etale-local-source}). Puisque $X'\to X$ est
un isomorphisme au-dessus du complémentaire de $x$, on conclut que
$X'/R$ possède la même propriété (en appliquant encore Courbes algébriques,
lemme \ref{curves-lemma-nodal-family-etale-local-source}).

\medskip\noindent
Enfin, il faut montrer qu'après un nombre fini de ces éclatements on atteint
un modèle régulier $X_m$. C'est assez clair, car l'« invariant » $n$
diminue de $2$ sous l'éclatement décrit ci-dessus ; voir le calcul de
l'exemple \ref{example-blowup}. Toutefois, pour éviter de définir
précisément cet invariant et d'en établir les propriétés, raisonnons comme
suit. Si $n=2$, alors $W'$ est régulier, donc $X'$ est régulier en tous les
points situés au-dessus de $x$, et le nombre de points singuliers de $X$ a
diminué de $1$. Si $n>2$, l'unique point singulier $w'$ de $W'$ situé
au-dessus de $w$ vérifie $\kappa(w)=\kappa(w')$. Ainsi $U'$ possède un
unique point singulier $u'$ situé au-dessus de $u$, avec
$\kappa(u)=\kappa(u')$. Il est clair que cela implique que $X'$ possède un
unique point singulier $x'$ situé au-dessus de $x$, à savoir l'image de
$u'$. On peut donc raisonner exactement comme ci-dessus et obtenir un
diagramme commutatif
$$
\xymatrix{
X'' \ar[d] &
U'' \ar[l] \ar[d] \ar[r] &
W'' \ar[d] \\
X' & U' \ar[l] \ar[r] & W'
}
$$
à carrés cartésiens, où les flèches verticales sont les éclatements de
$x',u',w'$ dans $X',U',W'$. En poursuivant ainsi, on obtient une suite
compatible d'éclatements qui s'arrête après $\lfloor n/2\rfloor$ étapes.
À l'issue de ce processus, le schéma
$X^{(\lfloor n/2\rfloor)}$ possède un point singulier de moins que $X$.
Une récurrence sur le nombre de points singuliers achève la preuve.
\end{proof}

\begin{lemma}
\label{lemma-blowdown-at-worst-nodal}
Soit $R$ un anneau de valuation discrète, de corps des fractions $K$ et de
corps résiduel $k$. Supposons que $X\to\Spec(R)$ soit à singularités au
plus nodales de dimension relative $1$ sur $R$. Soit $X\to X'$ la
contraction d'une courbe exceptionnelle de première espèce $E\subset X$.
Alors $X'$ est à singularités au plus nodales de dimension relative $1$
sur $R$.
\end{lemma}

\begin{proof}
Soit $x'\in X'$ l'image de $E$. Il suffit de voir que
$X'\to\Spec(R)$ est à singularités au plus nodales de dimension relative
$1$ dans un voisinage de $x'$. La fibre fermée de $X\to\Spec(R)$ est
réduite ; ainsi $\pi\in R$ s'annule à l'ordre $1$ sur $E$. Cela implique
immédiatement que $\pi$, considéré comme élément de
$\mathfrak m_{x'}\subset\mathcal{O}_{X',x'}$, n'appartient pas à
$\mathfrak m_{x'}^2$. Comme $\mathcal{O}_{X',x'}$ est régulier de
dimension $2$ (par la définition des contractions dans Résolution des
surfaces, section \ref{resolve-section-minus-one}), il s'ensuit que
$\mathcal{O}_{X'_k,x'}$ est régulier de dimension $1$ (Algèbre, lemme
\ref{algebra-lemma-regular-ring-CM}). D'autre part, la courbe $E$ doit
rencontrer au moins une autre composante, disons $C$, de la fibre fermée
$X_k$. Soit $x\in E\cap C$. Alors $x$ est un nœud de la fibre spéciale
$X_k$, donc $\kappa(x)/k$ est finie séparable ; voir Courbes algébriques,
lemme \ref{curves-lemma-nodal}. Comme $x\mapsto x'$, on conclut que
$\kappa(x')/k$ est finie séparable. D'après Algèbre, lemme
\ref{algebra-lemma-separable-smooth}, le morphisme
$X'_k\to\Spec(k)$ est donc lisse dans un voisinage ouvert de $x'$.
Combiné à la platitude, cela montre que $X'\to\Spec(R)$ est lisse dans un
voisinage de $x'$ (Morphismes, lemme
\ref{morphisms-lemma-smooth-at-point}). Ceci achève la preuve, car un
morphisme lisse de dimension relative $1$ est à singularités au plus
nodales de dimension relative $1$ (Courbes algébriques, lemme
\ref{curves-lemma-smooth-relative-dimension-1}).
\end{proof}

\begin{lemma}
\label{lemma-semistable}
Soit $R$ un anneau de valuation discrète de corps des fractions $K$.
Soit $C$ une courbe projective lisse sur $K$ telle que
$H^0(C,\mathcal{O}_C)=K$. Les conditions suivantes sont équivalentes :
\begin{enumerate}
\item il existe un modèle propre de $C$ qui est à singularités au plus
nodales de dimension relative $1$ sur $R$ ;
\item il existe un modèle minimal de $C$ qui est à singularités au plus
nodales de dimension relative $1$ sur $R$ ;
\item tout modèle minimal de $C$ est à singularités au plus nodales de
dimension relative $1$ sur $R$.
\end{enumerate}
\end{lemma}

\begin{proof}
Pour donner un sens à cet énoncé, rappelons qu'un modèle minimal est, par
définition, un modèle propre régulier sans courbes exceptionnelles de
première espèce (définition \ref{definition-minimal-model}), que les modèles
minimaux existent (proposition \ref{proposition-exists-minimal-model}), et
qu'ils sont uniques si le genre de $C$ est $>0$ (lemme
\ref{lemma-minimal-model-unique}). Cela étant, les implications
(2) $\Rightarrow$ (1) et (3) $\Rightarrow$ (2) sont claires.

\medskip\noindent
Supposons (1). Soit $X$ un modèle propre de $C$ à singularités au plus
nodales de dimension relative $1$ sur $R$. En appliquant le lemme
\ref{lemma-blowup-at-worst-nodal}, on voit que l'on peut en outre supposer
$X$ régulier. Soit
$$
X = X_m \to X_{m - 1} \to \ldots \to X_1 \to X_0
$$
comme dans le lemme \ref{lemma-pre-exists-minimal-model}. Le lemme
\ref{lemma-blowdown-at-worst-nodal} et une récurrence montrent alors que
$X_0$ est à singularités au plus nodales de dimension relative $1$ sur
$R$.

\medskip\noindent
Pour achever la preuve, il reste à montrer que (2) implique (3). Cela est
clair si le genre de $C$ est $>0$, puisque le modèle minimal est alors
unique (voir la discussion ci-dessus). D'autre part, si le modèle minimal
n'est pas unique, alors $X\to\Spec(R)$ est lisse pour tout modèle minimal,
car sa fibre spéciale est isomorphe à $\mathbf{P}^1_k$ d'après le lemme
\ref{lemma-nonuniqueness}.
\end{proof}

\begin{definition}
\label{definition-semistable}
Soit $R$ un anneau de valuation discrète de corps des fractions $K$.
Soit $C$ une courbe projective lisse sur $K$ telle que
$H^0(C,\mathcal{O}_C)=K$. On dit que $C$ a une
{\it réduction semi-stable} si les conditions équivalentes du lemme
\ref{lemma-semistable} sont satisfaites.
\end{definition}

\begin{lemma}
\label{lemma-good}
Soit $R$ un anneau de valuation discrète de corps des fractions $K$.
Soit $C$ une courbe projective lisse sur $K$ telle que
$H^0(C,\mathcal{O}_C)=K$. Les conditions suivantes sont équivalentes :
\begin{enumerate}
\item il existe un modèle propre et lisse de $C$ ;
\item il existe un modèle minimal de $C$ qui est lisse sur $R$ ;
\item tout modèle minimal est lisse sur $R$.
\end{enumerate}
\end{lemma}

\begin{proof}
Si $X$ est un modèle propre et lisse, sa fibre spéciale est connexe
(lemme \ref{lemma-regular-model-connected}) et lisse, donc irréductible.
Cela implique immédiatement que ce modèle est minimal. Ainsi (1) implique (2).
Pour achever la preuve, il reste à montrer que (2) implique (3). Cela est
clair si le genre de $C$ est $>0$, puisque le modèle minimal est alors
unique (lemme \ref{lemma-minimal-model-unique}). D'autre part, si le modèle
minimal n'est pas unique, le morphisme $X\to\Spec(R)$ est lisse pour tout
modèle minimal, car sa fibre spéciale est isomorphe à $\mathbf{P}^1_k$
d'après le lemme \ref{lemma-nonuniqueness}.
\end{proof}

\begin{definition}
\label{definition-good}
Soit $R$ un anneau de valuation discrète de corps des fractions $K$.
Soit $C$ une courbe projective lisse sur $K$ telle que
$H^0(C,\mathcal{O}_C)=K$. On dit que $C$ a une {\it bonne réduction} si
les conditions équivalentes du lemme \ref{lemma-good} sont satisfaites.
\end{definition}








\section{Réduction semi-stable en genre zéro}
\label{section-semistable-reduction-genus-zero}

\noindent
Dans cette section, nous démontrons le théorème de réduction semi-stable
(théorème \ref{theorem-semistable-reduction}) pour les courbes de genre
zéro.

\medskip\noindent
Soit $R$ un anneau de valuation discrète de corps des fractions $K$.
Soit $C$ une courbe projective lisse sur $K$ telle que
$H^0(C,\mathcal{O}_C)=K$. Si le genre de $C$ est $0$, alors $C$ est
isomorphe à une conique ; voir Courbes algébriques, lemme
\ref{curves-lemma-genus-zero}. Il existe donc une extension finie séparable
$K'/K$ de degré au plus $2$ telle que $C(K') \not = \emptyset$ ; voir Courbes
algébriques, lemme
\ref{curves-lemma-smooth-plane-curve-point-over-separable}. Soit
$R'\subset K'$ la clôture intégrale de $R$ ; voir la discussion de
Compléments d'algèbre, remarque
\ref{more-algebra-remark-finite-separable-extension}. Nous montrerons que
$C_{K'}$ a une réduction semi-stable sur $R'_{\mathfrak m}$ pour tout idéal
maximal $\mathfrak m$ de $R'$ (dans le cas présent, il y en a bien sûr au
plus deux). En remplaçant $R$ par $R'_{\mathfrak m}$ et $C$ par $C_{K'}$,
on se ramène au cas du paragraphe suivant.

\medskip\noindent
Dans ce paragraphe, $R$ est un anneau de valuation discrète de corps des
fractions $K$, $C$ est une courbe projective lisse sur $K$ telle que
$H^0(C,\mathcal{O}_C)=K$, de genre $0$, et $C$ possède un point
$K$-rationnel. Dans ce cas, $C\cong\mathbf{P}^1_K$ d'après Courbes
algébriques, proposition \ref{curves-proposition-projective-line}. On peut
donc prendre $\mathbf{P}^1_R$ comme modèle, ce qui montre que $C$ a à la
fois bonne réduction et réduction semi-stable.

\begin{example}
\label{example-extension-necessary-genus-zero}
Soit $R=\mathbf{R}[[\pi]]$ et considérons le schéma
$$
X = V(T_1^2 + T_2^2 - \pi T_0^2) \subset \mathbf{P}^2_R
$$
Le changement de base de $X$ à $\mathbf{C}[[\pi]]$ est isomorphe au schéma
défini dans l'exemple \ref{example-nonunique-in-genus-zero}, car on a la
factorisation $T_1^2+T_2^2=(T_1+iT_2)(T_1-iT_2)$ sur $\mathbf{C}$.
Ainsi $X$ est régulier et sa fibre spéciale est irréductible mais
singulière ; par conséquent, $X$ est l'unique modèle minimal de sa fibre
générique (utiliser le lemme \ref{lemma-nonuniqueness}). Il s'ensuit qu'une
extension est nécessaire même en genre $0$.
\end{example}



\section{Réduction semi-stable en genre un}
\label{section-semistable-reduction-genus-one}

\noindent
Dans cette section, nous démontrons le théorème de réduction semi-stable
(théorème \ref{theorem-semistable-reduction}) pour les courbes de genre un.
Nous conseillons au lecteur de lire d'abord la preuve dans le cas du genre
$\geq2$ (section \ref{section-semistable-reduction-genus-at-least-two}).
Nous utiliserons autant que possible la classification des types numériques
minimaux de genre $1$ donnée dans le lemme \ref{lemma-genus-one}.

\medskip\noindent
Soit $R$ un anneau de valuation discrète de corps des fractions $K$.
Soit $C$ une courbe projective lisse sur $K$ telle que
$H^0(C,\mathcal{O}_C)=K$. Supposons que le genre de $C$ soit $1$.
Choisissons un nombre premier $\ell\geq7$ distinct de la caractéristique
de $k$. Choisissons une extension finie séparable $K'/K$ telle que
$C(K') \not = \emptyset$ et
$\Pic(C_{K'})[\ell]\cong(\mathbf{Z}/\ell\mathbf{Z})^{\oplus2}$.
Voir Courbes algébriques, lemme
\ref{curves-lemma-torsion-picard-becomes-visible}. Soit $R'\subset K'$ la
clôture intégrale de $R$ ; voir la discussion de Compléments d'algèbre,
remarque \ref{more-algebra-remark-finite-separable-extension}. On peut
remplacer $R$ par $R'_{\mathfrak m}$ pour un idéal maximal $\mathfrak m$
de $R'$, et $C$ par $C_{K'}$. On se ramène ainsi au cas du paragraphe
suivant.

\medskip\noindent
Dans le reste de cette section, $R$ est un anneau de valuation discrète de
corps des fractions $K$, $C$ est une courbe projective lisse sur $K$ telle
que $H^0(C,\mathcal{O}_C)=K$, de genre $1$, possédant un point
$K$-rationnel et telle que
$\Pic(C)[\ell]\cong(\mathbf{Z}/\ell\mathbf{Z})^{\oplus2}$ pour un nombre
premier $\ell\geq7$ distinct de la caractéristique de $k$. Nous allons
montrer que $C$ a une réduction semi-stable.

\medskip\noindent
Soit $X$ un modèle minimal de $C$ ; voir la proposition
\ref{proposition-exists-minimal-model}. Soit
$T=(n, m_i, (a_{ij}), w_i, g_i)$ le type numérique associé à $X$
(définition \ref{definition-numerical-type-model}). Alors $T$ est un type
numérique minimal (lemme \ref{lemma-numerical-type-minimal-model}).
Puisque $C$ possède un point rationnel, il existe un $i$ tel que
$m_i=w_i=1$, d'après le lemme
\ref{lemma-numerical-type-rational-point}. En examinant la classification
des types numériques minimaux de genre $1$ du lemme
\ref{lemma-genus-one}, on voit que $m=w=1$ et que les cas
(\ref{item-up4}),
(\ref{item-equal-down3}),
(\ref{item-up-up}),
(\ref{item-down-up}),
(\ref{item-up-equal-up}),
(\ref{item-down-equal-up}),
(\ref{item-triple-with-down}),
(\ref{item-equal-equal-down-equal}),
(\ref{item-up-equal-equal-up}),
(\ref{item-down-equal-equal-up}),
(\ref{item-triple-extended-down}),
(\ref{item-up-chain-equal-up}),
(\ref{item-down-chain-equal-up}) et
(\ref{item-Dn-extended-down}) sont exclus (car il n'existe aucun indice
pour lequel $w_i$ et $m_i$ sont tous deux égaux à $1$). Soit $e$ le nombre de paires $(i,j)$
avec $i<j$ et $a_{ij}>0$. Pour les cas restants, on a
\begin{enumerate}
\item[(A)] $e=n-1$ dans les cas
(\ref{item-one}),
(\ref{item-two-cycle}),
(\ref{item-equal-up3}),
(\ref{item-up-down}),
(\ref{item-up-equal-down}),
(\ref{item-triple-with-up}),
(\ref{item-equal-equal-up-equal}),
(\ref{item-up-equal-equal-down}),
(\ref{item-quadruple}),
(\ref{item-triple-extended-up}),
(\ref{item-up-chain-equal-down}),
(\ref{item-Dn-extended-up}),
(\ref{item-double-triple}),
(\ref{item-E6-completed}),
(\ref{item-E7-completed}) et
(\ref{item-E8-completed}) ;
\item[(B)] $e=n$ dans les cas
(\ref{item-three-cycle}),
(\ref{item-four-cycle}),
(\ref{item-five-cycle}) et
(\ref{item-n-cycle}).
\end{enumerate}
Nous traiterons ces deux cas séparément.

\medskip\noindent
Cas (A). Dans ce cas, $\Pic(T)[\ell]$ est trivial (le groupe de Picard d'un
type numérique est défini dans la section \ref{section-picard-group}).
Cette annulation résulte de ce que $\Pic(T)\subset\Coker(A)$ (lemme
\ref{lemma-picard-T-and-A}) et de ce que $\Coker(A)[\ell]=0$ d'après le
lemme \ref{lemma-recurring-symmetric-integer} et le fait que $\ell$ a été
choisi premier aux $a_{ij}$ et aux $m_i$. D'après les lemmes
\ref{lemma-sequence-torsion} et \ref{lemma-torsion-embeds}, on conclut qu'il
existe un plongement
$$
(\mathbf{Z}/\ell \mathbf{Z})^{\oplus 2} \subset \Pic((X_k)_{red})[\ell].
$$
D'après Courbes algébriques, lemme
\ref{curves-lemma-bound-torsion-simple}, on obtient
$$
2 \leq \dim_k H^1((X_k)_{red}, \mathcal{O}_{(X_k)_{red}}) +
g_{geom}((X_k)_{red}/k)
$$
D'après Courbes algébriques, lemmes
\ref{curves-lemma-bound-geometric-genus} et
\ref{curves-lemma-bound-geometric-genus-curve}, on a
$g_{geom}((X_k)_{red}/k)\leq\sum w_ig_i$. Les hypothèses du lemme
\ref{lemma-genus-reduction-smaller} sont satisfaites d'après le lemme
\ref{lemma-numerical-type-rational-point}, et l'on conclut que
$\dim_kH^1((X_k)_{red},\mathcal{O}_{(X_k)_{red}})\leq g=1$.
En combinant ces inégalités, on obtient
$$
2 \leq 1 + \sum w_i g_i
$$
L'examen de la liste montre que le type numérique est donné par
$n=1$, $w_1=m_1=g_1=1$. Comme toutes les inégalités sont des égalités,
on voit que $g_{geom}(C_1/k)=1$. D'autre part, on sait que $C_1$ possède
un point $k$-rationnel $x$ tel que $C_1\to\Spec(k)$ soit lisse en $x$.
Il s'ensuit que $C_1$ est géométriquement intègre (Variétés, lemme
\ref{varieties-lemma-variety-with-smooth-rational-point}). Ainsi
$g_{geom}(C_1/k)=1$ est à la fois égal au genre de la normalisation de
$C_{1,\overline{k}}$ et au genre de $C_{1,\overline{k}}$. Il s'ensuit que
le morphisme de normalisation
$C_{1,\overline{k}}^\nu\to C_{1,\overline{k}}$ est un isomorphisme
(Courbes algébriques, lemme \ref{curves-lemma-genus-normalization}). On
conclut que $C_1$ est lisse sur $k$, comme voulu.

\medskip\noindent
Cas (B). Ici, on conclut seulement qu'il existe un plongement
$$
\mathbf{Z}/\ell \mathbf{Z} \subset \Pic(X_k)[\ell]
$$
La classification des types montre que $m_i=w_i=1$ et $g_i=0$ pour tout
$i$. Chaque $C_i$ est donc une courbe de genre zéro sur $k$. De plus, pour
tout $i$, il existe un $j$ tel que $C_i\cap C_j$ soit un point
$k$-rationnel. Par conséquent, $C_i\cong\mathbf{P}^1_k$ d'après Courbes
algébriques, proposition \ref{curves-proposition-projective-line}.
En particulier, puisque $X_k$ est la réunion schématique des $C_i$, on voit
que $X_{\overline{k}}$ est la réunion schématique des
$C_{i,\overline{k}}$. Ainsi $X_{\overline{k}}$ est un schéma propre réduit
connexe de dimension $1$ sur $\overline{k}$ tel que
$\dim_{\overline{k}}H^1(X_{\overline{k}},
\mathcal{O}_{X_{\overline{k}}})=1$. De plus, d'après Variétés, lemme
\ref{varieties-lemma-change-fields-pic}, et ce qui précède, on a encore
$$
\dim_{\mathbf{F}_\ell}(\Pic(X_{\overline{k}}) \geq 1
$$
D'après Courbes algébriques, proposition
\ref{curves-proposition-torsion-picard-reduced-proper}, on voit que
$X_{\overline{k}}$ n'a que des singularités en croisement multiple.
Mais $X_k$ est de Gorenstein (lemme \ref{lemma-gorenstein}), donc
$X_{\overline{k}}$ l'est aussi (Dualité pour les schémas, lemme
\ref{duality-lemma-gorenstein-base-change}). On conclut que les singularités
de $X_{\overline{k}}$ sont au plus nodales d'après Courbes algébriques,
lemme \ref{curves-lemma-multicross-gorenstein-is-nodal}. Ceci achève la
preuve dans le cas (B).

\begin{example}
\label{example-extension-necessary-genus-one}
Soit $k$ un corps algébriquement clos. Soit $Z$ une courbe projective lisse
sur $k$ de genre positif $g$. Soit $n\geq1$ un entier premier à la
caractéristique de $k$. Soit $\mathcal{L}$ un $\mathcal{O}_Z$-module
inversible d'ordre $n$ ; voir Courbes algébriques, lemme
\ref{curves-lemma-torsion-picard-smooth-projective}. Choisissons un
isomorphisme $\varphi:\mathcal{L}^{\otimes n}\to\mathcal{O}_Z$.
Posons $R=k[[\pi]]$, de corps des fractions $K=k((\pi))$. Notons $Z_R$ le
changement de base de $Z$ à $R$, et $\mathcal{L}_R$ l'image réciproque de
$\mathcal{L}$ sur $Z_R$. Considérons le morphisme fini plat
$$
p : X \longrightarrow Z_R
$$
tel que
$$
p_*\mathcal{O}_X =
\text{Sym}^*_{\mathcal{O}_{Z_R}}(\mathcal{L}_R)/(\varphi - \pi) =
\mathcal{O}_{Z_R} \oplus \mathcal{L}_R \oplus
\mathcal{L}_R^{\otimes 2} \oplus \ldots \oplus \mathcal{L}_R^{\otimes n - 1}
$$
Plus précisément, si $U=\Spec(A)\subset Z$ est un ouvert affine tel que
$\mathcal{L}|_U$ soit trivialisé par une section $s$ vérifiant
$\varphi(s^{\otimes n})=f$ (où $f$ est une unité), alors
$$
p^{-1}(U_R) = \Spec\left(
(A \otimes_R R[[\pi]])[x]/(x^n - \pi f)
\right)
$$
Le lecteur vérifie que le morphisme $X_K\to Z_K$ entre les fibres
génériques est fini étale. La description du faisceau structural montre
que $H^0(X,\mathcal{O}_X)=R$ et $H^0(X_K,\mathcal{O}_{X_K})=K$.
D'après Riemann--Hurwitz (Courbes algébriques, lemme
\ref{curves-lemma-rhe}), le genre de $X_K$ est $n(g-1)+1$. En particulier,
$X_K$ est de genre $1$ si $Z$ est de genre $1$. D'autre part, le schéma
$X$ est régulier d'après l'équation locale ci-dessus, et la fibre spéciale
$X_k$ est, comme diviseur de Cartier effectif, égale à $n$ fois sa fibre
spéciale réduite. Il s'ensuit que toute extension finie $K'/K$ sur laquelle
$X_K$ acquiert une réduction semi-stable doit être ramifiée avec un indice
de ramification au moins $n$ (certains détails sont omis). Il n'existe donc
pas de borne universelle pour le degré d'une extension sur laquelle une
courbe de genre $1$ acquiert une réduction semi-stable.
\end{example}






\section{Réduction semi-stable en genre au moins deux}
\label{section-semistable-reduction-genus-at-least-two}

\noindent
Dans cette section, nous démontrons le théorème de réduction semi-stable
(théorème \ref{theorem-semistable-reduction}) pour les courbes de genre
$\geq2$. Fixons $g\geq2$.

\medskip\noindent
Soit $R$ un anneau de valuation discrète de corps des fractions $K$.
Soit $C$ une courbe projective lisse sur $K$ telle que
$H^0(C,\mathcal{O}_C)=K$. Supposons que le genre de $C$ soit $g$.
Choisissons un nombre premier $\ell>768g$ distinct de la caractéristique de
$k$. Choisissons une extension finie séparable $K'/K$ telle que
$C(K') \not = \emptyset$ et
$\Pic(C_{K'})[\ell]\cong(\mathbf{Z}/\ell\mathbf{Z})^{\oplus2g}$.
Voir Courbes algébriques, lemme
\ref{curves-lemma-torsion-picard-becomes-visible}. Soit $R'\subset K'$ la
clôture intégrale de $R$ ; voir la discussion de Compléments d'algèbre,
remarque \ref{more-algebra-remark-finite-separable-extension}. On peut
remplacer $R$ par $R'_{\mathfrak m}$ pour un idéal maximal $\mathfrak m$
de $R'$, et $C$ par $C_{K'}$. On se ramène ainsi au cas du paragraphe
suivant.

\medskip\noindent
Dans le reste de cette section, $R$ est un anneau de valuation discrète de
corps des fractions $K$, $C$ est une courbe projective lisse sur $K$ telle
que $H^0(C,\mathcal{O}_C)=K$, de genre $g$, possédant un point
$K$-rationnel et telle que
$\Pic(C)[\ell]\cong(\mathbf{Z}/\ell\mathbf{Z})^{\oplus2g}$ pour un nombre
premier $\ell\geq768g$ distinct de la caractéristique de $k$. Nous allons
montrer que $C$ a une réduction semi-stable.

\medskip\noindent
Dans le reste de cette section, nous utiliserons sans autre mention le fait
que les conclusions du lemme \ref{lemma-numerical-type-rational-point}
sont vraies.

\medskip\noindent
Soit $X$ un modèle minimal de $C$ ; voir la proposition
\ref{proposition-exists-minimal-model}. Soit
$T=(n, m_i, (a_{ij}), w_i, g_i)$ le type numérique associé à $X$
(définition \ref{definition-numerical-type-model}). Alors $T$ est un type
numérique minimal de genre $g$ (lemme
\ref{lemma-numerical-type-minimal-model}). D'après la proposition
\ref{proposition-bound-picard-group}, on a
$$
\dim_{\mathbf{F}_\ell} \Pic(T)[\ell] \leq g_{top}
$$
D'après les lemmes \ref{lemma-sequence-torsion} et
\ref{lemma-torsion-embeds}, on conclut qu'il existe un plongement
$$
(\mathbf{Z}/\ell \mathbf{Z})^{\oplus 2g - g_{top}}
\subset \Pic((X_k)_{red})[\ell].
$$
D'après Courbes algébriques, lemme
\ref{curves-lemma-bound-torsion-simple}, on obtient
$$
2g - g_{top} \leq
\dim_k H^1((X_k)_{red}, \mathcal{O}_{(X_k)_{red}}) + g_{geom}(X_k/k)
$$
D'après les lemmes \ref{lemma-genus-reduction-smaller} et
\ref{lemma-genus-reduction-bigger-than}, on a
$$
g \geq \dim_k H^1((X_k)_{red}, \mathcal{O}_{(X_k)_{red}}) \geq
g_{top} + g_{geom}(X_k/k)
$$
La théorie élémentaire des nombres montre que ces $3$ inégalités ne
peuvent être satisfaites que si elles sont toutes des égalités.
L'examen du lemme \ref{lemma-genus-reduction-smaller} montre que
$m_i=1$ pour tout $i$. Celui du lemme
\ref{lemma-equality-genus-reduction-bigger-than} montre que toute composante
irréductible de $X_k$ est lisse sur $k$.

\medskip\noindent
En particulier, puisque $X_k$ est la réunion schématique de ses composantes
irréductibles $C_i$, on voit que $X_{\overline{k}}$ est la réunion
schématique des $C_{i,\overline{k}}$. Ainsi $X_{\overline{k}}$ est un
schéma propre réduit connexe de dimension $1$ sur $\overline{k}$, tel que
$\dim_{\overline{k}}H^1(X_{\overline{k}},
\mathcal{O}_{X_{\overline{k}}})=g$. De plus, d'après Variétés, lemme
\ref{varieties-lemma-change-fields-pic}, et ce qui précède, on a encore
$$
\dim_{\mathbf{F}_\ell}(\Pic(X_{\overline{k}})[\ell])
\geq 2g - g_{top} =
\dim_{\overline{k}} H^1(X_{\overline{k}}, \mathcal{O}_{X_{\overline{k}}})
+ g_{geom}(X_{\overline{k}})
$$
D'après Courbes algébriques, proposition
\ref{curves-proposition-torsion-picard-reduced-proper}, on voit que
$X_{\overline{k}}$ n'a que des singularités en croisement multiple.
Mais $X_k$ est de Gorenstein (lemme \ref{lemma-gorenstein}), donc
$X_{\overline{k}}$ l'est aussi (Dualité pour les schémas, lemme
\ref{duality-lemma-gorenstein-base-change}). On conclut que les singularités
de $X_{\overline{k}}$ sont au plus nodales d'après Courbes algébriques,
lemme \ref{curves-lemma-multicross-gorenstein-is-nodal}. Ceci achève la
preuve.









\section{Réduction semi-stable des courbes}
\label{section-semistable-reduction-theorem}

\noindent
Dans cette section, nous achevons la preuve du théorème. Pour $g\geq2$,
soient $768g<\ell'<\ell$ les deux premiers nombres premiers $>768g$, et
posons
\begin{equation}
\label{equation-bound}
B_g = (2g - 2)(\ell^{2g})!
\end{equation}
La forme précise de $B_g$ n'a pas d'importance ; le point essentiel est
qu'il ne dépend que de $g$.

\begin{theorem}
\label{theorem-semistable-reduction}
\begin{reference}
\cite[Corollary 2.7]{DM}
\end{reference}
Soit $R$ un anneau de valuation discrète de corps des fractions $K$.
Soit $C$ une courbe projective lisse sur $K$ telle que
$H^0(C,\mathcal{O}_C)=K$. Il existe alors une extension d'anneaux de
valuation discrète $R\subset R'$ qui induit une extension finie séparable
$K'/K$ des corps des fractions et telle que $C_{K'}$ ait une réduction
semi-stable. Plus précisément :
\begin{enumerate}
\item Si le genre de $C$ est zéro, il existe une extension séparable de
degré $2$, notée $K'/K$, telle que $C_{K'}\cong\mathbf{P}^1_{K'}$ ; ainsi
$C_{K'}$ est isomorphe à la fibre générique du schéma projectif lisse
$\mathbf{P}^1_{R'}$ sur la clôture intégrale $R'$ de $R$ dans $K'$.
\item Si le genre de $C$ est un, il existe une extension finie séparable
$K'/K$ telle que $C_{K'}$ ait une réduction semi-stable sur
$R'_\mathfrak m$ pour tout idéal maximal $\mathfrak m$ de la clôture
intégrale $R'$ de $R$ dans $K'$. De plus, la fibre spéciale du modèle
minimal (unique) de $C_{K'}$ sur $R'_\mathfrak m$ est soit une courbe
lisse de genre un, soit un cycle de courbes rationnelles.
\item Si le genre $g$ de $C$ est supérieur à un, il existe une extension
finie séparable $K'/K$ de degré au plus $B_g$ (\ref{equation-bound}) telle
que $C_{K'}$ ait une réduction semi-stable sur $R'_\mathfrak m$ pour
tout idéal maximal $\mathfrak m$ de la clôture intégrale $R'$ de $R$ dans
$K'$.
\end{enumerate}
\end{theorem}

\begin{proof}
Pour le cas du genre zéro, voir la section
\ref{section-semistable-reduction-genus-zero}. Pour celui du genre un, voir
la section \ref{section-semistable-reduction-genus-one}. Pour celui du
genre supérieur à un, voir la section
\ref{section-semistable-reduction-genus-at-least-two}. Pour obtenir la borne
sur le degré $[K':K]$, on peut utiliser la borne sur le degré de l'extension
nécessaire pour rendre visible toute la $\ell$-torsion ou toute la
$\ell'$-torsion, démontrée dans Courbes algébriques, lemme
\ref{curves-lemma-torsion-picard-becomes-visible}. (La raison d'utiliser
$\ell$ et $\ell'$ est qu'il faut éviter la caractéristique du corps
résiduel $k$.)
\end{proof}

\begin{remark}[Amélioration de la borne]
\label{remark-improving-bound}
Les résultats de la littérature suggèrent que l'on peut améliorer la borne
donnée dans l'énoncé du théorème \ref{theorem-semistable-reduction}.
Par exemple, \cite{DM} montre que la réduction semi-stable de $C$ équivaut
à celle de sa jacobienne lorsque le corps résiduel est parfait, et cela est
vraisemblablement encore vrai pour un corps résiduel quelconque.
Une variété abélienne a une réduction semi-stable si l'action de Galois sur
sa $\ell$-torsion est triviale pour un $\ell\geq3$ quelconque distinct de
la caractéristique résiduelle. On peut donc vraisemblablement choisir
$\ell=5$ dans la formule (\ref{equation-bound}) définissant $B_g$ (mais la
preuve demanderait beaucoup plus de travail ; si nous en avons un jour
besoin, nous donnerons ici un énoncé précis et une preuve).
\end{remark}





\begin{multicols}{2}[\section{Autres chapitres}]
\noindent
Préliminaires
\begin{enumerate}
\item \hyperref[introduction-section-phantom]{Introduction}
\item \hyperref[conventions-section-phantom]{Conventions}
\item \hyperref[sets-section-phantom]{Théorie des ensembles}
\item \hyperref[categories-section-phantom]{Catégories}
\item \hyperref[topology-section-phantom]{Topologie}
\item \hyperref[sheaves-section-phantom]{Faisceaux sur les espaces}
\item \hyperref[sites-section-phantom]{Sites et faisceaux}
\item \hyperref[stacks-section-phantom]{Champs}
\item \hyperref[fields-section-phantom]{Corps}
\item \hyperref[algebra-section-phantom]{Algèbre commutative}
\item \hyperref[brauer-section-phantom]{Groupes de Brauer}
\item \hyperref[homology-section-phantom]{Algèbre homologique}
\item \hyperref[derived-section-phantom]{Catégories dérivées}
\item \hyperref[simplicial-section-phantom]{Méthodes simpliciales}
\item \hyperref[more-algebra-section-phantom]{Compléments d'algèbre}
\item \hyperref[smoothing-section-phantom]{Lissage des morphismes d'anneaux}
\item \hyperref[modules-section-phantom]{Faisceaux de modules}
\item \hyperref[sites-modules-section-phantom]{Modules sur les sites}
\item \hyperref[injectives-section-phantom]{Objets injectifs}
\item \hyperref[cohomology-section-phantom]{Cohomologie des faisceaux}
\item \hyperref[sites-cohomology-section-phantom]{Cohomologie sur les sites}
\item \hyperref[dga-section-phantom]{Algèbre différentielle graduée}
\item \hyperref[dpa-section-phantom]{Algèbre à puissances divisées}
\item \hyperref[sdga-section-phantom]{Faisceaux différentiels gradués}
\item \hyperref[hypercovering-section-phantom]{Hyperrecouvrements}
\end{enumerate}
Schémas
\begin{enumerate}
\setcounter{enumi}{25}
\item \hyperref[schemes-section-phantom]{Schémas}
\item \hyperref[constructions-section-phantom]{Constructions de schémas}
\item \hyperref[properties-section-phantom]{Propriétés des schémas}
\item \hyperref[morphisms-section-phantom]{Morphismes de schémas}
\item \hyperref[coherent-section-phantom]{Cohomologie des schémas}
\item \hyperref[divisors-section-phantom]{Diviseurs}
\item \hyperref[limits-section-phantom]{Limites de schémas}
\item \hyperref[varieties-section-phantom]{Variétés}
\item \hyperref[topologies-section-phantom]{Topologies sur les schémas}
\item \hyperref[descent-section-phantom]{Descente}
\item \hyperref[perfect-section-phantom]{Catégories dérivées de schémas}
\item \hyperref[more-morphisms-section-phantom]{Compléments sur les morphismes}
\item \hyperref[flat-section-phantom]{Compléments sur la platitude}
\item \hyperref[groupoids-section-phantom]{Schémas en groupoïdes}
\item \hyperref[more-groupoids-section-phantom]{Compléments sur les schémas en groupoïdes}
\item \hyperref[etale-section-phantom]{Morphismes étales de schémas}
\end{enumerate}
Thèmes de théorie des schémas
\begin{enumerate}
\setcounter{enumi}{41}
\item \hyperref[chow-section-phantom]{Homologie de Chow}
\item \hyperref[intersection-section-phantom]{Théorie de l'intersection}
\item \hyperref[pic-section-phantom]{Schémas de Picard des courbes}
\item \hyperref[weil-section-phantom]{Théories cohomologiques de Weil}
\item \hyperref[adequate-section-phantom]{Modules adéquats}
\item \hyperref[dualizing-section-phantom]{Complexes dualisants}
\item \hyperref[duality-section-phantom]{Dualité pour les schémas}
\item \hyperref[discriminant-section-phantom]{Discriminants et différentes}
\item \hyperref[derham-section-phantom]{Cohomologie de de Rham}
\item \hyperref[local-cohomology-section-phantom]{Cohomologie locale}
\item \hyperref[algebraization-section-phantom]{Géométrie algébrique et formelle}
\item \hyperref[curves-section-phantom]{Courbes algébriques}
\item \hyperref[resolve-section-phantom]{Résolution des surfaces}
\item \hyperref[models-section-phantom]{Réduction semi-stable}
\item \hyperref[functors-section-phantom]{Foncteurs et morphismes}
\item \hyperref[equiv-section-phantom]{Catégories dérivées de variétés}
\item \hyperref[pione-section-phantom]{Groupes fondamentaux de schémas}
\item \hyperref[etale-cohomology-section-phantom]{Cohomologie étale}
\item \hyperref[crystalline-section-phantom]{Cohomologie cristalline}
\item \hyperref[proetale-section-phantom]{Cohomologie pro-étale}
\item \hyperref[relative-cycles-section-phantom]{Cycles relatifs}
\item \hyperref[more-etale-section-phantom]{Compléments de cohomologie étale}
\item \hyperref[trace-section-phantom]{La formule des traces}
\end{enumerate}
Espaces algébriques
\begin{enumerate}
\setcounter{enumi}{64}
\item \hyperref[spaces-section-phantom]{Espaces algébriques}
\item \hyperref[spaces-properties-section-phantom]{Propriétés des espaces algébriques}
\item \hyperref[spaces-morphisms-section-phantom]{Morphismes d'espaces algébriques}
\item \hyperref[decent-spaces-section-phantom]{Espaces algébriques décents}
\item \hyperref[spaces-cohomology-section-phantom]{Cohomologie des espaces algébriques}
\item \hyperref[spaces-limits-section-phantom]{Limites d'espaces algébriques}
\item \hyperref[spaces-divisors-section-phantom]{Diviseurs sur les espaces algébriques}
\item \hyperref[spaces-over-fields-section-phantom]{Espaces algébriques sur des corps}
\item \hyperref[spaces-topologies-section-phantom]{Topologies sur les espaces algébriques}
\item \hyperref[spaces-descent-section-phantom]{Descente et espaces algébriques}
\item \hyperref[spaces-perfect-section-phantom]{Catégories dérivées d'espaces}
\item \hyperref[spaces-more-morphisms-section-phantom]{Compléments sur les morphismes d'espaces}
\item \hyperref[spaces-flat-section-phantom]{Platitude sur les espaces algébriques}
\item \hyperref[spaces-groupoids-section-phantom]{Groupoïdes dans les espaces algébriques}
\item \hyperref[spaces-more-groupoids-section-phantom]{Compléments sur les groupoïdes dans les espaces}
\item \hyperref[bootstrap-section-phantom]{Amorçage}
\item \hyperref[spaces-pushouts-section-phantom]{Sommes amalgamées d'espaces algébriques}
\end{enumerate}
Thèmes de géométrie
\begin{enumerate}
\setcounter{enumi}{81}
\item \hyperref[spaces-chow-section-phantom]{Groupes de Chow des espaces}
\item \hyperref[groupoids-quotients-section-phantom]{Quotients de groupoïdes}
\item \hyperref[spaces-more-cohomology-section-phantom]{Compléments sur la cohomologie des espaces}
\item \hyperref[spaces-simplicial-section-phantom]{Espaces simpliciaux}
\item \hyperref[spaces-duality-section-phantom]{Dualité pour les espaces}
\item \hyperref[formal-spaces-section-phantom]{Espaces algébriques formels}
\item \hyperref[restricted-section-phantom]{Algébrisation des espaces formels}
\item \hyperref[spaces-resolve-section-phantom]{Résolution des surfaces revisitée}
\end{enumerate}
Théorie des déformations
\begin{enumerate}
\setcounter{enumi}{89}
\item \hyperref[formal-defos-section-phantom]{Théorie formelle des déformations}
\item \hyperref[defos-section-phantom]{Théorie des déformations}
\item \hyperref[cotangent-section-phantom]{Le complexe cotangent}
\item \hyperref[examples-defos-section-phantom]{Problèmes de déformation}
\end{enumerate}
Champs algébriques
\begin{enumerate}
\setcounter{enumi}{93}
\item \hyperref[algebraic-section-phantom]{Champs algébriques}
\item \hyperref[examples-stacks-section-phantom]{Exemples de champs}
\item \hyperref[stacks-sheaves-section-phantom]{Faisceaux sur les champs algébriques}
\item \hyperref[criteria-section-phantom]{Critères de représentabilité}
\item \hyperref[artin-section-phantom]{Axiomes d'Artin}
\item \hyperref[quot-section-phantom]{Espaces de Quot et de Hilbert}
\item \hyperref[stacks-properties-section-phantom]{Propriétés des champs algébriques}
\item \hyperref[stacks-morphisms-section-phantom]{Morphismes de champs algébriques}
\item \hyperref[stacks-limits-section-phantom]{Limites de champs algébriques}
\item \hyperref[stacks-cohomology-section-phantom]{Cohomologie des champs algébriques}
\item \hyperref[stacks-perfect-section-phantom]{Catégories dérivées de champs}
\item \hyperref[stacks-introduction-section-phantom]{Introduction aux champs algébriques}
\item \hyperref[stacks-more-morphisms-section-phantom]{Compléments sur les morphismes de champs}
\item \hyperref[stacks-geometry-section-phantom]{La géométrie des champs}
\end{enumerate}
Thèmes de théorie des espaces de modules
\begin{enumerate}
\setcounter{enumi}{107}
\item \hyperref[moduli-section-phantom]{Champs de modules}
\item \hyperref[moduli-curves-section-phantom]{Espaces de modules de courbes}
\end{enumerate}
Divers
\begin{enumerate}
\setcounter{enumi}{109}
\item \hyperref[examples-section-phantom]{Exemples}
\item \hyperref[exercises-section-phantom]{Exercices}
\item \hyperref[guide-section-phantom]{Guide bibliographique}
\item \hyperref[desirables-section-phantom]{Desiderata}
\item \hyperref[coding-section-phantom]{Style de programmation}
\item \hyperref[obsolete-section-phantom]{Obsolète}
\item \hyperref[fdl-section-phantom]{Licence de documentation libre GNU}
\item \hyperref[index-section-phantom]{Index généré automatiquement}
\end{enumerate}
\end{multicols}


\bibliography{my}
\bibliographystyle{amsalpha}

\end{document}
